% L’existence et la mort — Philosophie Tle Tech — Cours signé OnBuch+
% Programme officiel MINESEC. Compiler avec : tectonic N11-existence-mort.tex
\newcommand{\DOCMATIERE}{Philosophie}
\newcommand{\DOCNIVEAU}{Tle Tech}
\newcommand{\DOCMODULE}{Philosophie – classe de Terminale technique}
\newcommand{\DOCLECON}{N11}
\newcommand{\DOCTITRE}{L’existence et la mort}
% OnBuch+ — préambule « cours signés » ENRICHI (compilé avec tectonic / XeLaTeX)
% Système visuel « Pop » d'OnBuch+ : fond crème (jamais blanc), bordures encre
% épaisses, ombres « dures » décalées, titres Archivo Black en capitales.
% Version MATHÉMATIQUES : graphes (pgfplots/tikz), boîtes pédagogiques
% (définition, propriété, méthode, attention, le savais-tu, exercice résolu,
% exercice, TP, bilan), page de garde et signature OnBuch+ ancrée en bas.
%
% Macros attendues dans main.tex AVANT \input{preamble} :
%   \DOCMATIERE  \DOCNIVEAU  \DOCMODULE  \DOCLECON (numéro)  \DOCTITRE
\documentclass[11pt]{article}

\usepackage{fontspec}
\usepackage[french]{babel}
\usepackage[a4paper,margin=2cm,top=3.4cm,bottom=2.8cm,headheight=1.4cm,footskip=1.3cm]{geometry}
\usepackage{amsmath,amssymb,amsfonts}
\usepackage{mathtools} % \xmapsto, \coloneqq... souvent utilisés par les modèles
\usepackage{cancel} % simplifications barrées (bilans de forces)
% \abs{z} (module, valeur absolue) et \norm{u} (norme), utilisés par les modèles
\providecommand{\abs}{}\let\abs\relax\DeclarePairedDelimiter{\abs}{\lvert}{\rvert}
\providecommand{\norm}{}\let\norm\relax\DeclarePairedDelimiter{\norm}{\lVert}{\rVert}
\usepackage[table]{xcolor} % [table] : fournit aussi \rowcolors (alternance de lignes)
\usepackage{pagecolor}
\usepackage{tikz}
\usetikzlibrary{arrows.meta,positioning,calc,shapes.geometric,shapes.misc,shapes.arrows,decorations.pathmorphing,decorations.pathreplacing,decorations.markings,patterns,shadows,mindmap,trees,fit,backgrounds,matrix,intersections,angles,quotes}
\usepackage{pgfplots}
\usepackage[version=4]{mhchem} % équations nucléaires \ce{^{235}_{92}U}
\usepackage[european,siunitx]{circuitikz} % schémas électriques
\pgfplotsset{compat=1.18}
\usepgfplotslibrary{fillbetween,groupplots} % aires entre courbes (calcul intégral)
\usepackage{enumitem}
\usepackage{fancyhdr}
% [most] charge toutes les bibliothèques tcolorbox (listings, minted, raster,
% documentation...) et gonfle inutilement la mémoire TeX sur un document long
% (~300 boîtes) : on ne charge que ce qui est réellement utilisé.
\usepackage[skins,breakable]{tcolorbox}
\usepackage{graphicx}
\usepackage{colortbl}
\usepackage{booktabs}
\usepackage{tabularx}
\usepackage{multirow}
\usepackage{diagbox} % case d'en-tête diagonale des tableaux à double entrée
% Colonnes extensibles centrée / gauche / droite (C, L, R), souvent utilisées par les modèles
\newcolumntype{C}{>{\centering\arraybackslash}X}
\newcolumntype{L}{>{\raggedright\arraybackslash}X}
\newcolumntype{R}{>{\raggedleft\arraybackslash}X}
\usepackage{array}
\usepackage{multicol}
\usepackage{siunitx}
% parse-numbers=false : accepte aussi \SI{2,5 \times 10^{-3}}{...}
\sisetup{locale=FR,output-decimal-marker={,},group-minimum-digits=4,parse-numbers=false}
\DeclareSIUnit\ohm{\ensuremath{\symup{\Omega}}} % U+2126 absent de la police
\DeclareSIUnit\division{div} % réglages d'oscilloscope (ms/div, V/div)
\DeclareSIUnit\tr{tr} % tours (vitesse de rotation, ex. \SI{1500}{\tr\per\minute})
\DeclareSIUnit\molar{M} % concentration molaire (ex. \SI{3}{\molar}, \SI{1}{\micro\molar})
\DeclareSIUnit\an{an} % année (le modèle confond parfois avec \year, un compteur TeX) — ex. \SI{5}{\kilo\meter\per\an}
\DeclareSIUnit\FCFA{FCFA} % franc CFA (ex. \SI{500}{\FCFA\per\kilo\gram})
\DeclareSIUnit\degreeBrix{\textdegree Brix} % teneur en sucre d'un sirop/jus (réfractométrie)
\DeclareSIUnit\month{mois} % le modèle utilise parfois \month, un compteur TeX, comme unité de durée
\DeclareSIUnit\microliter{\micro\liter} % le modèle écrit parfois \microliter en un seul token
\DeclareSIUnit\million{millions} % dénombrement (ex. \SI{15}{\million\per\mL} spermatozoïdes)
\usepackage{qrcode}
% \degree : commande fréquemment utilisée par les modèles pour un angle ou une
% température en dehors de \SI{}{} (non fournie par les paquets chargés).
\providecommand{\degree}{^{\circ}}
% \qed (fin de démonstration), utilisé par les modèles sans amsthm
\providecommand{\qed}{\hfill$\square$}
% \barre : double barre des valeurs interdites dans un tableau de variations (tkz-tab)
\providecommand{\barre}{\ensuremath{\|}}
% environnement proof (amsthm non chargé) utilisé par les modèles
\ifdefined\proof\else\newenvironment{proof}[1][Démonstration]{\par\noindent\textit{#1.}\ }{\qed\par}\fi
\usepackage{lastpage}
\usepackage{hyperref}
\hypersetup{hidelinks}

% ============================================================
% Palette « Pop » OnBuch+ — mêmes hex que la classe P (plus_ui.dart)
% ============================================================
\definecolor{poporange}{HTML}{F59321}
\definecolor{popdark}{HTML}{E07A0C}
\definecolor{popcream}{HTML}{FFF6E8}
\definecolor{popink}{HTML}{1C1714}
\definecolor{poppurple}{HTML}{7A5AE0}
\definecolor{popgreen}{HTML}{2E9E6A}
\definecolor{poppink}{HTML}{E0457B}
\definecolor{popblue}{HTML}{2F7DE1}
\definecolor{popgold}{HTML}{C98A12}
\definecolor{popmuted}{HTML}{8A7F76}
\definecolor{popline}{HTML}{EADFD2}
\definecolor{popgray}{HTML}{8A7F76}
\colorlet{popgrey}{popgray}
% Alias tolérants (le modèle nomme parfois les couleurs différemment)
\colorlet{popwhite}{white}
\colorlet{popblack}{popink}
\colorlet{popred}{poppink}
\colorlet{popyellow}{popgold}
% Teintes claires (fonds de tableaux/encadrés) — le modèle nomme parfois
% les couleurs avec un suffixe L (ex. popblueL) sans les avoir définies.
\colorlet{poporangeL}{poporange!20}
\colorlet{popdarkL}{popdark!20}
\colorlet{poppurpleL}{poppurple!20}
\colorlet{popgreenL}{popgreen!20}
\colorlet{poppinkL}{poppink!20}
\colorlet{popblueL}{popblue!20}
\colorlet{popgoldL}{popgold!20}
\colorlet{popredL}{popred!20}
\colorlet{popgrayL}{popgray!20}
% Teintes claires pour fonds de tableaux / zones
\colorlet{popblueL}{popblue!10!white}
\colorlet{poporangeL}{poporange!14!white}
\colorlet{popgreenL}{popgreen!12!white}
\colorlet{poppurpleL}{poppurple!12!white}
\colorlet{poppinkL}{poppink!12!white}
\colorlet{popgoldL}{popgold!14!white}

\pagecolor{popcream}

% ============================================================
% Typographie
% ============================================================
\defaultfontfeatures{Ligatures=TeX}
\setmainfont{PlusJakartaSans-Medium.ttf}[Path=../../../fonts/,BoldFont=PlusJakartaSans-Bold.ttf]
% Maths en sans-serif (Fira Math) avec chiffres et lettres droites de Plus
% Jakarta, pour que les formules se fondent dans le texte.
\usepackage[math-style=ISO,bold-style=ISO]{unicode-math}
\setmathfont{FiraMath-Regular.otf}[Path=../../../fonts/]
\setmathfont{PlusJakartaSans-Medium.ttf}[Path=../../../fonts/,range={up/num,up/{Latin,latin},\mathup}]
\usepackage{icomma} % virgule décimale sans espace en mode math
% Caractères mathématiques Unicode absents des polices de texte (Plus Jakarta,
% Archivo Black) : rendus par leur équivalent mathématique, en texte comme en
% formule — sinon ils s'affichent en carré vide (ex. « Arithmétique dans ℤ »).
\usepackage{newunicodechar}
\newunicodechar{ℝ}{\ensuremath{\mathbb{R}}}
\newunicodechar{ℤ}{\ensuremath{\mathbb{Z}}}
\newunicodechar{ℂ}{\ensuremath{\mathbb{C}}}
\newunicodechar{ℕ}{\ensuremath{\mathbb{N}}}
\newunicodechar{ℚ}{\ensuremath{\mathbb{Q}}}
\newunicodechar{𝔻}{\ensuremath{\mathbb{D}}}
\newunicodechar{α}{\ensuremath{\alpha}}
\newunicodechar{β}{\ensuremath{\beta}}
\newunicodechar{γ}{\ensuremath{\gamma}}
\newunicodechar{δ}{\ensuremath{\delta}}
\newunicodechar{ε}{\ensuremath{\varepsilon}}
\newunicodechar{θ}{\ensuremath{\theta}}
\newunicodechar{λ}{\ensuremath{\lambda}}
\newunicodechar{μ}{\ensuremath{\mu}}
\newunicodechar{σ}{\ensuremath{\sigma}}
\newunicodechar{τ}{\ensuremath{\tau}}
\newunicodechar{φ}{\ensuremath{\varphi}}
\newunicodechar{ψ}{\ensuremath{\psi}}
\newunicodechar{ω}{\ensuremath{\omega}}
\newunicodechar{Σ}{\ensuremath{\Sigma}}
% majuscules produites par \MakeUppercase dans les titres (α-aminés -> Α)
\newunicodechar{Α}{\ensuremath{\alpha}}
\newunicodechar{Β}{\ensuremath{\beta}}
\newunicodechar{Γ}{\ensuremath{\Gamma}}
\newunicodechar{Δ}{\ensuremath{\Delta}}
\newunicodechar{Ε}{\ensuremath{\varepsilon}}
\newunicodechar{Θ}{\ensuremath{\Theta}}
\newunicodechar{Λ}{\ensuremath{\Lambda}}
\newunicodechar{Μ}{\ensuremath{\mu}}
\newunicodechar{Τ}{\ensuremath{\tau}}
\newunicodechar{Φ}{\ensuremath{\Phi}}
\newunicodechar{Ψ}{\ensuremath{\Psi}}
\newunicodechar{Ω}{\ensuremath{\Omega}}
\newunicodechar{↦}{\ensuremath{\mapsto}}
\newunicodechar{∈}{\ensuremath{\in}}
\newunicodechar{∉}{\ensuremath{\notin}}
\newunicodechar{∧}{\ensuremath{\wedge}}
\newunicodechar{⇔}{\ensuremath{\Leftrightarrow}}
\newunicodechar{⟺}{\ensuremath{\Longleftrightarrow}}
\newunicodechar{⇒}{\ensuremath{\Rightarrow}}
\newunicodechar{⟹}{\ensuremath{\Longrightarrow}}
\newunicodechar{∘}{\ensuremath{\circ}}
\newunicodechar{∪}{\ensuremath{\cup}}
\newunicodechar{∩}{\ensuremath{\cap}}
\newunicodechar{≡}{\ensuremath{\equiv}}
\newunicodechar{⊂}{\ensuremath{\subset}}
\newunicodechar{⊥}{\ensuremath{\perp}}
\newunicodechar{∃}{\ensuremath{\exists}}
\newunicodechar{∀}{\ensuremath{\forall}}
\newunicodechar{∛}{\ensuremath{\sqrt[3]{\,}}}
\newunicodechar{∜}{\ensuremath{\sqrt[4]{\,}}}
\newunicodechar{⃗}{\ensuremath{{}^{\to}}}
\newunicodechar{ⁿ}{\ifmmode^{n}\else\textsuperscript{\ensuremath{n}}\fi}
\newunicodechar{ˣ}{\ifmmode^{x}\else\textsuperscript{\ensuremath{x}}\fi}
\newunicodechar{ᵇ}{\ifmmode^{b}\else\textsuperscript{\ensuremath{b}}\fi}
\newunicodechar{ʳ}{\ifmmode^{r}\else\textsuperscript{\ensuremath{r}}\fi}
\newunicodechar{ᵘ}{\ifmmode^{u}\else\textsuperscript{\ensuremath{u}}\fi}
\newunicodechar{ʸ}{\ifmmode^{y}\else\textsuperscript{\ensuremath{y}}\fi}
\newunicodechar{ᵃ}{\ifmmode^{a}\else\textsuperscript{\ensuremath{a}}\fi}
\newunicodechar{ᵉ}{\ifmmode^{e}\else\textsuperscript{\ensuremath{e}}\fi}
\newunicodechar{ᵏ}{\ifmmode^{k}\else\textsuperscript{\ensuremath{k}}\fi}
\newunicodechar{ᵖ}{\ifmmode^{p}\else\textsuperscript{\ensuremath{p}}\fi}
\newunicodechar{ᴺ}{\ifmmode^{N}\else\textsuperscript{\ensuremath{N}}\fi}
\newunicodechar{ᶜ}{\ifmmode^{c}\else\textsuperscript{\ensuremath{c}}\fi}
\newunicodechar{ʰ}{\ifmmode^{h}\else\textsuperscript{\ensuremath{h}}\fi}
\newunicodechar{ᵠ}{\ifmmode^{q}\else\textsuperscript{\ensuremath{q}}\fi}
\newunicodechar{ₙ}{\ifmmode_{n}\else\textsubscript{\ensuremath{n}}\fi}
\newunicodechar{ₐ}{\ifmmode_{a}\else\textsubscript{\ensuremath{a}}\fi}
\newunicodechar{ᵢ}{\ifmmode_{i}\else\textsubscript{\ensuremath{i}}\fi}
\newunicodechar{ᵣ}{\ifmmode_{r}\else\textsubscript{\ensuremath{r}}\fi}
\newunicodechar{ₖ}{\ifmmode_{k}\else\textsubscript{\ensuremath{k}}\fi}
\newunicodechar{ᵦ}{\ifmmode_{\beta}\else\textsubscript{\ensuremath{\beta}}\fi}
\newunicodechar{ₓ}{\ifmmode_{x}\else\textsubscript{\ensuremath{x}}\fi}
\newfontfamily\popdisplay{ArchivoBlack-Regular.ttf}[Path=../../../fonts/]
\newfontfamily\poplogo{SpaceGrotesk-Bold.ttf}[Path=../../../fonts/]
\newfontfamily\popsemi{PlusJakartaSans-SemiBold.ttf}[Path=../../../fonts/]
\newfontfamily\popxbold{PlusJakartaSans-ExtraBold.ttf}[Path=../../../fonts/]
\newfontfamily\popxboldls{PlusJakartaSans-ExtraBold.ttf}[Path=../../../fonts/,LetterSpace=3.0]

\setlist[enumerate]{leftmargin=1.7em,itemsep=5pt,topsep=4pt}
\setlist[itemize]{leftmargin=1.7em,itemsep=4pt,topsep=4pt,label={\color{poporange}\textbullet}}
\setlength{\parindent}{0pt}
\setlength{\parskip}{5pt}
\renewcommand{\arraystretch}{1.25}

% pgfplots : style Pop par défaut
\pgfplotsset{
  every axis/.append style={axis line style={popink,line width=1pt},tick style={popink},
    grid style={popline,line width=0.5pt},label style={font=\small},tick label style={font=\footnotesize}},
  popcurve/.style={poporange,line width=1.8pt,no markers,smooth},
}
% Même style disponible dans un \draw TikZ simple (hors pgfplots)
\tikzset{popcurve/.style={poporange,line width=1.8pt}}
\tikzset{popgrid/.style={popline,very thin}}
\colorlet{popcurve}{poporange} % le nom du style est parfois utilisé comme couleur

% ============================================================
% Pill / squiggle
% ============================================================
\newcommand{\poppill}[3]{%
  \tikz[baseline=-0.55ex]\node[draw=popink,line width=1.2pt,rounded corners=2.6pt,%
    fill=#2,inner xsep=6.5pt,inner ysep=3pt]{%
    \popxboldls\fontsize{7}{7}\selectfont\color{#3}\MakeUppercase{#1}};%
}
\newcommand{\popsquiggle}[1]{%
  \tikz[baseline=-0.4ex]{
    \draw[#1,line width=2.6pt,line cap=round]
      plot [smooth] coordinates {(0,0.09) (0.34,0.01) (0.68,0.21) (1.02,0.01) (1.36,0.10)};
  }%
}
% Mot-clé surligné (marqueur orange sous le texte)
\newcommand{\cle}[1]{{\popxbold\color{popdark}#1}}

% ============================================================
% En-tête / pied
% ============================================================
\pagestyle{fancy}
\fancyhf{}
\renewcommand{\headrulewidth}{0pt}
\renewcommand{\footrulewidth}{0pt}
\lhead{\raisebox{-0.15ex}{\poplogo\fontsize{13}{13}\selectfont\color{popink}OnBuch\color{poporange}+}}
\rhead{\poppill{\DOCMATIERE\ \textbullet\ \DOCNIVEAU\ \textbullet\ Leçon \DOCLECON}{white}{popink}}
\lfoot{\popxbold\fontsize{7.6}{7.6}\selectfont\color{popmuted}\MakeUppercase{Cours signé OnBuch+}\ \textbullet\ {\popsemi\DOCTITRE}}
\rfoot{\tikz[baseline=-0.6ex]\node[rounded corners=8.5pt,fill=popink,minimum height=17pt,minimum width=17pt,inner xsep=5pt,inner sep=0]{\popxbold\fontsize{8}{8}\selectfont\color{popcream}\thepage\,/\,\pageref*{LastPage}};}

% ============================================================
% Titres
% ============================================================
\newcommand{\chaptitre}[2]{%
  \begin{tcolorbox}[enhanced,colback=poporange,colframe=popink,boxrule=2.6pt,arc=11pt,
      drop shadow=popink,left=15pt,right=15pt,top=12pt,bottom=11pt,before skip=3pt,after skip=18pt]
    \poppill{#2}{popink}{popcream}\par\vspace{9pt}
    {\raggedright\hyphenpenalty=10000\exhyphenpenalty=10000\popdisplay\fontsize{22}{24}\selectfont\color{popcream}\MakeUppercase{#1}\par}
  \end{tcolorbox}
}
% Section numérotée : pastille orange avec le numéro + titre Archivo
\newcounter{popsec}
\newcommand{\coursec}[1]{%
  \stepcounter{popsec}\setcounter{popsub}{0}%
  \par\vspace{16pt}\noindent
  \begin{minipage}{\linewidth}
  \tikz[baseline=-0.7ex]\node[draw=popink,line width=1.6pt,fill=poporange,rounded corners=4pt,minimum size=22pt,inner sep=0]{\popdisplay\fontsize{12}{12}\selectfont\color{popcream}\thepopsec};%
  \hspace{8pt}{\popdisplay\fontsize{15}{17}\selectfont\color{popink}\MakeUppercase{#1}}\par
  \vspace{4pt}\noindent{\color{poporange}\rule{36pt}{3.6pt}}
  \end{minipage}\par\nopagebreak\vspace{8pt}
}
\newcounter{popsub}
\newcommand{\courssub}[1]{%
  \stepcounter{popsub}%
  \par\vspace{9pt}\noindent{\popxbold\fontsize{11.5}{13}\selectfont\color{popink}{\color{poporange}\thepopsec.\thepopsub}\hspace{6pt}#1}\par\nopagebreak\vspace{4pt}
}

% ============================================================
% Boîtes pédagogiques — même recette : fond blanc, bordure épaisse colorée,
% ombre dure encre, badge plein. Titre optionnel [#1].
% ============================================================
\tcbset{popbase/.style={breakable,enhanced,colback=white,boxrule=2.3pt,arc=9pt,
  shadow={3pt}{-3pt}{0pt}{popink},left=14pt,right=14pt,top=11pt,bottom=11pt,before skip=11pt,after skip=15pt,
  fontupper=\popsemi\fontsize{10.6}{14.5}\selectfont\color{popink},
  fontlower=\popsemi\fontsize{10.6}{14.5}\selectfont\color{popink}}}
% Coche dessinée (le glyphe ✓ n'existe pas dans les polices du document)
\newcommand{\popcheck}[1]{\tikz[baseline=-0.5ex]\draw[#1,line width=1.6pt,line cap=round,line join=round](-0.13,0.0)--(-0.04,-0.09)--(0.13,0.1);}
\providecommand{\coche}{\popcheck{popgreen}}
\providecommand{\resultat}[1]{\par\smallskip\noindent\fcolorbox{popgreen}{popgreenL}{\parbox{\dimexpr\linewidth-2\fboxsep-2\fboxrule\relax}{\textbf{Résultat.} #1}}\par\smallskip}
\newcommand{\popboxhead}[3]{\poppill{#1}{#2}{white}\ifx\relax#3\relax\else\hspace{6pt}{\popxbold\fontsize{10.6}{12}\selectfont\color{popink}#3}\fi\par\vspace{7pt}}

\newtcolorbox{aretenir}[1][]{popbase,colframe=popgreen,before upper={\popboxhead{À retenir}{popgreen}{#1}}}
\newtcolorbox{exemplebox}[1][]{popbase,colframe=poporange,before upper={\popboxhead{Exemple}{poporange}{#1}}}
\newtcolorbox{definition}[1][]{popbase,colframe=popblue,before upper={\popboxhead{Définition}{popblue}{#1}}}
\newtcolorbox{propriete}[1][]{popbase,colframe=poppurple,before upper={\popboxhead{Propriété}{poppurple}{#1}}}
\newtcolorbox{methode}[1][]{popbase,colframe=popgold,before upper={\popboxhead{Méthode}{popgold}{#1}}}
\newtcolorbox{attention}[1][]{popbase,colframe=poppink,before upper={\popboxhead{Attention}{poppink}{#1}}}
\newtcolorbox{experience}[1][]{popbase,colframe=popblue,colback=popblueL,before upper={\popboxhead{Expérience}{popblue}{#1}}}
% « Le savais-tu ? » : carte violette pleine (contexte, culture, Cameroun)
\newtcolorbox{savaistu}[1][]{popbase,colback=poppurple,colframe=popink,
  fontupper=\popsemi\fontsize{10.4}{14.2}\selectfont\color{white},
  before upper={\poppill{Le savais-tu\,?}{white}{poppurple}\ifx\relax#1\relax\else\hspace{6pt}{\popxbold\color{white}#1}\fi\par\vspace{7pt}}}
% Exercice résolu : énoncé en haut, \tcblower puis la solution sur fond crème
\newcounter{popexo}\newcounter{popexr}
\newtcolorbox{exoresolu}[1][]{popbase,colframe=poporange,colbacklower=popcream,
  segmentation style={popink,line width=1.2pt,dashed},
  before upper={\stepcounter{popexr}\popboxhead{Exercice résolu \thepopexr}{poporange}{#1}},
  before lower={\poppill{Solution}{popink}{popcream}\par\vspace{6pt}}}
% Exercice (non résolu en place) — la correction est dans la partie Corrigés
\newtcolorbox{exercice}[1][]{popbase,colframe=popink,
  before upper={\stepcounter{popexo}\popboxhead{Exercice \thepopexo}{popink}{#1}}}
% Corrigé
\newtcolorbox{corrige}[1][]{popbase,colframe=popgreen,colback=popgreenL,
  before upper={\popboxhead{Corrigé}{popgreen}{#1}}}
% Objectifs / prérequis / bilan
\newtcolorbox{objectifs}{popbase,colframe=popink,colback=poporangeL,
  before upper={\popboxhead{Objectifs de la leçon}{popink}{}}}
\newtcolorbox{prerequis}{popbase,colframe=popink,colback=popblueL,
  before upper={\popboxhead{Prérequis}{popblue}{}}}
\newtcolorbox{bilan}{popbase,colframe=popink,colback=white,boxrule=2.8pt,
  before upper={\popboxhead{Fiche bilan}{poporange}{L'essentiel à connaître pour l'examen}}}
% Figure encadrée avec légende
\newtcolorbox{popfigure}[1][]{enhanced,breakable=false,colback=white,colframe=popline,boxrule=1.4pt,arc=8pt,
  left=10pt,right=10pt,top=10pt,bottom=8pt,before skip=10pt,after skip=12pt,halign=center,
  lower separated=false,halign lower=center,
  fontlower=\popsemi\fontsize{9}{11.5}\selectfont\color{popmuted},
  #1}
\newcommand{\legende}[1]{\par\vspace{6pt}{\popsemi\fontsize{9}{11.5}\selectfont\color{popmuted}#1}}

% ============================================================
% Signature OnBuch+
% ============================================================
\newcommand{\poplogoinline}[1]{{\poplogo\fontsize{#1}{#1}\selectfont\color{popink}OnBuch\color{poporange}+}}
\newcommand{\signatureonbuch}{%
  \begin{tcolorbox}[enhanced,colback=popink,colframe=popink,boxrule=2.2pt,arc=10pt,
      drop shadow=poporange,left=14pt,right=14pt,top=10pt,bottom=10pt,before skip=0pt,after skip=0pt]
    \tikz[baseline=-0.7ex]\node[circle,fill=popgreen,draw=popcream,line width=1.3pt,minimum size=22pt,inner sep=0]{\popcheck{white}};%
    \hspace{9pt}\begin{minipage}[c]{0.62\linewidth}
      {\popxbold\fontsize{12}{13}\selectfont\color{popcream}Signé\ \textbullet\ {\poplogo OnBuch\color{poporange}+}}\par\vspace{2pt}
      {\raggedright\popsemi\fontsize{8}{10.5}\selectfont\color{popline}\MakeUppercase{Équipe pédagogique OnBuch+}\\\MakeUppercase{Conforme au programme officiel MINESEC}\par}
    \end{minipage}\hfill
    {\poplogo\fontsize{22}{22}\selectfont\color{popcream}OnBuch\color{poporange}+}
  \end{tcolorbox}%
}

% ============================================================
% Page de garde
% #1 titre · #2 matière · #3 niveau · #4 module · #5 n° leçon · #6 durée ·
% #7 description / objectifs (texte ou itemize)
% ============================================================
\newcommand{\pagegarde}[7]{%
  \thispagestyle{empty}
  \newgeometry{a4paper,margin=1.7cm,top=1.6cm,bottom=1.4cm}%
  \begin{tikzpicture}[remember picture,overlay]
    \fill[poporange!18] ([xshift=-3.2cm,yshift=-3.6cm]current page.north east) circle (4.6cm);
    \fill[poppurple!14] ([xshift=2.4cm,yshift=7.6cm]current page.south west) circle (3.8cm);
  \end{tikzpicture}%
  {\poplogo\fontsize{30}{30}\selectfont\color{popink}OnBuch\color{poporange}+}\hfill
  \raisebox{0.6ex}{\poppill{Cours signé \textbullet\ Premium}{popink}{popcream}}\par
  \vspace{1pt}\popsquiggle{poppurple}\par
  \vspace{8pt}
  \begin{tcolorbox}[enhanced,colback=poporange,colframe=popink,boxrule=2.8pt,arc=13pt,
      drop shadow=popink,left=18pt,right=18pt,top=12pt,bottom=13pt,after skip=10pt]
    \poppill{#2 \textbullet\ #3}{popink}{popcream}\hspace{5pt}\poppill{#4}{white}{popink}\par\vspace{8pt}
    {\popdisplay\fontsize{14}{14}\selectfont\color{popink}LEÇON #5}\par\vspace{4pt}
    {\raggedright\hyphenpenalty=10000\exhyphenpenalty=10000\popdisplay\fontsize{25}{27}\selectfont\color{popcream}\MakeUppercase{#1}\par}
    \vspace{8pt}
    {\popxbold\fontsize{9.5}{9.5}\selectfont\color{popink}\MakeUppercase{Durée conseillée : #6}}
  \end{tcolorbox}
  \begin{tcolorbox}[enhanced,colback=white,colframe=popink,boxrule=2.2pt,arc=10pt,
      drop shadow=popink,left=16pt,right=16pt,top=10pt,bottom=10pt,after skip=8pt]
    \poppill{Dans cette leçon}{poppurple}{white}\par\vspace{5pt}
    \popsemi\fontsize{9.3}{12.6}\selectfont\color{popink}#7
  \end{tcolorbox}
  \noindent\begin{minipage}[t]{0.487\linewidth}
    \begin{tcolorbox}[enhanced,colback=popgreen,colframe=popink,boxrule=2.2pt,arc=10pt,
        drop shadow=popink,left=11pt,right=11pt,top=10pt,bottom=10pt,height=3.1cm,valign=center]
      \tcbox[colback=white,colframe=popink,boxrule=1.3pt,arc=5pt,boxsep=0pt,nobeforeafter,
        left=3pt,right=3pt,top=3pt,bottom=3pt,valign=center]{\qrcode[height=1.9cm]{https://play.google.com/store/apps/details?id=com.luvvix.onbuch&hl=fr}}%
      \hspace{8pt}\begin{minipage}[c]{0.5\linewidth}
        {\popdisplay\fontsize{11}{12.5}\selectfont\color{white}\MakeUppercase{Télécharge OnBuch}}\par\vspace{4pt}
        {\raggedright\popsemi\fontsize{8.4}{11}\selectfont\color{white}Gratuit \textbullet\ Google Play\\Tuteur IA, annales, concours, résultats\par}
      \end{minipage}
    \end{tcolorbox}
  \end{minipage}\hfill
  \begin{minipage}[t]{0.487\linewidth}
    \begin{tcolorbox}[enhanced,colback=poppurple,colframe=popink,boxrule=2.2pt,arc=10pt,
        drop shadow=popink,left=11pt,right=11pt,top=10pt,bottom=10pt,height=3.1cm,valign=center]
      \tikz[baseline=-0.7ex]\node[circle,fill=white,draw=popink,line width=1.3pt,minimum size=1.6cm,inner sep=0]{\popdisplay\fontsize{24}{24}\selectfont\color{poppurple}+};%
      \hspace{8pt}\begin{minipage}[c]{0.6\linewidth}
        {\popdisplay\fontsize{11}{12.5}\selectfont\color{white}\MakeUppercase{Passe à OnBuch+}}\par\vspace{4pt}
        {\raggedright\popsemi\fontsize{8.4}{11}\selectfont\color{white}Tous les cours signés, Tuteur IA illimité, fiches PDF — onglet OnBuch+ de l'app\par}
      \end{minipage}
    \end{tcolorbox}
  \end{minipage}
  \vspace{8pt}
  \popcontenu
  \vfill
  \signatureonbuch
  \restoregeometry
  \newpage
}

% Rangée « Ce cours contient » (page de garde)
\newcommand{\poptile}[3]{% #1 couleur #2 gros texte #3 légende
  \begin{tcolorbox}[enhanced,colback=#1,colframe=popink,boxrule=2pt,arc=9pt,shadow={2.5pt}{-2.5pt}{0pt}{popink},
      left=6pt,right=6pt,top=9pt,bottom=9pt,halign=center,valign=center,height=1.75cm,width=\linewidth,nobeforeafter]
    {\popdisplay\fontsize{12.5}{13}\selectfont\color{white}\MakeUppercase{#2}}\par\vspace{3pt}
    {\centering\popsemi\fontsize{7.8}{9.5}\selectfont\color{white}#3\par}
  \end{tcolorbox}}
\newcommand{\popcontenu}{%
  \par\noindent{\popxboldls\fontsize{7.5}{7.5}\selectfont\color{popmuted}CE COURS SIGNÉ CONTIENT}\par\vspace{7pt}
  \noindent\begin{minipage}[t]{0.235\linewidth}\poptile{popblue}{Cours}{complet, illustré, pas à pas}\end{minipage}\hfill
  \begin{minipage}[t]{0.235\linewidth}\poptile{poporange}{Méthodes}{et exercices résolus}\end{minipage}\hfill
  \begin{minipage}[t]{0.235\linewidth}\poptile{poppink}{Exercices}{type examen + corrigés}\end{minipage}\hfill
  \begin{minipage}[t]{0.235\linewidth}\poptile{popgreen}{Fiche bilan}{l'essentiel à retenir}\end{minipage}\par}

% Sommaire
\newtcolorbox{sommaire}{enhanced,colback=white,colframe=popink,boxrule=2.2pt,arc=10pt,
  drop shadow=popink,left=18pt,right=18pt,top=13pt,bottom=13pt,before skip=6pt,after skip=20pt,
  before upper={\poppill{Sommaire}{poppurple}{white}\par\vspace{9pt}\popsemi\fontsize{10.6}{14.5}\selectfont\color{popink}}}

% Signature de fin de cours — poussée tout en bas de la dernière page
\newcommand{\findecours}{%
  \par\vfill\nopagebreak
  \begin{center}\popsquiggle{poporange}\end{center}\vspace{4pt}
  \signatureonbuch
}
\AtBeginDocument{\def\llbracket{\mathopen{[\mkern-3.5mu[}}\def\rrbracket{\mathclose{]\mkern-3.5mu]}}} % intervalles d'entiers (glyphe ⟦ absent de la police)
% Glyphes absents de Fira Math : redéfinis à partir de glyphes disponibles
\AtBeginDocument{%
  \def\setminus{\mathbin{\backslash}}\let\smallsetminus\setminus
  \def\vdots{\mathord{\vbox{\baselineskip3.2pt\lineskiplimit0pt\kern4pt\hbox{.}\hbox{.}\hbox{.}}}}%
  \def\ddots{\mathinner{\mkern1mu\raise6.5pt\hbox{.}\mkern2mu\raise3.6pt\hbox{.}\mkern2mu\raise0.7pt\hbox{.}\mkern1mu}}%
  \def\triangle{\mathord{\scalebox{0.9}{$\symup{\Delta}$}}}%
  \def\nabla{\mathord{\raisebox{\depth}{\scalebox{1}[-1]{$\symup{\Delta}$}}}}%
  \def\checkmark{\popcheck{popdark}}%
  \def\star{\mathbin{\ast}}\def\bigstar{\mathord{\ast}}%
  \def\diamond{\mathbin{\scalebox{0.6}{$\Diamond$}}}%
  \def\bigcup{\mathop{\mathchoice{\raisebox{-0.3em}{\scalebox{1.7}{$\cup$}}}{\scalebox{1.25}{$\cup$}}{\cup}{\cup}}\displaylimits}%
  \def\bigcap{\mathop{\mathchoice{\raisebox{-0.3em}{\scalebox{1.7}{$\cap$}}}{\scalebox{1.25}{$\cap$}}{\cap}{\cap}}\displaylimits}%
  \def\lesseqqgtr{\mathrel{\lessgtr}}\def\bigoplus{\mathop{\oplus}\displaylimits}%
}
\providecommand{\ssi}{si et seulement si} % abréviation fréquente
\AtBeginDocument{\providecommand{\wideparen}{\overparen}} % arc de cercle

%% FIGFIX-BEGIN v3 — correctifs globaux des figures (docs/onbuch-generation/tools/figfix)
\usepackage{adjustbox}\usepackage{etoolbox}
% toute figure TikZ/pgfplots plus large que la ligne est réduite à la largeur disponible
\BeforeBeginEnvironment{tikzpicture}{\begin{adjustbox}{max width=\linewidth}}
\AfterEndEnvironment{tikzpicture}{\end{adjustbox}}
% légendes pgfplots lisibles par-dessus les courbes
\pgfplotsset{every axis/.append style={legend style={fill=white,fill opacity=0.92,text opacity=1,draw=black!15,inner sep=3pt}}}
% étiquettes d'arêtes (syntaxe edge["..."]) avec fond blanc
\tikzset{every edge quotes/.append style={fill=white,inner sep=1.2pt}}
% bulles de cartes mentales : texte à la ligne dans la bulle, bulle un peu plus grande
\tikzset{every concept/.append style={text width=2.0cm,align=center,minimum size=2.3cm,inner sep=2pt}}
%% FIGFIX-END
\begin{document}

\pagegarde{L’existence et la mort}{Philosophie}{Tle Tech}{Philosophie – classe de Terminale technique}{N11}{4 séances (fiche de progression harmonisée)}{Cette leçon interroge les concepts d'existence et de mort à travers les grands courants philosophiques et les réalités socioculturelles camerounaises. Elle outille l'élève pour argumenter rigoureusement sur le sens de la condition humaine face à sa finitude.
\vspace{4pt}
\begin{multicols}{2}\small
\begin{itemize}
\item À la fin de cette leçon, je sais définir rigoureusement les concepts d'existence, de vie biologique, de mort clinique et de mort sociale.
\item Je sais distinguer l'existence humaine (projet, conscience, liberté) de la simple vie biologique partagée avec l'animal et le végétal.
\item Je sais analyser et comparer les réponses philosophiques majeures (religieuse, existentialiste, absurde, matérialiste) sur le sens de l'existence.
\item Je sais expliquer comment la mort, par sa certitude, structure le sens de l'existence (Heidegger, Sartre, Camus, Épicure).
\item Je sais mobiliser des exemples concrets camerounais (rites funéraires, bioéthique, espérance de vie) pour illustrer la tension entre mort biologique et survie culturelle.
\item Je sais appliquer la méthode de la dissertation philosophique : analyser un sujet, problématiser, construire un plan dialectique et rédiger une introduction et une conclusion conformes aux attendus du Baccalauréat technique.
\end{itemize}\end{multicols}}

\begin{sommaire}
\begin{multicols}{2}
\begin{enumerate}
\item 1. Définir et distinguer : Vie, Existence, Mort
\item 2. Le sens de l'existence : Les grandes réponses philosophiques
\item 3. La mort : Fin radicale ou horizon de sens ?
\item 4. Existence, Mort et Ancestralité : Le cas camerounais (Situations concrètes)
\item 5. Méthodologie : Dissertation et Explication de texte (Savoir-faire rédactionnel)
\item Activité d'intégration
\item Méthodes et exercices résolus
\item Exercices
\item Corrigés des exercices
\item Fiche bilan
\end{enumerate}
\end{multicols}
\end{sommaire}

\begin{objectifs}
\begin{itemize}
\item À la fin de cette leçon, je sais définir rigoureusement les concepts d'existence, de vie biologique, de mort clinique et de mort sociale.
\item Je sais distinguer l'existence humaine (projet, conscience, liberté) de la simple vie biologique partagée avec l'animal et le végétal.
\item Je sais analyser et comparer les réponses philosophiques majeures (religieuse, existentialiste, absurde, matérialiste) sur le sens de l'existence.
\item Je sais expliquer comment la mort, par sa certitude, structure le sens de l'existence (Heidegger, Sartre, Camus, Épicure).
\item Je sais mobiliser des exemples concrets camerounais (rites funéraires, bioéthique, espérance de vie) pour illustrer la tension entre mort biologique et survie culturelle.
\item Je sais appliquer la méthode de la dissertation philosophique : analyser un sujet, problématiser, construire un plan dialectique et rédiger une introduction et une conclusion conformes aux attendus du Baccalauréat technique.
\item Je sais expliquer un texte philosophique en identifiant la thèse, les arguments, les concepts clés et en évaluant la portée de la démonstration.
\item Je sais justifier une position personnelle sur des dilemmes éthiques contemporains (euthanasie, acharnement thérapeutique, don d'organes) en articulant raison philosophique et contexte légal camerounais.
\end{itemize}
\end{objectifs}

\begin{prerequis}
\begin{itemize}
\item La distinction Sujet / Objet et la notion de conscience (classe de Première).
\item La notion de liberté et de responsabilité (début Terminale).
\item Les techniques de base de l'argumentation (thèse, argument, exemple, objection, réfutation).
\item Les repères historiques majeurs : Antiquité (Épicure), Moyen Âge (théocentrisme), Lumières, XIXe-XXe siècles (Existentialisme).
\item La lecture de tableaux statistiques simples (espérance de vie, taux de mortalité).
\end{itemize}
\end{prerequis}

\newpage

\coursec{Définir et distinguer : Vie, Existence, Mort}

\textbf{Situation d'accroche.} Au chevet de son père mourant à l'hôpital Laquintinie de Douala, Marie doit signer un formulaire de « limitation de l'acharnement thérapeutique » pendant que la famille maternelle exige le transfert au village pour les rites de « veillée des ancêtres ». Le médecin annonce : « Votre père est cliniquement mort, le tracé cérébral est plat. » La tante réplique : « Tant que nous n'avons pas accompli les rites, il n'est pas mort, il voyage encore. » Qui a raison ? Entre la mort clinique certifiée par le médecin et la mort sociale rituelle, quelle existence s'éteint vraiment ?

\textbf{Problématique.} Cette scène, familière à beaucoup de familles camerounaises, révèle que les mots « vie », « existence » et « mort » ne recouvrent pas une seule réalité. Vivre, est-ce seulement fonctionner biologiquement ? Exister, est-ce autre chose que vivre ? Et mourir, est-ce un événement ponctuel ou un processus à plusieurs étages ? Telle est la question que cette première section doit clarifier, en distinguant rigoureusement trois plans : le plan \cle{biologique}, le plan \cle{existentiel} et le plan \cle{socio-culturel}.

\courssub{La vie biologique : le fait de fonctionner}

\textbf{Intuition d'abord.} Observez une mangue qui germe, une chèvre qui broute à Bafoussam, un homme qui respire : tous trois partagent quelque chose. Ce quelque chose, c'est la \cle{vie} au sens le plus élémentaire : un ensemble de fonctions qui s'entretiennent elles-mêmes. Nul besoin de philosophie pour la reconnaître ; le biologiste la décrit, le médecin la mesure.

\begin{definition}[Vie biologique]
La \cle{vie biologique} est l'ensemble des fonctions physiologiques — \cle{nutrition}, \cle{respiration}, \cle{métabolisme}, \cle{reproduction}, \cle{excrétion} — communes à tous les êtres organisés, de la bactérie à l'homme. Elle se définit par des critères matériels observables : échanges de matière et d'énergie avec le milieu, homéostasie (maintien des constantes internes), croissance et capacité de reproduction.
\end{definition}

\textbf{Démarche d'investigation.} Comment savoir qu'un organisme est vivant ? La démarche scientifique procède par observation de critères mesurables. Un médecin urgentiste à l'hôpital central de Yaoundé vérifie : le pouls (fonction cardiaque), la respiration, la température, les réflexes pupillaires. Si ces fonctions sont présentes, même de façon artificiellement assistée (respirateur, perfusion), la vie biologique persiste. La vie est donc une \emph{donnée} : elle se constate, elle ne se décide pas.

\begin{exemplebox}[La vie sans conscience]
Un patient plongé dans le coma végétatif après un accident de moto-taxi sur l'axe Yaoundé--Douala continue de respirer, de digérer, de cicatriser. Sa vie biologique est intacte. Pourtant, il ne parle plus, ne projette rien, ne reconnaît personne. Ce cas extrême montre que la vie biologique peut subsister \emph{sans} ce qui fait la spécificité humaine. D'où la nécessité d'un second concept : l'existence.
\end{exemplebox}

\begin{attention}[Piège fréquent]
Ne pas confondre « vie » et « existence » dans une copie. La \cle{vie} est une donnée biologique que l'homme partage avec tous les organismes ; l'\cle{existence} est une tâche proprement humaine. Écrire « l'homme est un être vivant comme les autres » est biologiquement vrai mais philosophiquement insuffisant : c'est effacer la question du sens.
\end{attention}

\courssub{L'existence humaine : une vie qui se comprend elle-même}

\textbf{Intuition d'abord.} Un élève de Terminale technique à Douala se demande : « Que ferai-je après le Bac ? Deviendrai-je mécanicien, infirmier, informaticien ? » Aucun animal ne se pose cette question. Le zèbre vit, mais il n'\emph{existe} pas au sens où il devrait décider de ce qu'il sera. L'homme, lui, est obligé de se construire.

\begin{definition}[Existence]
Étymologiquement, \emph{ex-sistere} signifie « se tenir hors de soi », « se projeter ». L'\cle{existence} est l'acte par lequel un être se déploie dans le temps en se donnant un sens. L'existence humaine est le mode d'être spécifique de l'homme, caractérisé par quatre traits :
\begin{enumerate}
\item la \cle{conscience intentionnelle} : l'homme est toujours conscience \emph{de} quelque chose et conscience \emph{de soi} ;
\item la \cle{temporalité} : l'homme vit tendu vers l'avenir, il fait des \emph{projets} (Heidegger parle du \emph{Dasein}, l'« être-là » qui se soucie de son être) ;
\item la \cle{liberté} : l'homme peut choisir, dire non, transformer sa situation ;
\item la \cle{question du sens} : l'homme est le seul être qui s'interroge sur le pourquoi de sa vie.
\end{enumerate}
\end{definition}

\begin{propriete}[L'existence précède l'essence (Sartre)]
Selon Jean-Paul Sartre (\emph{L'existentialisme est un humanisme}, 1946), « l'existence précède l'essence » : l'homme n'a pas de nature prédéfinie, aucun « plan » ne décide à l'avance de ce qu'il doit être. Il \emph{existe} d'abord, puis il se définit par ses choix et ses actes. Conséquence décisive : la mort, comme limite certaine et finitude, est ce qui rend \emph{urgent} de choisir. Parce que mon temps est compté, mes choix ont du poids. Une vie infinie serait une vie où tout pourrait être remis à demain, donc une vie sans enjeu.
\end{propriete}

\textbf{La distinction fondamentale.} Martin Heidegger (\emph{Être et Temps}, 1927) formule la coupure décisive : \emph{« l'animal vit, l'homme existe »}. L'animal est enfermé dans son milieu et ses instincts ; l'homme, lui, entretient avec son propre être un rapport de \cle{souci} (\emph{Sorge}) : il se comprend lui-même, il se demande ce qu'il est, il peut même mettre en question sa vie. L'existence est donc \textbf{une vie qui se comprend elle-même et se soucie de son être}. C'est pourquoi l'on peut dire, sans paradoxe, qu'un homme « existe à peine » (quand il subit sa vie sans la choisir) ou qu'il « existe pleinement » (quand il assume ses projets).

\begin{exemplebox}[Exister ou simplement survivre]
Deux jeunes diplômés chômeurs à Bafoussam : l'un subit ses journées devant la télévision en attendant « que ça vienne » ; l'autre, dans la même précarité, monte un atelier de réparation de téléphones, échoue, recommence. Biologiquement, leurs vies sont identiques. Existentiellement, elles diffèrent : le second fait de sa vie un \emph{projet}, il se « tient hors de soi » vers un avenir qu'il veut construire. L'existence n'est pas un luxe de riche : c'est un rapport à soi.
\end{exemplebox}

\courssub{La mort biologique : critères cliniques}

\textbf{Intuition d'abord.} Si la vie est un fonctionnement, la mort biologique est l'arrêt de ce fonctionnement. Mais arrêt de quoi, exactement ? Le cœur peut être relancé ; la respiration peut être assistée par une machine. La médecine moderne a dû préciser ses critères.

\begin{definition}[Mort clinique et mort encéphalique]
La \cle{mort clinique} (ou mort cardiaque) est la cessation des fonctions cardiaque et respiratoire ; elle peut être \emph{réversible} dans un délai court (réanimation). La \cle{mort encéphalique} (mort du tronc cérébral) est la destruction irréversible de l'ensemble du cerveau, constatée par des critères stricts : coma absolu, abolition des réflexes du tronc cérébral, absence de respiration spontanée, confirmée par deux électroencéphalogrammes plats ou une angiographie. C'est le critère légal du décès retenu par l'OMS et appliqué au Cameroun (cadre réglementaire encadrant le prélèvement d'organes, notamment l'arrêté ministériel de 2012) : un patient en mort encéphalique, même si une machine maintient artificiellement son cœur battant, est légalement \emph{décédé}.
\end{definition}

\textbf{Pourquoi ce critère ?} Parce que le cerveau est le centre intégrateur de toutes les fonctions vitales. Sans tronc cérébral, l'organisme n'est plus une unité vivante mais une collection d'organes maintenus artificiellement. La mort encéphalique est donc la mort de l'organisme \emph{comme tout}, même si des cellules (peau, cheveux) survivent encore quelques heures. La mort biologique n'est pas un instant mais un \emph{processus} : mort de l'organisme, puis mort des organes, puis mort des cellules.

\begin{attention}[Erreur fréquente en dissertation]
Ne pas écrire « le cœur s'arrête donc l'homme est mort » sans nuance. Un arrêt cardiaque peut être réanimé ; un cœur peut battre sous assistance chez un patient en mort encéphalique. Le critère décisif, médical et légal, est la \cle{mort encéphalique}. Citer ce critère dans une copie montre la maîtrise de la distinction vie/mort.
\end{attention}

\courssub{La mort sociale et culturelle : le groupe et le défunt}

\textbf{Intuition d'abord.} Revenons à Laquintinie. Le médecin a signé le certificat de décès : biologiquement, le père de Marie est mort. Pourtant la famille dit : « Il n'est pas encore parti. » De quoi parle-t-elle ? D'une autre mort, que ni le stéthoscope ni l'électroencéphalogramme ne peuvent mesurer.

\begin{definition}[Mort sociale et culturelle]
La \cle{mort sociale} est le processus par lequel le groupe retire au défunt son statut de personne vivante et le réintègre dans un nouveau statut (ancêtre, esprit, souvenir). Elle s'accomplit par des \cle{rites} : veillée, deuil, funérailles, héritage, partage des biens, remariage éventuel de la veuve. Tant que les rites ne sont pas accomplis, le défunt occupe une position liminale : il n'est plus vivant, mais pas encore pleinement « mort » pour la communauté. La mort sociale peut donc \emph{suivre} la mort biologique (funérailles différées, « célébration » des mois après, comme dans les funérailles bamiléké ou les \emph{ngondo} sawa) — et dans certains cas extrêmes la \emph{précéder} (exclusion, bannissement d'un vivant traité comme mort par son groupe).
\end{definition}

\begin{savaistu}[Le cadavre en droit camerounais]
En droit camerounais (droit civil hérité du droit français et droit coutumier), la \cle{personnalité juridique} commence à la naissance de l'enfant « vivant et viable » et cesse à la mort. Mais le cadavre n'est pas une « chose » comme les autres : il bénéficie d'une protection spécifique (respect de la sépulture, sanctions pénales contre la violation de tombe). Le droit reconnaît ainsi, à sa manière, que l'homme mort n'est pas réductible à de la matière : il reste pris dans un tissu de respect social. C'est la trace juridique de la mort sociale.
\end{savaistu}

\begin{exemplebox}[Données chiffrées : la précarité de l'existence concrète]
Selon les données récentes de l'INS (projections RGPH et EDS-MICS), l'espérance de vie à la naissance au Cameroun est d'environ \num{60,6} ans (hommes : \num{59,2} ans ; femmes : \num{62,1} ans), et le taux de mortalité infantile avoisine \num{48} pour mille, soit près d'un enfant sur vingt qui meurt avant cinq ans. Ces chiffres ne sont pas de simples statistiques : ils situent la \emph{finitude concrète} de l'existence camerounaise. Philosophiquement, ils rappellent que la question du sens de l'existence n'est pas abstraite : elle se pose à des êtres dont le temps est objectivement limité et inégalement distribué.
\end{exemplebox}

\courssub{Synthèse : trois plans, trois seuils}

Les trois notions se combinent sans se confondre. On peut être biologiquement vivant sans exister pleinement (coma végétatif) ; on peut être biologiquement mort sans être encore socialement mort (rites non accomplis) ; on peut être socialement « mort » tout en étant biologiquement vivant (exil, exclusion). Le schéma suivant visualise ces imbrications.

\begin{popfigure}
\begin{center}
\begin{tikzpicture}[scale=0.95]
% Grand cercle : vie biologique
\draw[fill=popgreenL,draw=popgreen,thick] (0,0) circle (3.1);
\node[popdark,font=\bfseries] at (0,2.55) {Vie biologique};
% Cercle existence
\draw[fill=popblueL,draw=popblue,thick] (-0.5,0.2) circle (1.9);
\node[popdark,font=\bfseries] at (-0.5,1.55) {Existence humaine};
% Cercle personne morale (chevauchant)
\draw[fill=poporangeL,draw=poporange,thick,fill opacity=0.55] (1.15,-0.3) circle (1.25);
\node[popdark,font=\bfseries,align=center] at (1.35,-0.35) {Personne\\morale/juridique};
% Flèches annotées
\draw[->,very thick,popink] (-4.6,0.2) -- (-3.15,0.2) node[midway,above,font=\small\itshape]{Naissance};
\draw[->,very thick,popink] (3.15,1.3) -- (4.6,1.3) node[midway,above,font=\small\itshape]{Mort clinique};
\draw[->,very thick,popink,dashed] (3.15,-1.3) -- (4.6,-1.3) node[midway,below,font=\small\itshape]{Mort sociale};
\end{tikzpicture}
\end{center}
\legende{Diagramme des trois plans. Le grand cercle (vert) : la vie biologique, commune à tous les organismes. Le cercle bleu : l'existence humaine, incluse dans la vie mais dotée de conscience, de projet et de liberté. Le cercle orange : la personne juridique et morale, qui chevauche l'existence (elle naît avec l'enfant « vivant et viable » et survit partiellement au corps par la mémoire, l'héritage, les rites). Les flèches marquent les seuils : naissance, mort clinique, mort sociale.}
\end{popfigure}

Le schéma heuristique suivant permet de traiter les cas limites qui tombent régulièrement au Baccalauréat technique : l'embryon, le grand brûlé inconscient, le patient en état végétatif.

\begin{popfigure}
\begin{center}
\begin{tikzpicture}[scale=1.0]
% Axes
\draw[->,thick,popdark] (0,0) -- (8.6,0) node[below left,font=\small]{Vie biologique $\rightarrow$};
\draw[->,thick,popdark] (0,0) -- (0,5.4) node[above,font=\small]{Conscience / Reconnaissance sociale $\rightarrow$};
% Points
\fill[popgreen] (1.2,0.9) circle (3pt) node[right,align=left,font=\small]{Embryon : vie forte,\\conscience nulle,\\reconnaissance débattue};
\fill[popblue] (6.2,4.4) circle (3pt) node[right,align=left,font=\small]{Homme éveillé : vie,\\conscience, reconnaissance};
\fill[poporange] (6.8,1.1) circle (3pt) node[right,align=left,font=\small]{Patient végétatif : vie maintenue,\\conscience absente,\\reconnaissance affective};
\fill[poppurple] (4.6,2.6) circle (3pt) node[right,align=left,font=\small]{Grand brûlé sous coma artificiel :\\existence \emph{suspendue}};
\fill[popink] (7.6,0.5) circle (3pt) node[right,align=left,font=\small]{Mort encéphalique :\\corps perfusé, personne décédée};
\end{tikzpicture}
\end{center}
\legende{Schéma heuristique des cas limites. En abscisse, l'intensité de la vie biologique ; en ordonnée, le degré de conscience et de reconnaissance sociale. Chaque cas clinique occupe une position différente : c'est la preuve que « vivant », « existant » et « personne » ne sont pas synonymes.}
\end{popfigure}

\begin{aretenir}[L'essentiel de la section]
\begin{itemize}
\item La \cle{vie} est une donnée biologique (fonctions physiologiques) commune à tous les organismes.
\item L'\cle{existence} est le mode d'être proprement humain : conscience de soi, temporalité, projet, liberté, question du sens. « L'animal vit, l'homme existe » (Heidegger).
\item « L'existence précède l'essence » (Sartre) : l'homme se définit par ses choix ; la mort, comme limite, rend le choix urgent.
\item La \cle{mort biologique} se définit légalement par la mort encéphalique, non par le seul arrêt cardiaque.
\item La \cle{mort sociale} est un processus rituel (deuil, funérailles, héritage) qui, au Cameroun, suit souvent la mort biologique et donne au défunt un nouveau statut (ancêtre).
\end{itemize}
\end{aretenir}

\begin{exoresolu}[Appliquer la distinction à un cas concret]
\textbf{Énoncé.} À l'hôpital Laquintinie, le père de Marie est déclaré en mort encéphalique ; les machines maintiennent son cœur battant. La famille refuse de le considérer comme mort avant les rites du village. En vous appuyant sur les distinctions de la leçon, analysez ce conflit en cinq à huit lignes : qui « meurt », et à quel plan ?
\tcblower
\textbf{Solution détaillée.} Le conflit oppose deux plans de réalité que la leçon a distingués. Sur le plan \emph{biologique}, le médecin a raison : la mort encéphalique, critère légal du décès, signifie que l'organisme comme tout est irréversiblement détruit ; le cœur battant n'est qu'une survie artificielle d'organes. Sur le plan \emph{existentiel}, l'existence du père — conscience, projet, liberté — est déjà éteinte : il ne se « tient plus hors de soi » vers aucun avenir. Mais sur le plan \emph{socio-culturel}, la famille a également raison à son niveau : pour elle, la mort n'est pas un instant clinique mais un processus qui exige les rites pour que le défunt change de statut et devienne ancêtre. La solution philosophique n'est donc pas de départager un « vrai » et un « faux », mais de reconnaître que la mort humaine est un événement \emph{à plusieurs dimensions} : le corps meurt une fois, la personne sociale meurt par étapes. Le drame de Marie vient précisément de ce décalage temporel entre les plans.
\end{exoresolu}

\begin{exercice}[Pour s'entraîner]
Un journal titre : « Grâce aux machines, cet homme est resté en vie trois ans après sa mort. » Relevez les confusions conceptuelles de ce titre et reformulez-le correctement en distinguant vie biologique, existence et mort encéphalique. (Savoir-faire visé : appliquer les notions de l'unité à une situation concrète et justifier sa réponse.)
\end{exercice}

\begin{corrige}[Exercice 1]
Le titre confond trois plans. (1) « En vie » : seules des fonctions biologiques partielles (cœur, poumons assistés) ont été maintenues artificiellement ; il ne s'agit pas de la vie de l'organisme comme tout. (2) « Après sa mort » : si la mort encéphalique a été constatée, l'homme était légalement et médicalement décédé ; ce sont des \emph{organes} qui ont survécu, non la personne. (3) L'« existence » au sens philosophique (conscience, projet) était éteinte depuis la destruction cérébrale. Reformulation correcte : « Grâce aux machines, les fonctions cardiaque et respiratoire de ce patient en mort encéphalique ont été artificiellement maintenues trois ans après le constat légal de son décès. » La rigueur du vocabulaire est ici elle-même un acte philosophique.
\end{corrige}

\coursec{Le sens de l'existence : Les grandes réponses philosophiques}

\courssub{Introduction : Une question qui engage toute une vie}

Après avoir distingué la \cle{vie} (phénomène biologique universel), l'\cle{existence} (manière singulière, consciente et libre dont l'homme habite le monde) et la \cle{mort} (cessation de la vie biologique, mais aussi horizon qui structure l'existence), nous abordons maintenant le cœur du problème : \textbf{questa vie a-t-elle un sens, et si oui, comment le découvrir ou le construire ?}

Cette question n'est pas une curiosité intellectuelle : elle guide nos choix éthiques, notre rapport à la souffrance, notre engagement professionnel et notre manière d'habiter la société camerounaise. Comme l'écrivait \textbf{Victor Frankl}, rescapé des camps, « celui qui a un \emph{pourquoi} pour vivre peut supporter presque n'importe quel \emph{comment} ». Les réponses philosophiques se répartissent en grandes familles selon l'origine qu'elles attribuent au sens : \emph{donné} (transcendance), \emph{découvert} (raison), \emph{inventé} (liberté), \emph{absent} (absurde), \emph{nié} (nihilisme), \emph{biologique} (matérialisme) ou \emph{relationnel} (éthique).

\begin{popfigure}
\begin{tikzpicture}[
    timeline/.style={very thick, popblue},
    era/.style={font=\small\bfseries, text=popdark},
    bubble/.style={draw, thick, rounded corners=4pt, fill=popcream, text width=2.8cm, align=center, font=\small, inner sep=3pt},
    arrow/.style={-Stealth, thick, popline},
    pinline/.style={thin, popline, dashed}
]
% Axe principal
\draw[timeline] (0,0) -- (16,0) node[right, era] {Aujourd'hui};

% Repères temporels
\foreach \x/\label in {0/Antiquité, 4/Moyen Âge, 8/Moderne, 12/Contemporain, 16/Actuel} {
    \draw (\x,0.15) -- (\x,-0.15) node[below, era, yshift=-4pt] {\label};
}

% Bulles philosophiques : (position x, y, couleur bordure, texte)
% Épicure
\node[bubble, draw=popgreen!80!black, align=center] (epi) at (1, 2.5){\textbf{Épicure} (341--270 av. J.-C.)\\ Plaisir modéré / Ataraxie\\ \textit{Matérialisme}};
\draw[pinline] (1,0) -- (epi.south);

% Augustin
\node[bubble, draw=poppurple!80!black, align=center] (aug) at (4, -2.5){\textbf{St Augustin} (354--430)\\ Sens donné par Dieu\\ \textit{Théocentrisme}};
\draw[pinline] (4,0) -- (aug.north);

% Descartes
\node[bubble, draw=popblue!80!black, align=center] (des) at (7, 2.5){\textbf{Descartes} (1596--1650)\\ Raison / Vérité / Vertu\\ \textit{Rationalisme classique}};
\draw[pinline] (7,0) -- (des.south);

% Nietzsche
\node[bubble, draw=poporange!80!black, align=center] (nie) at (9.5, -2.5){\textbf{Nietzsche} (1844--1900)\\ Mort de Dieu / Surhomme\\ \textit{Nihilisme actif}};
\draw[pinline] (9.5,0) -- (nie.north);

% Heidegger
\node[bubble, draw=poppink!80!black, align=center] (hei) at (11.5, 2.5){\textbf{Heidegger} (1889--1976)\\ Être-pour-la-mort\\ \textit{Existentialisme ontologique}};
\draw[pinline] (11.5,0) -- (hei.south);

% Sartre
\node[bubble, draw=popink!80!black, align=center] (sar) at (13, -2.5){\textbf{Sartre} (1905--1980)\\ Existence précède essence\\ Engagement / Liberté radicale};
\draw[pinline] (13,0) -- (sar.north);

% Camus
\node[bubble, draw=popgold!80!black, align=center] (cam) at (14, 2.5){\textbf{Camus} (1913--1960)\\ Absurde / Révolte / Passion\\ \textit{Mythe de Sisyphe}};
\draw[pinline] (14,0) -- (cam.south);

% Ricœur
\node[bubble, draw=popgreen!60!black, align=center] (ric) at (15.5, -1.8){\textbf{Ricœur} (1913--2005)\\ Éthique / Narration\\ Responsabilité / Autre};
\draw[pinline] (15.5,0) -- (ric.north);

% Légendes familles
\node[anchor=west, font=\footnotesize, text=popmuted] at (16.2, 2.5) {\textcolor{popgreen}{Matérialisme / Sagesse antique}};
\node[anchor=west, font=\footnotesize, text=popmuted] at (16.2, 1.8) {\textcolor{poppurple}{Théocentrisme / Révélation}};
\node[anchor=west, font=\footnotesize, text=popmuted] at (16.2, 1.1) {\textcolor{popblue}{Rationalisme / Sujet cognoscitif}};
\node[anchor=west, font=\footnotesize, text=popmuted] at (16.2, 0.4) {\textcolor{poporange}{Critique des valeurs / Nihilisme}};
\node[anchor=west, font=\footnotesize, text=popmuted] at (16.2, -0.3) {\textcolor{poppink}{Analytique existentielle / Mort}};
\node[anchor=west, font=\footnotesize, text=popmuted] at (16.2, -1.0) {\textcolor{popink}{Existentialisme athée / Liberté}};
\node[anchor=west, font=\footnotesize, text=popmuted] at (16.2, -1.7) {\textcolor{popgold}{Philosophie de l'absurde / Révolte}};
\node[anchor=west, font=\footnotesize, text=popmuted] at (16.2, -2.4) {\textcolor{popgreen!60!black}{Humanisme laïque / Éthique}};
\end{tikzpicture}
\legende{Frise chronologique des grandes réponses philosophiques sur le sens de l'existence. Chaque bulle situe un auteur majeur dans son époque et sa famille de pensée (code couleur).}
\end{popfigure}

\courssub{L'approche religieuse et théocentrique : Le sens comme don et vocation}

\begin{definition}[Sens théocentrique de l'existence]
Dans la perspective religieuse (christianisme, islam, judaïsme, religions traditionnelles africaines), le sens de l'existence humaine n'est pas une invention humaine : il est \textbf{donné \emph{a priori}} par le Créateur (Dieu, \cle{Nzambé}, \cle{Roog}, les Ancêtres). L'homme est créé \emph{intentionnellement} pour une fin : connaître, aimer et servir Dieu ; participer à l'harmonie cosmique ; accomplir sa destinée (\cle{chi} dans la pensée igbo, \cle{nommo} chez les Dogon). La vie est un \emph{don}, une \emph{épreuve} ou un \emph{passage} vers l'éternité.
\end{definition}

\paragraph{Structure de l'argumentation.}
\begin{enumerate}
    \item \textbf{Origine :} L'homme ne se fait pas lui-même ; il est \emph{créature}. Son essence (sa nature, sa finalité) précède son existence (Aristote / St Thomas d'Aquin).
    \item \textbf{Nature de la vie :} Elle est sacrée, dépositaire d'un souffle divin. La souffrance et la mort ne sont pas des absurdités mais des mystères (épreuve de foi, purification, retour à la source).
    \item \textbf{Horizon :} La mort n'est pas une fin radicale mais un \emph{passage} (jugement, réincarnation, monde des ancêtres). Elle donne urgence et gravité à la vie présente.
\end{enumerate}

\begin{exemplebox}[Traditions africaines : L'existence comme alliance avec les Ancêtres]
Au Cameroun, chez les \textbf{Bamiléké}, les \textbf{Bassa}, les \textbf{Beti} ou les \textbf{Peuls}, l'existence s'inscrit dans une \cle{triade ontologique} : \emph{Vivants -- Ancêtres -- Divinités / Forces vitales}.
\begin{itemize}
    \item Naître, c'est recevoir une \emph{force vitale} des ancêtres ; on a une \cle{dette} envers eux (respect des rites, funérailles dignes, transmission de la vie).
    \item Le sens de l'existence : \textbf{maintenir l'harmonie} (paix sociale, fertilité, santé) par le respect des interdits (\cle{totems}, \cle{jours de marché}, rites agraires).
    \item La mort : « \emph{Aller rejoindre les ancêtres} ». Les funérailles (\emph{tsi}, \emph{ngondo}, \emph{ngoun}) ne sont pas un adieu mais une \emph{intronisation} du défunt comme ancêtre protecteur. L'existence trouve son accomplissement dans la \emph{mémoire} et la \emph{postérité}.
\end{itemize}
\end{exemplebox}

\begin{attention}[Piège fréquent : Confondre fatalisme et vocation]
Dire « c'est la volonté de Dieu / des ancêtres » ne signifie pas \emph{passivité} ou \emph{résignation béate}. Dans le christianisme comme dans les traditions africaines, l'homme est \textbf{co-responsable} : il doit \emph{discerner} sa vocation (prière, divination, conseil des sages) et \emph{y répondre} par l'action (travail, charité, justice). Le sens est \emph{donné}, mais il faut \emph{l'accueillir librement}.
\end{attention}

\courssub{L'approche rationaliste et classique : Le sens découvert par la raison}

\begin{definition}[Rationalisme classique (Descartes, Leibniz, Wolff)]
Pour le rationalisme des XVII\textsuperscript{e}--XVIII\textsuperscript{e} siècles, le monde est un \emph{cosmos} ordonné, intelligible, créé par un Dieu géomètre. La raison humaine, participante de la Raison divine, peut \textbf{découvrir} l'ordre du monde et s'y conformer. Le sens de l'existence réside dans la \cle{connaissance de la vérité} (sciences, métaphysique) et la \cle{pratique de la vertu} (sagesse, maîtrise des passions).
\end{definition}

\paragraph{Descartes (\emph{Discours de la méthode}, \emph{Passions de l'âme}).}
\begin{itemize}
    \item \textbf{Dualisme substanceel :} \cle{res cogitans} (âme, pensée) et \cle{res extensa} (corps, étendue). La mort est \emph{séparation} de l'âme et du corps ; l'âme est immortelle par nature.
    \item \textbf{Morale par provision :} En attendant la vérité certaine, obéir aux lois/religion du pays ; être ferme dans ses résolutions ; se conquérir soi-même plutôt que la fortune.
    \item \textbf{Sagesse finale :} La générosité (libre volonté de bien user de sa liberté) et la connaissance des vérités premières procurent une \emph{joie} supérieure aux plaisirs sensibles.
\end{itemize}

\paragraph{Leibniz (\emph{Théodicée}, \emph{Monadologie}).}
\begin{itemize}
    \item \textbf{Meilleur des mondes possibles :} Dieu a choisi le monde où le rapport \emph{variété / ordre} est maximal. Le mal (métaphysique, physique, moral) est nécessaire au bien supérieur.
    \item \textbf{Sens :} \emph{Amour de Dieu} = joie dans la perfection des autres. L'existence humaine vise la \emph{félicité} par la connaissance progressive de l'harmonie universelle.
\end{itemize}

\begin{aretenir}[Point commun et différence avec le religieux]
\textbf{Commun :} Le sens est \emph{objectif}, \emph{universel}, \emph{antérieur} à mon choix subjectif.
\textbf{Différence :} Le rationalisme ne demande pas la \emph{foi} (révélation) mais l'\emph{exercice de la raison} (démonstration, méthode). La mort devient un problème \emph{métaphysique} (séparation) plus qu'un mystère \emph{sacramentel}.
\end{aretenir}

\courssub{L'existentialisme athée : L'existence précède l'essence — L'invention du sens}

\begin{propriete}[Thèse fondamentale de l'existentialisme (Sartre, \emph{L'Existentialisme est un humanisme}, 1946)]
\emph{« L'existence précède l'essence. »} Cela signifie : l'homme \textbf{apparaît d'abord sur la scène du monde} (existence), se rencontre, se définit \emph{après coup} par ses choix (essence). Il n'y a \textbf{pas de nature humaine prédéfinie} (ni par Dieu, ni par la biologie, ni par la société). L'homme \emph{est} ce qu'il \emph{fait} de lui-même.
\end{propriete}

\paragraph{Conséquences anthropologiques majeures (Sartre).}
\begin{enumerate}
    \item \textbf{Condamné à être libre :} Pas de déterminisme (Dieu, instincts, classe sociale) qui m'excuse. Je suis \emph{totalement responsable} de ce que je suis.
    \item \textbf{Angoisse :} Je choisis \emph{pour tous les hommes} (universalisation implicite de mon choix). Choisir d'être prof, c'est dire « c'est bien d'être prof ».
    \item \textbf{Mauvaise foi :} Fuir sa liberté en se faisant \emph{objet} (ex : « Je suis \emph{né} timide », « Je \emph{dois} obéir aux ordres »).
    \item \textbf{Engagement :} Le sens naît de l'\emph{action} projetée vers l'avenir (\cle{projet fondamental}). « L'homme n'est rien d'autre que son projet. »
\end{enumerate}

\paragraph{Heidegger (\emph{Être et Temps}, 1927) : L'analytique du \emph{Dasein}.}
\begin{itemize}
    \item L'homme n'est pas un « sujet » face à des « objets », mais \cle{Dasein} (l'être-là : être-au-monde, être-avec, être-pour-la-mort).
    \item \textbf{Le « On » (\emph{Das Man}) :} Mode d'existence \emph{inauthentique} : on parle comme « on » parle, on pense comme « on » pense, on fuit la mort en la banalisant (« tout le monde meurt »).
    \item \textbf{L'anticipation de la mort (\emph{Vorlaufen}) :} La mort est ma \emph{possibilité la plus propre, infranchissable, certaine, indéterminée}. L'assumer (\cle{résolution}) m'arrache au « On » pour une existence \emph{authentique}.
\end{itemize}

\begin{propriete}[Théorème / Thèse clé : Heidegger — La mort comme « possibilité de l'impossibilité absolue »]
La mort n'est pas un événement futur parmi d'autres (comme un examen), mais la \textbf{possibilité de ne plus pouvoir être là}. Elle est \emph{impossibilité absolue} de toute possibilité. L'\textbf{anticipation résolue} de cette fin (non pas peur panique, mais lucidité sereine) \textbf{individualise} le Dasein : personne ne peut mourir à ma place. Cette solitude face à la mort fonde l'\textbf{authenticité} : je deviens \emph{ma propre loi} au lieu de subir la loi du « On ».
\end{propriete}

\begin{savaistu}[Sartre et la définition classique]
La formule « L'existence précède l'essence » (1946) renverse explicitement la métaphysique classique (Aristote, St Thomas) : pour celle-ci, l'\emph{essence} (la définition de ce qu'est l'homme : « animal raisonnable ») précède l'\emph{existence} (la réalisation concrète). Sartre dit : \emph{il n'y a pas de nature humaine}, il n'y a que des \emph{histoires humaines}.
\end{savaistu}

\courssub{La philosophie de l'Absurde : Vivre sans appel (Camus)}

\begin{definition}[L'Absurde chez Camus (\emph{Le Mythe de Sisyphe}, 1942)]
L'Absurde n'est \textbf{ni} dans l'homme (sa soif de sens) \textbf{ni} dans le monde (son silence), mais dans \textbf{leur confrontation} : le \emph{divorce} entre l'acteur qui crie « Pourquoi ? » et le décor qui ne répond pas. C'est une \cle{tension}, non une thèse nihiliste.
\end{definition}

\paragraph{Les trois conséquences de l'Absurde (Camus).}
\begin{enumerate}
    \item \textbf{La Révolte :} Refus lucide de se suicider (physiquement ou philosophiquement / espoir religieux). « Je me révolte, donc nous existons. » C'est une \emph{valeur} : la dignité humaine.
    \item \textbf{La Liberté :} Puisque pas d'au-delà, pas de maître, pas de destin écrit. Je suis libre de \emph{vivre maintenant}, intensément, \emph{sans appel} à l'éternité.
    \item \textbf{La Passion (Quantité d'expériences) :} Non pas « qualité » (sainteté, héroïsme), mais \emph{quantité} : multiplier les vies, les amours, les voyages, les soleils. « Il n'y a pas de destin qui ne se surmonte par le mépris. »
\end{enumerate}

\paragraph{Le Mythe de Sisyphe.}
Sisyphe, condamné par les dieux à rouler un rocher au sommet d'une montagne pour le voir redescendre éternellement, est le \emph{héros absurde}. Sa torture : la \emph{lucidité} pendant la redescente. Mais \textbf{cette lucidité est sa victoire} : il est supérieur à son destin. \emph{« Il faut imaginer Sisyphe heureux. »} Le sens n'est pas \emph{donné} (Dieu), ni \emph{construit} (projet sartrien), mais \emph{vécu} dans la révolte lucide.

\begin{attention}[Piège majeur : Nihilisme vs Absurde]
\begin{itemize}
    \item \textbf{Nihilisme (Nietzsche passif, ou sens commun) :} « Rien n'a de sens, donc tout est permis / tout est vain / je meurs. » \emph{Négation} des valeurs.
    \item \textbf{Absurde (Camus) :} « Le monde n'a pas de sens \emph{donné}, \emph{mais} ma soif de sens est indestructible. \emph{Cette tension même} fonde une éthique : révolte, liberté, passion, solidarité (les hommes unis par le même sort). » \emph{Affirmation} de la vie \emph{dans} son absence de sens ultime.
\end{itemize}
\emph{Ne pas écrire : « Camus dit que la vie n'a pas de sens. » Écrire : « Camus dit que la vie n'a pas de sens \emph{transcendant}, mais que cette absence fonde une éthique de l'intensité et de la fraternité. »}
\end{attention}

\courssub{L'approche matérialiste et scientifique : La mort n'est rien pour nous}

\begin{definition}[Matérialisme antique (Épicure, Lucrèce) et biologisme contemporain]
L'âme n'est pas une substance spirituelle immortelle, mais un \emph{arrangement d'atomes} (Épicure) ou un \emph{produit de la complexité neuronale} (neurosciences). La conscience \emph{césse} avec la désorganisation de la matière. \textbf{La mort n'est pas un malheur pour le sujet} : « Quand nous sommes, la mort n'est pas ; quand la mort est, nous ne sommes plus » (\emph{Lettre à Ménécée}).
\end{definition}

\paragraph{Épicure : La sagesse comme ataraxie (absence de trouble).}
\begin{itemize}
    \item \textbf{Physique :} Atomes + vide. Âme = atomes fins répartis dans le corps. Mort = dispersion. Pas de survie, pas de châtiments.
    \item \textbf{Éthique :} Le bien = plaisir (\emph{hédonisme}), mais plaisir \emph{statique} (absence de douleur / \emph{ataraxie}) > plaisir \emph{cinétique} (mouvement, excès). Craindre la mort trouble l'ataraxie \emph{maintenant}. La science (physique) libère de la superstition (dieux, enfer).
\end{itemize}

\paragraph{Biologisme / Naturalisme contemporain (Dawkins, Changeux, etc.).}
\begin{itemize}
    \item L'homme est une \emph{machine de survie} pour ses gènes (théorie de l'évolution). Conscience = épiphénomène / fonction adaptative.
    \item Sens biologique : \emph{transmettre l'information génétique}, assurer la cohésion du groupe (altruisme réciproque).
    \item \textbf{Limite philosophique :} Explique \emph{comment} nous existons (causes efficientes), mais ne dit pas \emph{pourquoi} vivre \emph{cette} vie \emph{ici-maintenant} (finalité subjective).
\end{itemize}

\begin{exemplebox}[Matérialisme et soin palliatif au Cameroun]
Un médecin matérialiste convaincu, face à un patient en fin de vie, ne nie pas la souffrance. Il dit : « La mort est l'arrêt des fonctions cérébrales. Mon devoir (éthique laïque) est de \emph{soulager la douleur} (morphine, sédation) et d'\emph{accompagner} la personne dans sa lucidité, sans lui mentir sur un au-delà, mais en respectant sa dignité d'être humain qui \emph{a vécu}. » Le sens : \emph{réduire l'entropie souffrante} par la science et la compassion.
\end{exemplebox}

\courssub{L'humanisme laïque : Le sens naît de la relation à l'Autre (Ricœur, Levinas)}

\begin{definition}[Humanisme éthique et herméneutique (Paul Ricœur, Emmanuel Levinas)]
Sans postuler Dieu (théocentrisme) ni proclamer la liberté absolue (Sartre), cet courant situe le sens dans \textbf{l'intersubjectivité} : je deviens \emph{quelqu'un} (identité narrative) en répondant à \emph{quelqu'un d'autre} (responsabilité éthique).
\end{definition}

\paragraph{Levinas (\emph{Totalité et Infini}) : L'éthique comme philosophie première.}
\begin{itemize}
    \item Le \emph{Visage} de l'Autre (nu, démuni) m'adresse un commandement muet : « \emph{Tu ne tueras pas} », « \emph{Prends soin de moi} ».
    \item Je suis \emph{responsable} \emph{avant} d'avoir choisi (responsabilité \emph{pré-originale}). Le sens = \emph{réponse} à l'appel de l'Autre. L'existence est \emph{substitution} (porter la peine de l'autre).
\end{itemize}

\paragraph{Ricœur (\emph{Soi-même comme un autre}) : Identités narrative et éthique.}
\begin{itemize}
    \item \textbf{Identité narrative :} Je me raconte une histoire pour donner de la cohérence à ma vie fragmentée. Le sens = \emph{intelligibilité} de mon parcours (emploi, famille, engagements).
    \item \textbf{Attestation :} Je \emph{promets}, je \emph{témoigne}. La parole donnée fonde la confiance sociale.
    \item \textbf{Capacité / Vulnérabilité :} L'homme est \emph{capable} (agir, parler, raconter) et \emph{vulnérable} (souffrir, mourir). La justice et les institutions protègent la vulnérabilité.
\end{itemize}

\begin{aretenir}[Synthèse : Six réponses, un même questionnement]
\begin{center}
\begin{tabularx}{\linewidth}{|l|X|X|X|}\hline
\textbf{Approche} & \textbf{Origine du sens} & \textbf{Rôle de la liberté} & \textbf{Statut de la mort} \\\hline
Religieuse / Africaine & Donné (Dieu / Ancêtres) & Accueillir / Discerner vocation & Passage, jugement, ancêtres \\\hline
Rationaliste (Descartes) & Découvert (Raison / Ordre) & Maîtrise de soi / Vertu & Séparation âme/corps \\\hline
Existentialiste (Sartre) & Inventé (Projet / Engagement) & Absolue (Condamné à l'être) & Fin radicale qui donne urgence \\\hline
Absurde (Camus) & Absent (Silence du monde) & Révolte / Passion (Quantité) & Fin absurde, lucidité héroïque \\\hline
Matérialiste (Épicure) & Biologique / Plaisir modéré & Sagesse (Ataraxie) & Néant, « rien pour nous » \\\hline
Humaniste laïque (Ricœur) & Relationnel (Autre / Récit) & Responsabilité / Promesse & Fin qui appelle témoignage \\\hline
\end{tabularx}
\end{center}
\end{aretenir}

\courssub{Méthodologie : Comparer deux thèses sur un tableau à double entrée (Savoir-faire rédactionnel)}

\begin{methode}[Comparer deux thèses philosophiques — Exemple : Sartre vs Camus]
\emph{Consigne type Probatoire / Baccalauréat technique :} « Comparez la conception de la liberté chez Sartre et chez Camus. » ou « L'homme est-il condamné à être libre ? Confrontation Sartre / Camus. »

\textbf{Étape 1 : Construire le tableau d'analyse (brouillon).}
\begin{center}
\begin{tabularx}{\linewidth}{|l|X|X|}\hline
\textbf{Critères de comparaison} & \textbf{Sartre (\emph{L'Être et le Néant}, \emph{L'Existentialisme...})} & \textbf{Camus (\emph{Le Mythe de Sisyphe}, \emph{L'Homme révolté})} \\\hline
\textbf{Origine du sens} & Inventé par le projet libre (existence précède essence) & Absent (divorce homme/monde), mais tension féconde \\\hline
\textbf{Nature de la liberté} & \textbf{Ontologique} : structure de la conscience (néantisation). Totale, angoissante. & \textbf{Conquise} : liberté \emph{dans} l'absurde (révolte). Limitée par la condition humaine. \\\hline
\textbf{Rapport à la mort} & Fin radicale qui \emph{annule} le projet (inachevé). « La mort n'est pas vécue. » & Horizon qui \emph{illumine} le présent. « Vivre sans appel. » \\\hline
\textbf{Attitude face à l'absurde} & L'absurde n'est pas le concept central ; l'angoisse oui. & \textbf{Concept central}. Révolte, liberté, passion = réponses. \\\hline
\textbf{Éthique / Engagement} & Engagement subjectif (universalisation). Responsabilité totale. & Solidarité des hommes révoltés. Refus du meurtre (limite). \\\hline
\end{tabularx}
\end{center}

\textbf{Étape 2 : Rédiger la dissertation en 3 temps.}
\begin{enumerate}
    \item \textbf{Introduction :} Accroche (citation ou fait) $\rightarrow$ Définition des termes $\rightarrow$ Problématique commune $\rightarrow$ Annonce du plan dialectique (Thèse Sartre / Antithèse Camus / Synthèse dépassement).
    \item \textbf{Développement :} Un paragraphe = un argument = une ligne du tableau. \emph{Ne pas faire : « Sartre dit... Camus dit... »} Faire : « \textbf{Si} la liberté est ontologique (Sartre), \textbf{alors} elle est source d'angoisse ; \textbf{mais} pour Camus, elle est conquête lucide face à l'absurde. »
    \item \textbf{Conclusion :} Bilan nuancé $\rightarrow$ Ouverture (ex : Ricœur, éthique de la responsabilité).
\end{enumerate}
\end{methode}

\courssub{Situation concrète camerounaise : Face au diagnostic d'un cancer terminal}

\begin{exoresolu}[Analyse plurielle d'une situation limite]
\textbf{Énoncé :} \emph{M. Kamga, 48 ans, cadre bancaire à Douala, père de trois enfants, reçoit un diagnostic de cancer du pancréas métastatique (espérance de vie : 3--6 mois). Comment les différentes approches philosophiques éclairent-elles son vécu et ses choix ?}
\tcblower
\textbf{Solution détaillée :}

\begin{enumerate}
    \item \textbf{Approche religieuse / ancestrale (Foi catholique + tradition Bamiléké) :}
    \begin{quote}
    « C'est la \emph{volonté de Dieu} / \emph{les ancêtres m'appellent}. Je dois \emph{mettre ma maison en ordre} (testament, réconciliation familiale, rites de protection pour les enfants). La souffrance est \emph{partage de la Croix} / \emph{épreuve initiatique}. Je prie pour la guérison, mais j'accepte le \emph{retour au village des ancêtres}. Mes funérailles (\emph{tsi}) seront le moment où je \emph{deviendrai} ancêtre protecteur. Le sens : \emph{fidélité} et \emph{transmission}. »
    \end{quote}

    \item \textbf{Approche existentialiste (Sartre / Heidegger) :}
    \begin{quote}
    « Je suis \emph{condamné à être libre} face à cette mort. Je \emph{choisis} comment vivre ces mois : écrire mes mémoires pour mes enfants ? Créer une fondation ? Refuser l'acharnement thérapeutique pour garder ma lucidité ? \emph{L'anticipation de la mort} (Heidegger) m'arrache au « On » (les convenances sociales, le travail) pour une \emph{authenticité} finale. Je suis \emph{responsable} du sens que je donne à ma fin. »
    \end{quote}

    \item \textbf{Approche de l'Absurde (Camus) :}
    \begin{quote}
    « Cette maladie est \emph{injuste}, \emph{sans raison}. Le monde est \emph{sourd}. Je ne me suicide pas (physiquement ni par l'espoir illusoire d'un miracle). Je \emph{me révolte} : je vis \emph{intensément} chaque lever de soleil sur le Wouri, chaque rire de mes petits-enfants, chaque plat de \emph{ndolé} partagé. \emph{Quantité} d'expériences, \emph{passion} du présent. \emph{« Il faut imaginer Sisyphe heureux »} : je pousse mon rocher (soins, douleur) en \emph{maître} de mon destin. »
    \end{quote}

    \item \textbf{Approche humaniste laïque (Ricœur / Soins palliatifs) :}
    \begin{quote}
    « Mon équipe soignante (médecin, infirmière, psychologue, aumônier laïc) reconnaît ma \emph{vulnérabilité}. Le sens naît de la \emph{relation} : dire « merci » à ma femme, « pardon » à mon frère, « je t'aime » à mes enfants. \emph{Récit de vie} : on enregistre ma voix pour les générations futures. \emph{Responsabilité} : je signe mes directives anticipées pour ne pas être une charge. La mort est l'\emph{achèvement} de mon histoire, pas son effacement. »
    \end{quote}
\end{enumerate}

\textbf{Conclusion pédagogique :} \emph{Aucune approche n'exclut nécessairement les autres dans la vie réelle.} M. Kamga peut prier (religieux), choisir ses soins (existentialiste), savourer un coucher de soleil (absurde) et dicter ses dernières volontés (humaniste). La philosophie \emph{outille} pour \emph{nommer} et \emph{assumer} cette complexité.
\end{exoresolu}

\begin{popfigure}
\begin{tikzpicture}[
    mindmap,
    concept color=popblue,
    level 1 concept/.append style={level distance=4.5cm, sibling angle=60, font=\bfseries\small, concept color=popblue!70},
    level 2 concept/.append style={level distance=3cm, sibling angle=45, font=\small, concept color=popcream},
    every node/.style={align=center, text=popdark},
    grow cyclic
]
\node[concept, concept color=popblue, font=\bfseries\normalsize, align=center]{Sens de\\ l'existence}
    child[concept color=poppurple] { node[concept, align=center]{Donné\\(Religion /\\(Ancêtres)}
        child { node[concept color=poppurpleL] {Créateur / Nzambé} }
        child { node[concept color=poppurpleL] {Vocation / Destinée} }
        child { node[concept color=poppurpleL] {Mort = Passage} }
    }
    child[concept color=popink] { node[concept, align=center]{Construit\\(Existentialisme)}
        child { node[concept color=poppinkL] {Existence $>$ Essence} }
        child { node[concept color=poppinkL] {Projet / Engagement} }
        child { node[concept color=poppinkL] {Mort = Échéance} }
    }
    child[concept color=popgold] { node[concept, align=center]{Absent\\(Absurde / Camus)}
        child { node[concept color=popgoldL] {Divorce Homme/Monde} }
        child { node[concept color=popgoldL] {Révolte / Liberté} }
        child { node[concept color=popgoldL] {Passion / Quantité} }
    }
    child[concept color=poporange] { node[concept, align=center]{Niable\\(Nihilisme)}
        child { node[concept color=poporangeL] {Néant des valeurs} }
        child { node[concept color=poporangeL] {Tout est permis} }
        child { node[concept color=poporangeL] {Désespoir / Ironie} }
    }
    child[concept color=popgreen] { node[concept, align=center]{Biologique\\(Matérialisme)}
        child { node[concept color=popgreenL] {Atomes / Neurones} }
        child { node[concept color=popgreenL] {Plaisir / Ataraxie} }
        child { node[concept color=popgreenL] {Mort = Néant} }
    }
    child[concept color=popgreen!60!black] { node[concept, align=center]{Relationnel\\(Humanisme laïque)}
        child { node[concept color=popgreenL!60!black] {Visage de l'Autre} }
        child { node[concept color=popgreenL!60!black] {Récit / Promesse} }
        child { node[concept color=popgreenL!60!black] {Responsabilité} }
    };
\end{tikzpicture}
\legende{Carte mentale des six grandes familles de réponses au sens de l'existence. Chaque branche développe trois concepts clés.}
\end{popfigure}

\begin{exercice}[Entraînement Bac — Dissertation]
\textbf{Sujet :} « La mort donne-t-elle un sens à la vie ? »
\begin{enumerate}
    \item Analysez le sujet : termes, présupposés, problème.
    \item Construisez un plan dialectique en trois parties (Thèse / Antithèse / Synthèse) en mobilisant \emph{au moins trois} des approches étudiées.
    \item Rédigez l'introduction complète (accroche, définitions, problématique, annonce de plan).
\end{enumerate}
\end{exercice}

\begin{corrige}[Exercice n°1 — Éléments de correction]
\textbf{Analyse :} « La mort » (fin de la vie biologique / horizon existentiel), « donne » (origine externe ? structure interne ?), « sens » (finalité, valeur, intelligibilité), « à la vie » (existence humaine). \textbf{Problématique :} La mort est-elle une \emph{fin extérieure} qui vide la vie de sens (néant), ou bien l'\emph{horizon interne} qui seul permet de hiérarchiser nos choix et de fonder des valeurs ?
\textbf{Plan :}
\begin{itemize}
    \item \textbf{I. Thèse : La mort comme néant qui abolit le sens (Matérialisme / Épicure ; Absurde / Camus).} La mort est l'arrêt de la conscience ; l'univers est indifférent. Le sens ne peut venir de ce qui détruit le sujet.
    \item \textbf{II. Antithèse : La mort comme horizon qui fonde le sens (Religieux / Heidegger / Sartre / Ricœur).} Pour le croyant, la mort est passage/jugement $\rightarrow$ urgence éthique. Pour Heidegger, l'anticipation de la mort individualise et authentifie. Pour Sartre, la finitude donne urgence au projet. Pour Ricœur, la finitude appelle le témoignage et la transmission.
    \item \textbf{III. Synthèse : La mort ne « donne » pas un sens tout fait ; elle \emph{exige} que nous \emph{fassens sens} de notre vie finie.} La mort est la \emph{condition de possibilité} du sens (sans fin, pas de choix, pas de rareté, pas de valeur), mais le \emph{contenu} du sens reste notre responsabilité (foi, engagement, révolte, éthique).
\end{itemize}
\textbf{Introduction type :} « "La mort n'a pas d'importance, tant qu'on est vivant elle n'est pas là, et quand elle est là on n'est plus" (Épicure). Pourtant, Heidegger affirme que "l'anticipation de la mort est la possibilité de l'impossibilité absolue" qui seul permet l'authenticité. Entre le néant épicurien et l'horizon heuristique, la mort est-elle la négation du sens ou sa condition ? Nous verrons que si la mort ne \emph{contient} pas le sens, elle en est l'\emph{exigence} irréductible. »
\end{corrige}

\coursec{La mort : Fin radicale ou horizon de sens ?}

Nous avons vu, dans la section précédente, que la question du sens de l'existence humaine a reçu des réponses diverses : le bonheur, la liberté, le devoir, la foi. Mais toutes ces réponses se heurtent à un obstacle commun et redoutable : la \cle{mort}. Si tout s'achève dans le néant, à quoi bon vivre ? Faut-il craindre la mort, l'accepter, ou l'intégrer à notre projet de vie ? La mort est-elle une \emph{fin radicale} qui anéantit tout sens, ou au contraire un \emph{horizon de sens} qui donne son prix à l'existence ? Telle est la question que nous allons maintenant examiner, en mobilisant quatre grands philosophes — Épicure, Heidegger, Sartre et Jankélévitch — puis en abordant les débats bioéthiques contemporains et les représentations sociales de la mort au Cameroun.

\courssub{Épicure : « La mort n'est rien pour nous »}

Commençons par une intuition simple. Quand tu étais enfant, la nuit, tu as peut-être eu peur d'un bruit dans la cour. Mais une fois la lampe allumée, tu as découvert que ce n'était qu'une chèvre attachée. La peur venait de l'\emph{imagination}, non de la réalité. Pour Épicure (341--270 av. J.-C.), la crainte de la mort est exactement de ce type : une peur irrationnelle, fondée sur une erreur de raisonnement, que la philosophie doit soigner comme le médecin soigne une maladie.

\begin{propriete}[Thèse d'Épicure : la mort n'est rien pour nous]
Dans la \emph{Lettre à Ménécée}, Épicure affirme : « La mort n'est rien pour nous, car ce qui est dissous ne ressent rien, et ce qui ne ressent rien n'est rien pour nous. »
\textbf{Démonstration (à connaître pour la dissertation) :}
\begin{enumerate}
\item Le bien et le mal consistent dans la \cle{sensation} (le plaisir est bien, la douleur est mal).
\item Or la mort est la \cle{privation de toute sensation} (le corps dissous ne sent plus rien).
\item Donc la mort n'est ni un bien ni un mal pour nous.
\end{enumerate}
\textbf{Argument symétrique :} le temps \emph{avant} notre naissance ne nous cause aucune angoisse ; or le temps \emph{après} notre mort lui est parfaitement symétrique (non-existence dans les deux cas). Il est donc incohérent de redouter l'un et pas l'autre.
\textbf{Critique classique :} Épicure néglige l'\emph{angoisse de l'anticipation} — ce qui nous effraie n'est pas l'état d'être mort, mais le processus de mourir et la perte de tout ce que nous aimons.
\end{propriete}

La visée d'Épicure est \cle{thérapeutique} : la philosophie est une médecine de l'âme. En supprimant la crainte de la mort, elle permet d'atteindre l'\cle{ataraxie}, c'est-à-dire la tranquillité de l'âme, l'absence de trouble. Celui qui ne craint plus la mort jouit pleinement du présent, comme le jardinier de Bandjoun qui savoure sa récolte sans penser à la saison sèche.

\courssub{Heidegger : la mort comme possibilité qui individualise}

Changeons radicalement de perspective. Pour Martin Heidegger (1889--1976), dans \emph{Être et Temps} (1927), la mort n'est pas un événement extérieur qui surviendra un jour : elle est la structure la plus intime de notre existence. L'être humain, que Heidegger nomme le \cle{Dasein} (« être-là »), est un « \cle{être-pour-la-mort} ».

Heidegger définit la mort par quatre caractères : elle est la possibilité \emph{la plus propre} (personne ne peut mourir à ma place), \emph{la plus non-relative} (elle me coupe de toute relation), \emph{la plus certaine} (je mourrai, c'est la seule certitude absolue), et \emph{indépassable} (on ne peut pas aller au-delà d'elle).

\begin{popfigure}
\begin{center}
\begin{tikzpicture}[node distance=0.5cm, every node/.style={font=\small}]
\node[draw=popblue, fill=popblueL, rounded corners, text width=3.4cm, align=center] (n1) {\textbf{1. Imminence}\\ La mort est certaine, à tout moment};
\node[draw=poporange, fill=poporangeL, rounded corners, text width=3.4cm, align=center, right=of n1] (n2) {\textbf{2. Inaliénabilité}\\ Personne ne meurt à ma place};
\node[draw=poppurple, fill=poppurpleL, rounded corners, text width=3.4cm, align=center, right=of n2] (n3) {\textbf{3. Indépassabilité}\\ Fin de toute possibilité};
\node[draw=popgreen, fill=popgreenL, rounded corners, text width=4.5cm, align=center, below=1.1cm of n2] (n4) {\textbf{Résolution -- Authenticité}\\ Assumer sa finitude pour vivre librement};
\draw[-{Stealth}, thick, popdark] (n1) -- (n2);
\draw[-{Stealth}, thick, popdark] (n2) -- (n3);
\draw[-{Stealth}, thick, popgreen] (n1.south) |- ($(n4.west)+(0,0.15)$);
\draw[-{Stealth}, thick, popgreen] (n2.south) -- (n4.north);
\draw[-{Stealth}, thick, popgreen] (n3.south) |- ($(n4.east)+(0,0.15)$);
\end{tikzpicture}
\end{center}
\legende{Les trois temps de la mort chez Heidegger. En assumant ces trois caractères dans l'angoisse, le Dasein accède à une existence authentique.}
\end{popfigure}

L'idée centrale est la suivante : tant que je vis « comme tout le monde » (le « \emph{on} » : on se marie, on travaille, on meurt...), je vis de manière \emph{inauthentique}. Mais l'\cle{angoisse} (\emph{Angst}) devant ma mort — qui est \emph{ma} mort, que nul ne peut prendre à ma place — me \cle{individualise} : elle me révèle que mon existence m'appartient et que je dois en décider moi-même. La mort n'est donc pas une fin absurde : elle est l'\emph{horizon} qui donne son urgence et son sérieux à chacun de mes choix. Comme le dit le proverbe : « c'est au pied du mur qu'on voit le maçon » — c'est devant la mort que l'homme se révèle à lui-même.

\courssub{Sartre : la mort comme absurdité}

Jean-Paul Sartre (1905--1980) retourne la thèse de Heidegger. Dans \emph{L'Être et le Néant} (1943), il soutient que la mort ne peut donner \emph{aucun} sens à la vie, pour trois raisons :

\begin{itemize}
\item La mort est \emph{absurde} : elle survient de l'extérieur, par hasard (un accident de la route Yaoundé--Douala, une crise cardiaque), sans aucune justification. Elle est « la fin injuste ».
\item La mort est « le triomphe du point de vue d'autrui » : une fois mort, je deviens un \emph{objet} entre les mains des autres — ce sont eux qui racontent ma vie, jugent mes actes, organisent mes funérailles. Ma liberté est définitivement figée.
\item La mort n'attend pas que mes projets soient accomplis : elle interrompt toujours « en pleine course ». Elle n'achève pas le sens, elle l'\emph{annihile}.
\end{itemize}

\begin{aretenir}[Opposition fondamentale]
Pour \textbf{Heidegger}, la mort est \emph{à l'intérieur} de la vie : elle est ma possibilité la plus propre et donne sens à mon existence. Pour \textbf{Sartre}, la mort est \emph{à l'extérieur} de la vie : c'est un accident absurde qui prive la vie de sens. Conséquence sartrienne : le sens doit être construit \emph{avant} la mort, dans l'action et la liberté, sans attendre aucune finalité.
\end{aretenir}

\courssub{Jankélévitch : l'innommable qui donne du prix à la vie}

Vladimir Jankélévitch (1903--1985), dans \emph{La Mort} (1966), propose une position intermédiaire et nuancée. La mort est pour lui « l'\cle{innommable} » : on ne peut ni la penser (penser la mort, c'est encore la penser du point de vue du vivant), ni la nommer adéquatement. Elle est l'\cle{irréversible} par excellence : ce qui est fait est fait, ce qui est perdu est perdu à jamais.

Mais précisément, cette irréversibilité a une valeur positive : \emph{« la mort donne du prix à la vie »}. C'est parce que nos instants sont rares, comptés, irrécupérables, qu'ils ont de la valeur. Une vie immortelle serait une vie sans urgence, sans choix véritables, donc sans valeur — comme une monnaie qui existerait en quantité infinie ne vaudrait plus rien. La rareté crée le prix ; la finitude crée la valeur de l'existence.

\courssub{La mort au défi de la médecine : la bioéthique contemporaine}

La médecine moderne a bouleversé la définition même de la mort. Autrefois, on constatait la mort par l'arrêt du cœur et de la respiration (\cle{mort cardiaque}). Mais les machines de réanimation peuvent maintenir artificiellement la circulation sanguine d'un patient dont le cerveau est définitivement détruit. On parle alors de \cle{mort encéphalique} (ou mort cérébrale) : destruction irréversible de l'ensemble du cerveau, constatée par des examens cliniques et par l'absence d'activité électrique cérébrale.

Cette distinction pose un problème éthique redoutable : le patient en mort encéphalique, maintenu « cœur battant », est-il mort ou vivant ? C'est pourtant dans cet état que le prélèvement d'organes (rein, foie, cœur) est possible pour la transplantation. Le \emph{donneur d'organe à cœur battant} est juridiquement mort (cerveau détruit) mais biologiquement encore « chaud ». Où passe la frontière entre le respect du corps et le sauvetage d'autres vies ?

\begin{definition}[Acharnement thérapeutique et soins palliatifs]
L'\cle{acharnement thérapeutique} (ou obstination déraisonnable) est le maintien de traitements \emph{inutiles}, \emph{disproportionnés} ou ne servant qu'à prolonger artificiellement la vie d'un patient condamné. Il faut le distinguer des \cle{soins palliatifs}, qui ne visent pas à prolonger ni à abréger la vie, mais à \emph{soulager la douleur} et à \emph{accompagner} le mourant avec dignité. Refuser l'acharnement n'est pas donner la mort : c'est accepter la mort.
\end{definition}

On distingue ensuite deux attitudes face au mourant :
\begin{itemize}
\item \emph{Laisser mourir} : arrêter un traitement devenu inutile ou disproportionné (par exemple débrancher une machine qui ne fait que prolonger l'agonie). C'est une \emph{omission} qui laisse la maladie suivre son cours.
\item \emph{Faire mourir} : provoquer directement la mort, par exemple par injection létale. C'est l'\cle{euthanasie active}, une \emph{action}.
\end{itemize}

\begin{attention}[Erreur juridique fréquente au Cameroun]
Au Cameroun, l'\textbf{euthanasie active} (injection létale, même à la demande du patient) est juridiquement un \textbf{homicide volontaire}, réprimé par le Code pénal (articles 275 à 277, qui punissent le meurtre et l'assistance au suicide). En revanche, l'\textbf{arrêt de traitement} (« laisser mourir ») est toléré en éthique médicale lorsque le traitement est inutile ou disproportionné, mais le cadre juridique reste flou : aucun texte camerounais n'encadre clairement la fin de vie. Ne jamais écrire dans une copie que « l'euthanasie est autorisée si le patient consent » : c'est faux en droit camerounais.
\end{attention}

\courssub{La mort comme événement social : les représentations camerounaises}

Dans les sociétés camerounaises, la mort n'est jamais un simple fait biologique : c'est un événement cosmique et social qui doit être \emph{interprété} et \emph{ritualisé}. On distingue classiquement :

\begin{itemize}
\item La \cle{belle mort} : mourir très âgé, entouré de ses enfants, après avoir donné sa bénédiction. Elle assure une transition paisible vers le statut d'\emph{ancêtre}, protecteur de la lignée. Les funérailles sont alors une célébration autant qu'un deuil.
\item La \cle{mauvaise mort} : mort soudaine (accident, « sudden death »), suicide, mort en couches, mort par épidémie. Elle est perçue comme une rupture de l'ordre cosmique, souvent attribuée à une faute, une sorcellerie ou une malédiction. Elle exige des \emph{rites de purification} et d'apaisement pour empêcher le défunt de devenir un esprit errant et dangereux.
\end{itemize}

\begin{exemplebox}[La mortalité maternelle : un drame social chiffré]
Au Cameroun, le taux de mortalité maternelle est d'environ \textbf{406 décès pour \num{100000} naissances vivantes} (enquête EDS-MICS 2018). La mort d'une femme en couches reste un drame social majeur : socialement, elle est souvent perçue comme une « mauvaise mort » nécessitant des rites spécifiques ; médicalement, elle traduit les insuffisances de l'accès aux soins (distance des centres de santé, coût des accouchements, pénurie de personnel). Ici, la philosophie rejoint la santé publique : derrière chaque « représentation » de la mort se cachent des causes concrètes qu'il est possible de combattre.
\end{exemplebox}

\begin{popfigure}
\begin{center}
\begin{tikzpicture}
\begin{axis}[
  ybar, bar width=0.7cm,
  width=0.95\linewidth, height=6.2cm,
  ymin=0, ymax=60,
  ylabel={Part estimée des décès (en \%)},
  symbolic x coords={Maladies transmissibles, Maladies non transmissibles, Blessures},
  xtick=data,
  x tick label style={font=\small, align=center},
  nodes near coords, every node near coord/.append style={font=\small},
  axis line style={popdark},
]
\addplot[fill=popblue, draw=popdark] coordinates {(Maladies transmissibles,48) (Maladies non transmissibles,37) (Blessures,15)};
\end{axis}
\end{tikzpicture}
\end{center}
\legende{Causes de décès au Cameroun (ordres de grandeur OMS/INS, valeurs arrondies). \textbf{Maladies transmissibles} : paludisme, VIH, tuberculose (environ 48 \%). \textbf{Maladies non transmissibles} : maladies cardio-vasculaires, cancers (environ 37 \%). \textbf{Blessures} : accidents de la route, violences (environ 15 \%). Impact sur la représentation : les blessures et morts soudaines nourrissent l'image de la « mauvaise mort », tandis que les maladies longues laissent le temps de « préparer » la fin.}
\end{popfigure}

\begin{savaistu}[La veillée funèbre, un « procès social » du défunt]
La \emph{veillée} (wake-keeping) pratiquée chez les Bamiléké, les Bassa, les Douala, les Beti et bien d'autres peuples du Cameroun n'est pas un simple hommage au disparu. Elle remplit trois fonctions : c'est un \textbf{procès social} du défunt, où l'on fait publiquement la « comptabilité » de sa vie (ses biens, ses dettes, ses enfants, ses bonnes et mauvaises actions) ; c'est une \textbf{protection} rituelle de la famille contre les mauvais sorts ; et c'est enfin la \textbf{cérémonie d'intronisation} qui transforme le défunt en ancêtre, c'est-à-dire en membre actif et protecteur de la communauté. Mourir, au Cameroun, ce n'est donc pas disparaître : c'est changer de statut.
\end{savaistu}

\courssub{Méthode : analyser un texte philosophique sur la mort}

\begin{methode}[Expliquer un texte sur la mort (ex. : Heidegger, \emph{Être et Temps} § 50 ; Camus, \emph{Le Mythe de Sisyphe})]
\begin{enumerate}
\item \textbf{Identifier la thèse} : que soutient l'auteur ? (Ex. : « la mort est ma possibilité la plus propre » ; « la vie est absurde ».)
\item \textbf{Repérer les concepts techniques} : souligner les mots-clés (\emph{Dasein}, \emph{angoisse}, \emph{absurde}, \emph{ataraxie}) et les définir à partir du texte.
\item \textbf{Reconstituer l'argumentation} : quelles prémisses ? quelle conclusion ? Quel exemple l'auteur utilise-t-il ?
\item \textbf{Discuter la portée} : la thèse est-elle convaincante ? Peut-on lui opposer un contre-exemple (la mort d'un enfant, la mort au service d'une cause) ? Confronter avec un autre auteur (Heidegger contre Sartre, Épicure contre l'expérience de l'angoisse).
\end{enumerate}
\end{methode}

\begin{exoresolu}[Dissertation guidée : « Faut-il craindre la mort ? »]
\textbf{Énoncé.} Dégager l'intérêt philosophique du sujet suivant en le discutant : \emph{Faut-il craindre la mort ?}
\tcblower
\textbf{Corrigé rédigé (plan détaillé).}
\emph{Introduction.} La mort est la seule certitude absolue de l'existence humaine, et pourtant elle suscite une peur quasi universelle. Mais cette crainte est-elle rationnelle ? Si la philosophie est un apprentissage de la sagesse, elle doit nous apprendre à bien mourir — ou à ne plus craindre de mourir. Nous demanderons d'abord si la crainte de la mort ne repose pas sur une illusion (Épicure), puis si la mort ne doit pas plutôt être assumée comme l'horizon de notre liberté (Heidegger), avant d'examiner si cette crainte ne cache pas en réalité l'amour de la vie (Sartre, Jankélévitch).

\emph{I. La crainte de la mort est irrationnelle.} Selon Épicure, le mal suppose la sensation ; or la mort est privation de sensation ; donc la mort n'est pas un mal. L'argument symétrique (avant la naissance / après la mort) montre l'incohérence de notre peur. La philosophie a une fonction thérapeutique : délivrer de cette crainte pour atteindre l'ataraxie.

\emph{II. La mort ne doit pas être crainte mais assumée.} Chez Heidegger, l'angoisse devant la mort n'est pas une peur à éliminer mais une révélation : elle m'individualise et me renvoie à ma liberté. La mort, possibilité la plus propre, certaine et indépassable, donne son sérieux à l'existence. La crainte devient résolution.

\emph{III. Discussion : ce que nous craignons, c'est de perdre la vie.} Avec Sartre, on peut objecter que la mort est absurde et prive la vie de sens ; avec Jankélévitch, que c'est précisément parce qu'elle est irréversible qu'elle donne du prix à la vie. Craindre la mort, c'est donc aimer la vie : la crainte est légitime comme attachement, illégitime comme paralysie.

\emph{Conclusion.} Il ne faut pas craindre la mort comme un mal qui nous attend, mais l'assumer comme la limite qui donne sa valeur à chaque instant. La sagesse consiste à transformer la peur de mourir en volonté de vivre pleinement.
\end{exoresolu}

\begin{exercice}[Pour s'entraîner]
1. Définis les termes suivants : ataraxie, mort encéphalique, acharnement thérapeutique, euthanasie active, « mauvaise mort ».
2. Explique l'argument symétrique d'Épicure et formule une critique.
3. Confronte Heidegger et Sartre sur la question : la mort donne-t-elle un sens à la vie ? (tableau ou texte de 15 lignes).
4. Sujet de dissertation : \emph{La mort est-elle un événement de la vie ?}
\end{exercice}

\begin{corrige}[Exercice]
\textbf{1.} Ataraxie : tranquillité de l'âme, absence de trouble, but de la sagesse épicurienne. Mort encéphalique : destruction irréversible de l'ensemble du cerveau, critère médical moderne de la mort, distinct de la mort cardiaque. Acharnement thérapeutique : maintien de traitements inutiles ou disproportionnés qui prolongent artificiellement la vie. Euthanasie active : provoquer directement la mort d'un patient (interdite au Cameroun, assimilée à un homicide). « Mauvaise mort » : mort soudaine ou violente (accident, suicide, mort en couches) perçue comme perturbant l'ordre cosmique et exigeant des rites de purification.

\textbf{2.} Épicure remarque que le temps qui précède notre naissance est une non-existence qui ne nous cause aucune angoisse ; or le temps qui suit notre mort est une non-existence identique ; il est donc irrationnel de redouter l'une et pas l'autre. Critique : cette symétrie ignore que la mort prive d'un bien que l'on possède déjà (la vie, les êtres aimés), ce que la non-naissance ne faisait pas ; elle néglige aussi l'angoisse de l'anticipation.

\textbf{3.} Pour Heidegger, la mort est \emph{intérieure} à la vie : possibilité la plus propre, certaine, indépassable, elle individualise le Dasein et, par l'angoisse assumée, ouvre l'existence authentique ; elle est donc un horizon de sens. Pour Sartre, la mort est \emph{extérieure} et absurde : elle survient par hasard, interrompt les projets et livre le mort au point de vue d'autrui ; elle ne donne pas de sens, elle le supprime. Le sens doit donc être construit par la liberté, avant la mort. Les deux s'accordent sur un point : c'est la finitude qui oblige à choisir.

\textbf{4.} Pistes : I. La mort semble extérieure à la vie (elle la termine, elle est absurde — Sartre ; elle n'est « rien pour nous » — Épicure). II. Pourtant elle structure la vie de l'intérieur (être-pour-la-mort chez Heidegger ; elle donne du prix à la vie — Jankélévitch). III. Au-delà de l'individu, la mort est un événement social et symbolique (rites, veillées, ancestralité au Cameroun) : elle appartient à la vie de la communauté.
\end{corrige}

En définitive, la mort se situe à la croisée de deux lectures : fin radicale qui anéantit tout sens (Sartre), ou horizon qui confère à l'existence son urgence et son prix (Heidegger, Jankélévitch). La sagesse épicurienne nous invite à ne pas la craindre, la bioéthique nous oblige à la définir avec rigueur, et les traditions camerounaises nous rappellent qu'elle est aussi un passage, un événement social qui engage toute la communauté. Reste à voir, dans la section suivante, comment ces analyses s'incarnent concrètement dans les situations camerounaises de l'ancestralité.

\coursec{Existence, Mort et Ancestralité : Le cas camerounais (Situations concrètes)}

La philosophie occidentale, d'\emph{Épicure} à \emph{Heidegger} en passant par \emph{Sartre}, a longtemps interrogé la mort soit comme \textbf{finitude radicale} (retour au néant), soit comme \textbf{horizon qui donne du sens à la vie} (l'\emph{être-pour-la-mort}, \emph{Sein-zum-Tode}, chez Heidegger). Or, au Cameroun comme dans une grande partie de l'Afrique subsaharienne, l'expérience vécue de la mort ne s'arrête pas à l'arrêt des fonctions biologiques. La mort est pensée comme un \textbf{passage}, un changement d'état qui inscrit le défunt dans une \textbf{continuité lignagère}. Comprendre cette anthropologie de la mort n'est pas un exercice folklorique : c'est une exigence éthique, juridique et pratique pour tout technicien, soignant, administrateur ou ingénieur appelé à gérer des situations concrètes où se rencontrent, et parfois s'affrontent, \textbf{logique biomédicale} et \textbf{logique coutumière}.

\begin{attention}[Précaution méthodologique]
Les termes vernaculaires cités dans cette leçon (\emph{Nkwa}, \emph{Mboa}, \emph{Zam-zam}, etc.) varient fortement selon les langues et les groupes. Ils sont donnés ici à titre \textbf{illustratif}, à partir d'exemples documentés (pays bassa, Grassfields bamiléké, aire beti). L'essentiel pour l'examen est de retenir les \textbf{structures} (passage, liminalité, agrégation, dualité des normes), non de mémoriser un vocabulaire qui prétendrait à l'uniformité.
\end{attention}

\courssub{Continuité ontologique : la mort comme changement d'état et institution de l'ancêtre}

\begin{definition}[L'ancêtre : un défunt qualifié]
Dans de nombreuses conceptions camerounaises (bassa, bamiléké, beti, etc.), l'\textbf{ancêtre} n'est pas simplement un « mort ». C'est un \textbf{défunt qualifié} qui remplit classiquement trois conditions cumulatives :
\begin{enumerate}
    \item \textbf{Une « belle mort »} : mort à un âge avancé, entouré des siens, sans suicide ni mort violente non ritualisée.
    \item \textbf{La descendance} : avoir laissé des enfants, garants de la mémoire et du culte.
    \item \textbf{L'agrégation rituelle} : avoir reçu les rites funéraires puis les rites de fin de deuil, qui le séparent du monde des vivants pour l'insérer dans le \textbf{monde invisible}.
\end{enumerate}
L'ancêtre devient alors \textbf{intercesseur} auprès de l'Être suprême : il protège le lignage, veille sur la fécondité (femmes, champs, bétail) et sanctionne les transgressions. Il demeure, symboliquement, un \textbf{acteur social} à part entière.
\end{definition}

\begin{attention}[Écueil ethnophilosophique : l'essentialisme]
Il est faux de dire « les Africains pensent que... ». Le Cameroun compte \textbf{plus de 250 groupes ethniques} aux conceptions funéraires distinctes. De plus, la \textbf{modernité} (christianisme, islam, laïcité, urbanisation) recompose les rites : un cadre urbain de Douala peut être chrétien le dimanche, consulter un tradipraticien pour un problème d'héritage, et exiger d'être inhumé « au village » pour « rejoindre les siens ». Ne figez pas les cultures : analysez des \textbf{situations singulières}. Évitez aussi les termes péjoratifs hérités de la colonisation (par exemple « Kirdi », terme méprisant désignant les peuples du Nord, à remplacer par les noms propres des groupes : Mafa, Mofu, Toupouri, etc.).
\end{attention}

\begin{exemplebox}[Le caveau familial : lieu de mémoire et d'ancrage]
Chez les Bassa comme chez de nombreux groupes de l'Ouest, du Centre et du Littoral, la sépulture familiale (souvent appelée \emph{Mboa}, « la maison », « le chez-soi ») se situe dans la cour de la concession ou sur le terrain familial. Ce n'est pas un cimetière anonyme.
\begin{itemize}
    \item \textbf{Fonction symbolique} : le corps doit reposer sur la \textbf{terre des ancêtres} pour que le défunt veille sur les siens. Le « retour au village » du corps d'un citadin décédé en ville est une pratique massive : la mort sociale et spirituelle s'accomplit pleinement dans ce \textbf{retour à la terre}.
    \item \textbf{Tension urbaine} : les cimetières municipaux (Yaoundé, Douala) sont saturés et perçus comme des lieux sans ancrage. Les règlements municipaux imposent des concessions à durée limitée, ce qui heurte la temporalité coutumière, qui pense la sépulture comme définitive.
\end{itemize}
\end{exemplebox}

\courssub{Les rites de passage : structure universelle (Van Gennep) et variations locales}

\emph{Arnold Van Gennep} (\emph{Les Rites de passage}, 1909) a montré que tout changement de statut (naissance, puberté, mariage, mort) suit une structure tripartite : \textbf{séparation}, \textbf{marge} (ou liminalité), \textbf{agrégation}. Appliquons ce schéma aux funérailles camerounaises.

\begin{popfigure}
\begin{tikzpicture}[scale=0.95, every node/.style={font=\small}]
    \tikzset{
        phase/.style={draw=popblue, thick, fill=popblueL, rounded corners, minimum height=1.2cm, text width=3.6cm, align=center},
        arrow/.style={-{Stealth}, thick, popdark},
        rite/.style={draw=poporange, fill=poporangeL, rounded corners, text width=3.2cm, align=center, font=\footnotesize},
        actor/.style={draw=popgreen, fill=popgreenL, rounded corners, text width=2.6cm, align=center, font=\footnotesize},
    }
    \node[phase, align=center] (sep) at (0,0){\textbf{SÉPARATION}\\ Mort biologique};
    \node[phase, align=center] (marge) at (6.2,0){\textbf{MARGE / LIMINALITÉ}\\ Deuil, interdits, transmission};
    \node[phase, align=center] (agreg) at (12.4,0){\textbf{AGRÉGATION}\\ Fin de deuil, statut d'ancêtre};

    \draw[arrow] (sep.east) -- node[above, font=\footnotesize] {Veillée, inhumation} (marge.west);
    \draw[arrow] (marge.east) -- node[above, font=\footnotesize] {Cérémonie publique, levée des interdits} (agreg.west);

    \node[rite, align=center] (r1) at (0,-2.6){Annonce codifiée (tam-tam, cris)\\ Toilette et habillage du corps\\ Veillée (chants, prières, ripaille)};
    \draw[arrow] (sep.south) -- (r1.north);

    \node[rite, align=center] (r2a) at (6.2,-2){Deuil strict : pagne sombre, rasage\\ Interdits alimentaires et sociaux};
    \node[rite, align=center] (r2b) at (6.2,-3.8){Désignation de l'héritier\\ Transmission des biens et des secrets};
    \draw[arrow] (marge.south) -- (r2a.north);
    \draw[arrow] (r2a.south) -- (r2b.north);

    \node[rite, align=center] (r3) at (12.4,-2.6){Grande cérémonie (fête)\\ Sacrifices et libations\\ Danses, remise de l'héritage};
    \draw[arrow] (agreg.south) -- (r3.north);

    \node[actor, draw=poppurple, fill=poppurpleL, align=center] (anc) at (0,-5.2){\textbf{ANCÊTRE}\\ (monde invisible)};
    \draw[arrow, poppurple, dashed] (agreg.south) .. controls (14,-5.2) .. node[below, font=\footnotesize, poppurple] {Protection du lignage} (anc.east);
    \draw[arrow, poppurple, dashed] (anc.west) -- (-2.6,-5.2) -- (-2.6,0.8) -- node[above, font=\footnotesize, poppurple] {Consultations, offrandes, libations} (sep.north west);

    \node[actor] (a1) at (-3.4,0) {Héritier / chef de famille};
    \node[actor] (a2) at (15.8,0) {Notables, prêtres, tradipraticiens};
    \draw[arrow, dashed] (a1.east) -- (sep.west);
    \draw[arrow, dashed] (agreg.east) -- (a2.west);
\end{tikzpicture}
\legende{Cycle vie--mort--ancestralité (schéma simplifié inspiré des aires bassa et bamiléké). La mort n'est pas une sortie du système social, mais une réorganisation de la place de l'individu. La phase de marge (liminalité) est cruciale : le deuil est une « mort sociale » temporaire du survivant, qui renaît purifié à la levée des interdits.}
\end{popfigure}

\begin{exemplebox}[Le deuil codifié : une technologie psychique et sociale]
Le deuil coutumier (par exemple le \emph{Nkwa} en pays bassa) n'est pas une simple tristesse : c'est une \textbf{institution codifiée}.
\begin{itemize}
    \item \textbf{Marqueurs corporels} : port du pagne sombre, rasage du crâne, parfois port de cordes ou de feuilles aux poignets.
    \item \textbf{Interdits (tabous)} : restrictions alimentaires (nourriture « blanche », sans sel ni piment, signe de pureté), interdiction des rapports sexuels, de la danse, des travaux lourds.
    \item \textbf{Durée} : variable (quelques mois à deux ans) selon le rang du défunt et la proximité du lien.
    \item \textbf{Fonction} : la \textbf{transition}. Le survivant est provisoirement « mort » socialement (liminalité). La levée du deuil est une \textbf{renaissance} : bain rituel, onction, habits neufs, fête. Le rite structure ainsi collectivement le \textbf{travail de deuil psychique}, au lieu de laisser l'individu seul face à sa perte.
\end{itemize}
\end{exemplebox}

\courssub{Conflits de normes : médecine moderne, droit de l'État et coutume}

La situation concrète la plus fréquente pour un technicien (infirmier, administrateur, officier d'état civil, juriste) est le \textbf{conflit de qualification du corps} au moment du décès : pour l'hôpital, le corps est un objet sanitaire et juridique ; pour la famille, il reste une \textbf{personne} insérée dans le lignage.

\begin{methode}[Étude de cas guidée : le corps de M. Essomba]
\textbf{Situation} : M. Essomba, 62 ans, chef de famille, décède à l'hôpital central de Yaoundé à 14 h des suites d'un accident vasculaire cérébral.
\begin{enumerate}
    \item \textbf{Identifier les acteurs et leurs normes} :
    \begin{itemize}
        \item \textbf{Hôpital / État} : certificat médical de décès (obligatoire), déclaration à l'état civil, dépôt à la morgue (hygiène), autopsie possible si la mort est suspecte, délais légaux avant inhumation.
        \item \textbf{Famille / coutume} : réticence devant la morgue (« on ne laisse pas un chef seul dans le froid et le noir »), exigence de la veillée à domicile, rites de toilette par des personnes qualifiées, refus de l'autopsie (atteinte à l'intégrité du corps).
    \end{itemize}
    \item \textbf{Analyser les enjeux} :
    \begin{itemize}
        \item \textbf{Santé publique} : conservation du corps, risque infectieux éventuel.
        \item \textbf{Droit coutumier} : le corps relève du \textbf{lignage} ; le chef de famille (héritier) en a la garde.
        \item \textbf{Dignité} : pour la famille, la dignité = respect des rites ; pour l'hôpital, dignité = légalité, hygiène, rapidité.
    \end{itemize}
    \item \textbf{Proposer une médiation technique} (rôle du technicien) :
    \begin{itemize}
        \item Établir rapidement le certificat de décès lorsque la cause est claire, sans autopsie systématique.
        \item Obtenir l'autorisation administrative de remise du corps pour une veillée à domicile encadrée.
        \item Proposer un dispositif de conservation (glace, climatisation de la pièce mortuaire) compatible avec la veillée.
        \item S'appuyer sur un \textbf{médiateur} (notable, chef de quartier, aumônier) pour faciliter le dialogue.
    \end{itemize}
\end{enumerate}
\end{methode}

\begin{exoresolu}[Conflit autopsie contre rites : le cas d'un notable]
\textbf{Énoncé} : Un notable décède brutalement. Le procureur ordonne une autopsie (soupçon d'empoisonnement). La famille s'y oppose : « on ne coupe pas un chef ». L'hôpital bloque la remise du corps. Analysez le conflit et proposez une solution.
\tcblower
\textbf{Solution détaillée} :
\begin{enumerate}
    \item \textbf{Cadre légal} : en cas de mort suspecte, le droit pénal camerounais permet au procureur d'ordonner une autopsie ; la coutume, elle, sacralise l'intégrité du corps. Deux normes légitimes s'affrontent.
    \item \textbf{Argumentaire de médiation} : expliquer à la famille que l'autopsie peut \textbf{blanchir l'entourage} (l'héritier est souvent soupçonné en cas de mort brutale d'un chef) ; le refus peut passer, aux yeux de tous, pour un aveu.
    \item \textbf{Compromis technique} : autopsie réalisée par un médecin légiste, si possible en présence d'un représentant de la famille comme témoin, avec information préalable sur les gestes pratiqués.
    \item \textbf{Réparation symbolique} : après l'autopsie, rites de « réparation » du corps (toilette, onctions, prières) avant la veillée, afin de rétablir l'intégrité symbolique du défunt.
\end{enumerate}
\textbf{Leçon philosophique} : la médiation ne consiste pas à opposer « science » et « superstition », mais à articuler deux exigences de \textbf{dignité} : la vérité juridique et le respect du passage.
\end{exoresolu}

\courssub{Gestion des « mauvaises morts » : épidémies (choléra, sida, Covid-19)}

Les épidémies révèlent la tension extrême entre la \textbf{gestion étatique de la santé publique} (ce que Foucault appelle la \emph{biopolitique}) et l'\textbf{anthropologie du soin} (accompagnement rituel du passage).

\begin{popfigure}
\begin{tikzpicture}
\begin{axis}[
    ybar, bar width=14pt,
    width=\linewidth, height=7.5cm,
    symbolic x coords={Choléra, Sida, Covid-19},
    xtick=data,
    ylabel={Intensité (indice qualitatif de 1 à 10)},
    ymin=0, ymax=11,
    legend style={at={(0.5,-0.18)}, anchor=north, legend columns=-1, font=\footnotesize},
    nodes near coords, nodes near coords align={vertical},
    every node near coord/.append style={font=\footnotesize},
]
\addplot[fill=popblue!60] coordinates {(Choléra,8) (Sida,9) (Covid-19,10)};
\addplot[fill=poporange!70] coordinates {(Choléra,6) (Sida,7) (Covid-19,8)};
\legend{Conflit (enterrement sécurisé contre rites), Adaptation (négociation, rites modifiés)}
\end{axis}
\end{tikzpicture}
\legende{Gestion des « mauvaises morts » épidémiques au Cameroun : lecture qualitative. Le conflit entre protocoles sanitaires et rites reste élevé, mais les capacités d'adaptation (négociation, rites différés, outils numériques) progressent. Indices construits à titre pédagogique, sans valeur statistique officielle.}
\end{popfigure}

\textbf{Lecture du document, épidémie par épidémie :}
\begin{itemize}
    \item \textbf{Choléra} (épidémies récurrentes, notamment dans le Nord et le Littoral) : enterrements « sécurisés » par des équipes sanitaires (sacs hermétiques, chaux), sans toilette ni veillée traditionnelles ; d'où de vives résistances familiales, parfois des affrontements.
    \item \textbf{Sida} (pic de l'épidémie dans les années 1990--2000) : mort longtemps vécue comme « honteuse » ; dissimulation de la cause du décès, enterrements rapides et discrets, stigmatisation des veuves et des orphelins. Le défunt risquait d'être exclu symboliquement du statut d'ancêtre.
    \item \textbf{Covid-19} (2020--2022) : protocole strict (corps non visible, inhumation rapide par des équipes protégées). \textbf{Innovations observées} : veillées « virtuelles » par messagerie et visioconférence, bénédiction du cercueil fermé par un prêtre ou un imam en présence d'un notable, rites de fin de deuil \textbf{différés} après la crise.
\end{itemize}

\begin{aretenir}[Principes d'action pour le technicien face aux « mauvaises morts »]
\begin{enumerate}
    \item \textbf{Ne pas disqualifier le deuil} : dire « ce n'est qu'un virus » nie la réalité vécue du passage. Valider la douleur : « je sais que vous ne pouvez pas accomplir les rites habituels, c'est une épreuve terrible ».
    \item \textbf{Proposer des substituts rituels} : enterrement sécurisé, mais avec présence d'un représentant de la famille (à distance protégée), prière autorisée, objet symbolique (pagne, croix, Coran) déposé avec le corps.
    \item \textbf{Planifier le rite de réparation différé} : de nombreuses traditions prévoient des rites de « rattrapage » lorsque la « belle mort » a été empêchée. Les institutions (hôpital, mairie, organisations) gagnent à \textbf{accompagner} ces cérémonies différées, six mois ou un an plus tard : c'est une forme de \textbf{réparation symbolique} et de justice envers les familles.
\end{enumerate}
\end{aretenir}

\courssub{Héritage, succession et « mort sociale » du chef}

La mort biologique déclenche l'ouverture de la succession. Mais dans les sociétés à chefferies (Grassfields bamiléké et bamoun, sultanats peuls du Nord, certaines chefferies beti et sawa), la \textbf{mort sociale} (fin de la fonction) peut précéder ou suivre de loin la mort biologique.

\begin{propriete}[Dualité normative : droit coutumier et droit moderne]
\begin{itemize}
    \item \textbf{Filiations patrilinéaires} (bassa, bamiléké, beti, peul, etc.) : l'héritier principal est classiquement le \textbf{fils aîné} (ou le frère du défunt). Il reçoit la charge du patrimoine lignager (terres, cases, titres) et la \textbf{charge du culte des ancêtres}. Les cadets et les filles reçoivent des donations, mais la coutume leur reconnaît rarement une part successorale égale.
    \item \textbf{Filiations matrilinéaires} (certains groupes du Littoral et du Sud-Ouest) : l'héritier est le \textbf{neveu utérin} (fils de la sœur du défunt) ; le patrimoine suit le sang maternel.
    \item \textbf{Droit moderne} (droit civil hérité du droit français et du droit anglais selon les régions, lois sur l'état civil, jurisprudence) : tend vers l'égalité entre enfants et reconnaît des droits au conjoint survivant ; le testament écrit est admis.
\end{itemize}
\textbf{Conflit type} : le fils aîné (logique coutumière) s'oppose aux cadets et à la veuve (logique du droit moderne et du testament). Même lorsque le tribunal tranche le partage des biens, la \textbf{charge du culte ancestral} (entretien de la sépulture familiale, sacrifices annuels) reste dévolue à l'héritier coutumier, ce qui crée des tensions si celui-ci n'a pas les moyens financiers de l'assumer.
\end{propriete}

\begin{attention}[Exactitude juridique]
Au Cameroun, il n'existe pas de code successoral unifié : coexistent le droit civil (d'origine française à l'Est, anglaise à l'Ouest), le droit coutumier et le droit islamique au Nord. Les tribunaux appliquent souvent le droit coutumier aux successions « traditionnelles » et le droit moderne aux successions déclarées à l'état civil. Ne citez pas de numéro de loi sans certitude : au Baccalauréat technique, on attend l'\textbf{analyse du conflit de normes}, pas un exposé de droit positif.
\end{attention}

\begin{savaistu}[La « mort sociale » du chef : retraite rituelle et intronisation du successeur]
Dans les royaumes des Grassfields (bamiléké, bamoun, Nso), le chef (\emph{Fo}, \emph{Fon}, \emph{Sultan}) est le garant de l'ordre et de la fécondité du royaume. S'il devient très âgé, gravement malade ou incapable d'exercer la parole sacrée, il peut devenir, symboliquement, un danger pour la communauté.
\begin{itemize}
    \item \textbf{Mise à l'écart rituelle} : le conseil des notables peut décider sa « mort sociale » : retrait des insignes, interdiction de paraître en public, isolement dans le palais.
    \item \textbf{Intronisation du successeur} : le fils désigné est installé, parfois du vivant de son père. Le vieux chef devient un « ancêtre vivant ».
    \item \textbf{Secret de la mort} : la mort biologique d'un grand chef est parfois annoncée longtemps après (secret d'État) ; les funérailles officielles scellent alors la transition définitive.
\end{itemize}
\textbf{Leçon pour le gestionnaire} : ne jamais confondre \textbf{statut biologique} et \textbf{statut institutionnel}. Un chef « vivant » peut être socialement « mort » (dessaisi de sa fonction), et un corps « mort » peut rester politiquement « vivant » (symbole d'unité du groupe).
\end{savaistu}

\courssub{Boîte à outils : grille d'analyse anthropologique d'un rite funéraire}

\begin{center}
\begin{tabularx}{\linewidth}{|l|X|}\hline
\rowcolor{popblueL}\textbf{Dimension} & \textbf{Indicateurs d'observation / questions guides} \\ \hline
\textbf{1. Acteurs} & Qui décide ? (héritier, notables, médecin, prêtre, imam, mairie, police). Qui exécute ? (toilette, terrassement, rites secrets). Qui paie ? \\ \hline
\textbf{2. Espace} & Lieu du décès (hôpital, domicile, route). Lieu de la veillée (concession, salon, morgue). Lieu d'inhumation (sépulture familiale au village, cimetière municipal). Trajets du corps. \\ \hline
\textbf{3. Temps} & Délai entre décès et inhumation. Durée du deuil. Horaires rituels (libations à l'aube). Rites différés (fin de deuil, cérémonies de mémoire). \\ \hline
\textbf{4. Objets symboliques} & Corps (entier, os, cheveux). Tissus (pagne de deuil, tenues royales). Végétaux et animaux sacrifiés (poulet, chèvre, bœuf). Insignes (bâtons de commandement, colliers). \\ \hline
\textbf{5. Paroles} & Annonce (tam-tam, radio, réseaux sociaux). Éloges funèbres, prières, testaments oraux, chants de deuil puis chants de joie (fin de deuil). \\ \hline
\textbf{6. Fonction sociale} & Redistribution des biens et des titres. Réconciliation du lignage. Affirmation du statut social. Renforcement des alliances. \\ \hline
\textbf{7. Fonction psychique} & Contenir l'angoisse de la mort. Structurer le travail de deuil par phases. Gérer la culpabilité des survivants. Intégrer la veuve et les orphelins. \\ \hline
\end{tabularx}
\end{center}

\begin{exercice}[Application : analyse d'un terrain]
Assistez (ou visionnez un reportage fiable) à une veillée funèbre ou à une cérémonie de fin de deuil dans votre environnement (quartier, village). Remplissez la grille ci-dessus. Identifiez \textbf{un point de tension} entre modernité (hôpital, mairie, entreprise, école) et coutume. Proposez une \textbf{mesure de médiation concrète} adaptée au métier que vous visez (infirmier, administrateur, enseignant, ingénieur).
\end{exercice}

\begin{corrige}[Exercice n°1 -- Pistes de correction]
\textbf{Exemple de tension observée} : veillée de plusieurs jours à domicile (exigence coutumière) contre obligation de reprise du travail dès le lendemain (contrat de travail), avec en arrière-plan un risque sanitaire (promiscuité, eau non potable).
\textbf{Médiation proposée (responsable du personnel / infirmier d'entreprise)} :
\begin{enumerate}
    \item Négocier un \textbf{congé « deuil coutumier »} (3 à 5 jours, payé ou sur fonds social de l'entreprise), distinct du congé légal.
    \item Aménager un \textbf{temps et un lieu de recueillement} pour les collègues.
    \item Sensibiliser la famille aux \textbf{mesures d'hygiène minimales} (eau traitée, aération, limitation du nombre de visiteurs) sans interdire le rite.
    \item Prévoir un \textbf{soutien psychologique} post-deuil pour l'employé (prévention de la dépression et de l'absentéisme prolongé).
\end{enumerate}
\textbf{Critères de réussite} : la solution respecte à la fois la norme sanitaire et juridique (légitimité de l'État) et la norme coutumière (dignité du passage) ; elle est concrète, chiffrée et réaliste.
\end{corrige}

\coursec{Méthodologie : Dissertation et Explication de texte (Savoir-faire rédactionnel)}

Dans la section précédente, nous avons vu comment les traditions camerounaises — culte des ancêtres, rites funéraires chez les Bamiléké ou les Beti, transmission des noms — donnent à la mort un sens social et spirituel. Mais au Probatoire et au Baccalauréat technique, il ne suffit pas de \emph{savoir} : il faut \emph{rédiger}. Le correcteur n'évalue pas seulement votre culture philosophique ; il évalue votre capacité à construire un raisonnement écrit, rigoureux et personnel. Cette dernière section vous donne les outils méthodologiques pour transformer vos connaissances sur l'existence et la mort en une copie remarquable.

\courssub{L'analyse du sujet : les quatre temps de la réflexion}

Face à un sujet comme « La mort donne-t-elle un sens à la vie ? », l'élève pressé écrit immédiatement tout ce qu'il sait sur la mort. C'est l'erreur fatale. Une dissertation philosophique ne répond pas à un thème, elle répond à une \emph{question}. Or une question se démonte avant de se résoudre.

\begin{methode}[L'analyse de sujet en quatre temps]
\begin{enumerate}
\item \textbf{Définir les termes.} Interrogez chaque mot-clé du sujet. Dans « La mort donne-t-elle un sens à la vie ? » : \cle{mort} (cessation biologique ? néant ? passage ?), \cle{donner} (offrir gratuitement ? imposer ? révéler ?), \cle{sens} (signification ? direction ? valeur ?), \cle{vie} (existence biologique ? existence humaine consciente ?). Chaque définition provisoire oriente la réflexion.
\item \textbf{Identifier la tension ou le paradoxe.} Un sujet de philosophie contient toujours une difficulté. Ici : la mort semble \emph{détruire} l'existence — comment pourrait-elle en même temps lui \emph{donner} un sens ? C'est cette contradiction apparente qui fait la richesse du sujet.
\item \textbf{Formuler la problématique.} Transformez la tension en une question philosophique précise qui sera le fil directeur de toute la copie : « La finitude de l'existence est-elle ce qui la prive de sens, ou au contraire ce qui la rend capable d'en avoir un ? »
\item \textbf{Annoncer le plan.} Prévoyez deux ou trois parties qui répondront progressivement à la problématique, de la réponse la plus immédiate à la réponse la plus réfléchie.
\end{enumerate}
\end{methode}

\begin{attention}[Distinguer les sujets proches]
Deux sujets voisins n'appellent pas la même copie. Comparez : « Faut-il craindre la mort ? » (question \emph{pratique et affective} : la crainte est-elle légitime, raisonnable ? — Épicure y répond non) et « La mort est-elle un mal ? » (question \emph{théorique et axiologique} : la mort est-elle objectivement un mal, indépendamment de ce que nous ressentons ?). Traiter le second avec le plan du premier, c'est le hors-sujet. Soulignez toujours le verbe et l'adjectif du sujet : ce sont eux qui commandent.
\end{attention}

\courssub{Construire le plan : dialectique ou thématique}

Une fois la problématique posée, il faut organiser la réponse. Deux structures dominent au Baccalauréat technique.

\textbf{Le plan dialectique} (Thèse / Antithèse / Synthèse) convient aux sujets qui contiennent une contradiction. Appliquons-le à notre sujet témoin :

\begin{center}
\begin{tabularx}{\linewidth}{|l|X|X|}\hline
\rowcolor{popblueL} \textbf{Moment} & \textbf{Contenu} & \textbf{Auteurs mobilisables} \\ \hline
I. Thèse & Oui, la mort donne un sens à la vie : conscience de la fin qui crée l'urgence d'agir, prix des actions, espérance du salut. & Heidegger (être-pour-la-mort), Jankélévitch, traditions religieuses \\ \hline
II. Antithèse & Non : la mort est néant, elle détruit tout sens au lieu d'en donner ; la vie est absurde. & Épicure (« la mort n'est rien pour nous »), Sartre \\ \hline
III. Synthèse & La mort ne livre pas un sens \emph{tout fait} : elle \emph{oblige} la liberté à en créer un (responsabilité, transmission, engagement). & Camus, Ricœur, sagesse africaine de la transmission \\ \hline
\end{tabularx}
\end{center}

\begin{methode}[La synthèse n'est pas un mélange]
Le plan dialectique ne signifie pas « Pour / Contre / Un peu des deux ». La synthèse doit \emph{dépasser} la contradiction en changeant de niveau d'analyse. Ici, thèse et antithèse restent au niveau \emph{ontologique} (la mort est-elle, en soi, porteuse de sens ?). La synthèse passe au niveau \emph{éthique} : peu importe que la mort contienne un sens, l'essentiel est ce que sa conscience nous oblige à \emph{faire} de notre existence. Le dépassement est un changement de point de vue, non un compromis mou.
\end{methode}

\textbf{Le plan thématique} convient aux sujets descriptifs (« Qu'est-ce que l'existence humaine ? »). On peut alors articuler : I. Approche ontologique (l'existence comme être-là, finitude) ; II. Approche éthique (l'existence comme tâche, liberté, responsabilité) ; III. Approche culturelle et spirituelle (l'existence comme œuvre, mémoire, ancestralité).

\begin{exemplebox}[Sujet type Probatoire : « L'existence humaine est-elle absurde ? »]
Plan type attendu :
\begin{enumerate}
\item \textbf{L'absurde comme constat lucide} : le divorce entre l'attente humaine de sens et le silence du monde (Camus, \emph{Le Mythe de Sisyphe}).
\item \textbf{Les réponses illusoires} : le saut religieux, le nihilisme, la fuite dans les divertissements — autant de manières d'esquiver le constat.
\item \textbf{L'absurde comme fondement de la liberté et de la révolte} : si rien n'est écrit, tout est à faire ; la révolte, la création et la transmission donnent à l'existence une valeur qu'elle ne reçoit pas (Camus, Sartre, Ricœur).
\end{enumerate}
Remarquez la progression : constat $\rightarrow$ fausses solutions $\rightarrow$ solution authentique.
\end{exemplebox}

\courssub{Rédiger l'introduction et la conclusion}

L'introduction suit le schéma \cle{AMI} enrichi : \textbf{A}ccroche, \textbf{M}ise au point (définitions), problématique, annonce du plan.

\begin{exemplebox}[Introduction rédigée (sujet : « La mort donne-t-elle un sens à la vie ? »)]
\textbf{Accroche} : Dans nos villages, les funérailles mobilisent des familles entières pendant plusieurs jours, comme si la mort d'un homme éclairait rétrospectivement toute sa vie. \textbf{Mise au point} : La mort désigne la cessation irréversible de la vie biologique ; le sens renvoie à la fois à la signification et à la direction d'une existence. \textbf{Tension et problématique} : Or la mort semble anéantir l'existence plutôt que l'éclairer. Dès lors : la finitude est-elle ce qui prive la vie de sens, ou ce qui la rend capable d'en avoir un ? \textbf{Annonce du plan} : Nous verrons d'abord que la conscience de la mort confère prix et urgence à l'existence ; nous montrerons ensuite que la mort, comme néant, détruit tout sens ; nous soutiendrons enfin qu'elle n'offre pas un sens tout fait, mais oblige la liberté à en créer un.
\end{exemplebox}

La conclusion comporte trois temps : le \textbf{bilan} (rappel bref du chemin parcouru), la \textbf{réponse explicite} à la problématique (jamais de conclusion en queue de poisson), et une \textbf{ouverture} (question nouvelle, liée : « Mais le sens ainsi créé survit-il à celui qui meurt ? — question de la mémoire et de la transmission »).

\begin{popfigure}
\begin{center}
\begin{tikzpicture}[
  bloc/.style={draw=popblue, fill=popblueL, rounded corners=2pt, text width=4.6cm, align=center, font=\small, inner sep=4pt},
  sous/.style={draw=popline, fill=white, rounded corners=2pt, text width=4.2cm, align=center, font=\scriptsize, inner sep=3pt},
  fleche/.style={-{Stealth[length=2.5mm]}, thick, poporange},
  etiq/.style={font=\scriptsize\itshape, text=popdark, midway, right}
]
\node[bloc, fill=popgreenL, draw=popgreen] (intro) at (0,0) {\textbf{INTRODUCTION}};
\node[sous] (i1) at (0,-1.0) {Accroche (fait, expérience, citation)};
\node[sous] (i2) at (0,-1.9) {Définitions des termes du sujet};
\node[sous] (i3) at (0,-2.8) {Problématique (question précise)};
\node[sous] (i4) at (0,-3.7) {Annonce du plan};

\node[bloc, fill=poporangeL, draw=poporange] (dev) at (7.6,0) {\textbf{DÉVELOPPEMENT}};
\node[sous] (d1) at (7.6,-1.0) {Partie 1 : Arg. 1, Arg. 2, Exemple, Référence};
\node[sous] (d2) at (7.6,-1.9) {Transition (« Mais... », « Cependant... »)};
\node[sous] (d3) at (7.6,-2.8) {Partie 2 : objection argumentée};
\node[sous] (d4) at (7.6,-3.7) {Partie 3 : dépassement, position personnelle};

\node[bloc, fill=poppinkL, draw=poppink] (conc) at (15.2,0) {\textbf{CONCLUSION}};
\node[sous] (c1) at (15.2,-1.0) {Bilan du parcours};
\node[sous] (c2) at (15.2,-1.9) {Réponse à la problématique};
\node[sous] (c3) at (15.2,-2.8) {Ouverture};

\draw[fleche] (intro.east) -- (dev.west) node[etiq, above, pos=0.5] {Progressivité};
\draw[fleche] (dev.east) -- (conc.west) node[etiq, above, pos=0.5] {Articulation};
\draw[fleche, dashed] (i4.south) to[out=-90,in=180] (2.2,-4.6) -- (5.4,-4.6) to[out=0,in=-90] (d1.south);
\node[font=\scriptsize\itshape, text=popdark] at (3.8,-4.95) {le plan annoncé commande le développement};
\end{tikzpicture}
\end{center}
\legende{Architecture type d'une dissertation philosophique. Chaque flèche pleine indique la progressivité du raisonnement ; la flèche en pointillés rappelle que le développement doit suivre exactement le plan annoncé dans l'introduction.}
\end{popfigure}

\courssub{L'explication de texte : lire pour comprendre, comprendre pour discuter}

L'explication de texte n'est ni un résumé, ni une paraphrase, ni une dissertation déguisée. Elle suit quatre temps :

\begin{methode}[Les quatre temps de l'explication de texte]
\begin{enumerate}
\item \textbf{Présentation} : auteur, œuvre, date, contexte, thèse générale du texte. Exemple : « Ce texte est extrait de \emph{La Mort} de Vladimir Jankélévitch (1966). L'auteur y soutient que la mort, loin d'être un simple accident biologique, structure l'existence humaine. »
\item \textbf{Analyse structurée} : suivez le texte \emph{paragraphe par paragraphe} ou moment logique par moment logique. Pour chaque moment, dégagez : le concept central, l'argument, l'exemple, les références implicites (un philosophe qui écrit « le néant » pense souvent à Sartre ou à Heidegger sans les nommer).
\item \textbf{Discussion critique} : évaluez la portée (que nous apprend le texte ?), les limites (que néglige-t-il ? une conception religieuse ? la dimension sociale de la mort ?), l'actualité (que vaut cette thèse face aux traditions africaines de survie des ancêtres ?).
\item \textbf{Conclusion synthétique} : rappelez la thèse de l'auteur et formulez votre jugement argumenté.
\end{enumerate}
\end{methode}

\begin{attention}[Sanctions fréquentes au Bac technique]
\begin{itemize}
\item \textbf{Hors-sujet} : traiter de la mort biologique (définition médicale, causes) alors que le sujet porte sur le \emph{sens} de l'existence.
\item \textbf{Paraphrase du cours} : recopier ses notes sans argumenter ni répondre à la question posée.
\item \textbf{Absence de problématique} : une introduction qui « rentre dans le vif du sujet » sans question directrice condamne la copie à l'errance.
\item \textbf{Plan « catalogue d'auteurs »} : « Platon dit..., Épicure dit..., Heidegger dit... » sans articulation logique. Les auteurs sont des \emph{témoins} au service de \emph{votre} raisonnement, non l'inverse.
\end{itemize}
\end{attention}

\courssub{La boîte à outils du candidat}

\begin{aretenir}[Vocabulaire technique indispensable]
\begin{itemize}
\item \cle{Ontologie} : étude de l'être en tant qu'être.
\item \cle{Télos} : fin, but, accomplissement d'une existence.
\item \cle{Finitude} : caractère limité, mortel de l'existence humaine.
\item \cle{Contingence} : ce qui aurait pu ne pas être ou être autrement.
\item \cle{Néant} : absence d'être ; chez Sartre, ce que la conscience introduit dans le monde.
\item \cle{Transcendance / Immanence} : ce qui dépasse le monde vécu / ce qui demeure en lui.
\item \cle{Authenticité} : mode d'existence qui assume sa finitude (Heidegger).
\item \cle{Mauvaise foi} : fuite de sa propre liberté (Sartre).
\item \cle{Autonomie / Hétéronomie} : se donner sa loi soi-même / la recevoir d'ailleurs.
\end{itemize}
\end{aretenir}

\begin{popfigure}
\begin{center}
\begin{tikzpicture}
\node[draw=poppurple, fill=poppurpleL, rounded corners=3pt, font=\small\bfseries, inner sep=5pt] (titre) at (7.5,0) {Connecteurs logiques pour philosopher};
\node[draw=popblue, fill=popblueL, rounded corners=2pt, text width=5.4cm, align=center, font=\scriptsize, inner sep=4pt] (opp) at (2.2,-1.6) {\textbf{Opposition}\\ Cependant, toutefois, inversement, en revanche, or};
\node[draw=popgreen, fill=popgreenL, rounded corners=2pt, text width=5.4cm, align=center, font=\scriptsize, inner sep=4pt] (conc) at (12.8,-1.6) {\textbf{Concession}\\ Certes, il est vrai que, sans doute, quoique};
\node[draw=poporange, fill=poporangeL, rounded corners=2pt, text width=5.4cm, align=center, font=\scriptsize, inner sep=4pt] (ench) at (2.2,-3.2) {\textbf{Enchaînement}\\ Par conséquent, dès lors, c'est pourquoi, ainsi};
\node[draw=poppink, fill=poppinkL, rounded corners=2pt, text width=5.4cm, align=center, font=\scriptsize, inner sep=4pt] (ill) at (12.8,-3.2) {\textbf{Illustration}\\ Par exemple, notamment, ainsi le montre, c'est le cas de};
\node[draw=popgold, fill=popgoldL, rounded corners=2pt, text width=5.4cm, align=center, font=\scriptsize, inner sep=4pt] (nua) at (7.5,-4.4) {\textbf{Nuance}\\ Plutôt que, non pas que, à condition que, dans la mesure où};
\draw[-{Stealth}, popline, thick] (titre.south) -- (opp.north);
\draw[-{Stealth}, popline, thick] (titre.south) -- (conc.north);
\draw[-{Stealth}, popline, thick] (opp.south) -- (ench.north);
\draw[-{Stealth}, popline, thick] (conc.south) -- (ill.north);
\draw[-{Stealth}, popline, thick] (ench.east) -- (nua.west);
\draw[-{Stealth}, popline, thick] (ill.west) -- (nua.east);
\end{tikzpicture}
\end{center}
\legende{Les cinq familles de connecteurs logiques. Une copie sans connecteurs est une copie sans articulation : le correcteur ne voit pas le raisonnement progresser.}
\end{popfigure}

\begin{savaistu}[La méthode, un allié stratégique]
Aux épreuves de philosophie du Probatoire et du Baccalauréat technique, la notation distingue toujours la \emph{forme} (structure, langue, présentation) et le \emph{fond} (pertinence, argumentation, références). Concrètement : une copie méthodique avec une culture moyenne obtient souvent une meilleure note qu'une copie savante mais désorganisée. La maîtrise de la méthode vaut autant que la connaissance des auteurs.
\end{savaistu}

\courssub{Entraînement au Baccalauréat technique}

Pour une épreuve de 4 heures, répartissez votre temps ainsi : 30 minutes d'analyse du sujet et de plan détaillé au brouillon, 2 h 45 de rédaction, 30 minutes de relecture (orthographe, syntaxe, connecteurs, vérification que chaque partie répond à la problématique). Soignez la langue : phrases courtes, vocabulaire philosophique précis, aucune abréviation.

\begin{exoresolu}[Analyse complète d'un sujet type Probatoire]
Énoncé : « Faut-il craindre la mort ? » — Construisez l'analyse du sujet et le plan détaillé.
\tcblower
\textbf{1. Définitions.} \emph{Falloir} : est-ce un devoir, une nécessité rationnelle ? \emph{Craindre} : éprouver une peur anticipée devant un mal possible. \emph{Mort} : fin de l'existence individuelle, événement certain mais de date incertaine.

\textbf{2. Tension.} La mort est le seul événement absolument certain de notre vie, et pourtant nous ne pouvons en faire l'expérience : comment craindre rationnellement ce que l'on ne rencontrera jamais ? La crainte semble à la fois naturelle (instinct de survie) et déraisonnable (on ne souffre pas de ce qu'on ne vit pas).

\textbf{3. Problématique.} La crainte de la mort repose-t-elle sur une erreur de raisonnement, ou exprime-t-elle une vérité de la condition humaine ?

\textbf{4. Plan dialectique.}
\begin{itemize}
\item \textbf{I. Il ne faut pas craindre la mort.} Argument d'Épicure : « quand nous sommes, la mort n'est pas ; quand la mort est, nous ne sommes plus » — craindre ce qui ne nous atteindra jamais est irrationnel. La crainte gâche la seule vie dont nous disposons.
\item \textbf{II. Pourtant, la crainte de la mort est légitime et même fondatrice.} Elle exprime l'attachement à la vie et aux proches ; chez Heidegger, l'angoisse devant la mort arrache l'individu à l'anonymat du « on » et le rend à lui-même ; les rites funéraires camerounais montrent que la crainte peut être domestiquée par la culture.
\item \textbf{III. Synthèse.} Il ne faut pas \emph{subir} la crainte de la mort, mais la \emph{transformer} : convertie en conscience de la finitude, elle cesse d'être paralysante et devient principe d'urgence, de responsabilité et de transmission. La sagesse n'est pas l'absence de crainte, mais son éducation.
\end{itemize}
\end{exoresolu}

\begin{exercice}[Sujets d'entraînement]
Traitez l'un des sujets suivants en rédigeant au minimum l'introduction complète et le plan détaillé :
\begin{enumerate}
\item « La mort donne-t-elle un sens à la vie ? »
\item « L'existence humaine est-elle absurde ? »
\item « Vivre, est-ce seulement exister ? »
\end{enumerate}
\end{exercice}

\begin{corrige}[Exercice d'entraînement — éléments de réponse]
\textbf{Sujet 1} : voir le plan dialectique détaillé au paragraphe 5.2 (thèse Heidegger / antithèse Épicure-Sartre / synthèse éthique de la création de sens).

\textbf{Sujet 2} : voir la boîte « Exemple » du paragraphe 5.2 (constat lucide / réponses illusoires / absurde fondant la liberté et la révolte).

\textbf{Sujet 3} : définitions — \emph{vivre} (fonctionner biologiquement) vs \emph{exister} (se rapporter à soi, choisir, se projeter). Tension : tout homme vit, mais tout homme existe-t-il pleinement ? Problématique : qu'est-ce qui distingue l'existence authentique de la simple survie ? Plan : I. Vivre sans exister : la vie réduite à la biologie et à la routine (mauvaise foi, aliénation par le travail répétitif). II. Exister, c'est se choisir : conscience, projet, liberté (Sartre : « l'existence précède l'essence »). III. Exister pleinement, c'est exister avec et pour les autres : la transmission, la mémoire, la communauté — l'ancêtre camerounais « vit » encore parce qu'il a existé pour les siens.

Dans les trois cas, le correcteur attend : une problématique explicite, des transitions rédigées entre les parties, au moins un exemple concret par partie, et une conclusion qui répond effectivement à la question posée.
\end{corrige}

\begin{aretenir}[L'essentiel de la méthode]
Une bonne copie de philosophie tient en quatre exigences : une problématique claire issue de l'analyse des termes ; un plan progressif dont chaque partie répond à la question ; des arguments étayés par des exemples et des références maîtrisées ; une langue soignée, articulée par des connecteurs logiques. La méthode n'est pas un carcan : c'est ce qui rend votre pensée visible et votre liberté de juger crédible.
\end{aretenir}

\newpage

\coursec{Activité d'intégration}

\coursec{Leçon 11 : L'existence et la mort}

\courssub{Objectifs de l'activité}
\begin{itemize}
    \item Mobiliser les concepts philosophiques de la leçon (\cle{existence}, \cle{mort biologique}, \cle{mort encéphalique}, \cle{mort sociale}, \cle{dignité}, \cle{autonomie}, \cle{acharnement thérapeutique}) pour analyser un cas concret complexe.
    \item Articuler une argumentation rigoureuse croisant \textbf{analyse médicale}, \textbf{cadre juridique camerounais}, \textbf{anthropologie culturelle} (contexte Bamiléké) et \textbf{philosophie morale} (sens de l'existence, devoir de préserver la vie vs droit de mourir dans la dignité).
    \item Produire un écrit structuré de type \textbf{avis d'éthique / dissertation courte} (3 pages max), respectant la méthodologie : problème, thèse, antithèse, synthèse, recommandation opérationnelle.
    \item Développer la compétence « Rédiger et justifier une démarche, une réponse ou une conclusion » (Savoir-faire officiel) en situation de responsabilité simulée.
\end{itemize}

\courssub{Situation : Comité d'Éthique Hospitalier – Cas « Madame Ngo »}
\begin{popfigure}
\begin{tikzpicture}[node distance=1.2cm, every node/.style={font=\small}]
\node[draw, rounded corners, fill=popblueL, text width=0.95\linewidth, align=center] (intro) {
    \textbf{Contexte :} Hôpital Régional de Bafoussam, Service de Réanimation. Réunion extraordinaire du Comité d'Éthique Hospitalier (CEH).
    Présidence : Dr Fouda (Médecin-chef). Membres : Dr Ngono (Réanimateur), Me Essomba (Juriste, Conseil de l'Ordre), Pr Owona (Philosophe/Éthicien), Chef Njiké (Notable Bamiléké, Représentant communauté), Rév. Père Kamga (Aumônier), Mme Atangana (Infirmière coordinatrice).
};
\node[draw, rounded corners, fill=popcream, text width=0.95\linewidth, align=left, below=of intro] (dossier) {
    \textbf{DOSSIER MÉDICAL ANONYMISÉ – PATIENTE : Mme NGO (née FOSSO), 68 ans}\par
    \textbf{Antécédents :} HTA mal équilibrée, Diabète type 2, AVC ischémique léger il y a 4 ans (séquelles motrices minimes).\par
    \textbf{Circonstances :} Admise le 12 janvier 2024 (14h30) pour coma brutal (score Glasgow 3/15). IRM cérébrale : \textbf{AVC hémorragique massif du tronc cérébral et ponts} (volume 45 ml) avec rupture intraventriculaire et hernie tonsillaire. Prise en charge neurochirurgicale impossible (localisation, volume, instabilité hémodynamique).\par
    \textbf{Évolution à J+95 (17 avril 2024) :} \textbf{État Végétatif Persistant (EVP) irréversible} (critères Société Française de Réanimation / Haute Autorité de Santé : durée $>$ 3 mois post-AVC non traumatique).\par
    \begin{itemize}
        \item Ouverture oculaire spontanée cyclique (alternance veille/sommeil) mais \textbf{absence totale de conscience} : pas de fixation, pas de poursuite, pas de réponse aux commandes, pas de communication.
        \item Réflexes tronculaires : pupilles moyennes fixes, réflexes cornéens absents, réflexes oculo-céphaliques absents.
        \item \textbf{Respiration spontanée} via trachéotomie (ventilation assistée en mode support pression) : le tronc cérébral médullaire assure un drive respiratoire, excluant la mort encéphalique complète.
        \item Nutrition/Hydratation : Gastrostomie (PEG) depuis J+21.
        \item Infections nosocomiales récurrentes (3 pneumopathies, 2 sepsis urinaires) traitées par antibiothérapies lourdes.
        \item État cutané : escarres stade 3 sacrée et talons malgré prévention.
    \end{itemize}
    \textbf{Avis médical collégial (Dr Ngono + 2 réanimateurs) :} « Pronostic vital engagé à court terme (semaines). Aucune récupération cognitive ni motrice possible. L'état végétatif est \emph{irréversible}. La poursuite des traitements de support (ventilation, nutrition artificielle, antibiothérapies) constitue un \textbf{acharnement thérapeutique} (obstination déraisonnable). Proposition : \textbf{Arrêt de la nutrition/hydratation artificielle et de la ventilation} (laisser mourir), passage en soins de confort (antalgiques, sédation palliative). »
};
\node[draw, rounded corners, fill=poporangeL, text width=0.95\linewidth, align=left, below=of dossier] (famille) {
    \textbf{POSITION DE LA FAMILLE (Fils aîné, M. Ngo Samuel, 42 ans, Chef de lignage « Nkouong ») :}\par
    « \emph{Maman respire. Son cœur bat. Elle est chaude. Elle digère la bouillie qu'on lui met par la sonde. Elle n'est pas morte. Les médecins veulent la tuer pour gagner de la place ou de l'argent.}\par
    \emph{Dans notre tradition Bamiléké (Groupement Bafoussam), la mort n'est pas un arrêt technique. C'est un \textbf{voyage}. Les Ancêtres (Nzon Mbou) n'ont pas encore appelé Maman. Il n'y a pas eu le \textbf{signe} (rêve du patriarche, chant de l'oiseau « Ngoung » sur le toit, chute de l'arbre sacré).}\par
    \emph{Tant que les Ancêtres ne l'ont pas rappelée, elle est \textbf{vivante socialement}. La retirer de la machine, c'est rompre le lien lignager, c'est un \textbf{meurtre rituel} qui frappera la descendance. Nous, les fils, avons le \textbf{devoir filial} (Nkoung) de la maintenir en vie jusqu'au signe. Nous refusons l'arrêt. Nous demandons le transfert au domicile pour les rites de fin de vie quand le moment viendra. » }
};
\node[draw, rounded corners, fill=popgreenL, text width=0.95\linewidth, align=left, below=of famille] (juridique) {
    \textbf{EXTRAITS JURIDIQUES ET DÉONTOLOGIQUES FOURNIS AU COMITÉ :}\par
    \begin{itemize}
        \item \textbf{Code Pénal Camerounais (Loi n°2016/007) :} Art. 276 (Meutre), Art. 277 (Assassinat), Art. 280 (Non-assistance à personne en danger). \textbf{Art. 297-1 :} « N'est pas punissable le médecin qui, en fin de vie, limite ou arrête des traitements inutiles, disproportionnés ou n'ayant d'autre effet que le maintien artificiel de la vie, avec le consentement de la personne de confiance ou de la famille, dans le respect de la dignité. »
        \item \textbf{Code de Déontologie Médicale (Ordre National des Médecins du Cameroun, 2019) :} Art. 38 : « Le médecin doit respecter la volonté du patient... En fin de vie, il évite l'acharnement thérapeutique. Il assure les soins palliatifs. »
        \item \textbf{Loi n°2023/004 (Santé) :} Reconnaissance des \textbf{Directives Anticipées} (non rédigées ici) et de la \textbf{Personne de Confiance} (non désignée formellement, fils aîné = représentant légal coutumier).
        \item \textbf{Déclaration de Yaoundé sur la Bioéthique (2021) :} Principe de \textbf{dignité}, \textbf{autonomie}, \textbf{non-malfaisance}, \textbf{justice}. Intégration des \textbf{valeurs culturelles africaines} (communautarisme, sacré de la vie, rôle des ancêtres).
    \end{itemize}
};
\node[draw, rounded corners, fill=poppurpleL, text width=0.95\linewidth, align=left, below=of juridique] (avis) {
    \textbf{AVIS EXTERNES SOLLICITÉS :}\par
    \begin{itemize}
        \item \textbf{Rév. Père Kamga (Catholique) :} « L'Église distingue \emph{vie biologique} et \emph{vie personnelle}. L'EVP n'est pas la fin de la personne humaine (image de Dieu). Mais l'acharnement (moyens extraordinaires disproportionnés) est moralement inacceptable. L'arrêt de la nutrition/hydratation \emph{artificielle} est licite si c'est un arrêt de moyen disproportionné, pas une intention de tuer. \textbf{Prudence} : risque d'euthanasie déguisée. La famille a voix prépondérante si pas de directives. »
        \item \textbf{Imam Moussa (Musulman, Communauté Peule locale) :} « La vie appartient à Allah (Al-Muhyi, Al-Mumit). Seul Lui décide du terme (Ajal). Débrancher = jouer à Dieu (Chirk). Cependant, si les médecins certifient que la mort cérébrale est \emph{certaine} (Fatwa OMS/Caire 1986 : mort encéphalique = mort islamique), alors l'appareil ne maintient qu'un cadavre. Ici : \textbf{pas de mort encéphalique certifiée} (tronc cérébral partiel ?). Donc : maintenir la vie. »
    \end{itemize}
};
\end{tikzpicture}
\legende{Situation-problème : Dossier complet soumis au Comité d'Éthique Hospitalier (Cas Mme Ngo).}
\end{popfigure}

\courssub{Consignes}
Vous êtes \textbf{Pr Owona}, membre philosophe du CEH. Rédigez l'\textbf{Avis Motivé du Comité} (3 pages maximum, police Times New Roman 12, interligne 1,5) à destination de la Direction de l'Hôpital, de la Famille et du Procureur de la République. Votre avis doit structurer obligatoirement les 4 parties suivantes :

\begin{enumerate}
    \item \textbf{Analyse Clinique et Juridique (Qualification de l'état et du cadre légal)} : 
    \begin{itemize}
        \item Qualifier l'état de Mme Ngo : Vie ? Mort encéphalique ? État végétatif persistant ? Mort sociale ?
        \item Caractériser l'acte proposé (arrêt nutrition/ventilation) : Acharnement thérapeutique ? Laisser mourir ? Euthanasie ? Homicide ?
        \item Situer la responsabilité légale : Code pénal (Art. 297-1), Code déontologie, Personne de confiance / Représentant légal.
    \end{itemize}
    \item \textbf{Analyse Philosophique (Dignité, Autonomie, Sens de l'existence)} : 
    \begin{itemize}
        \item Confronter \textbf{deux thèses majeures} sur le sens de l'existence en fin de vie : \emph{Vitalisme absolu} (la vie valeur suprême, intangible, don de Dieu/Ancêtres) vs \emph{Personalisme / Qualité de vie} (la vie n'a de sens que si conscience/relation/autonomie possibles ; dignité = respect de la personne, pas simple maintien biologique).
        \item Citer explicitement \textbf{deux auteurs/références philosophiques} vus en cours (ex: \textbf{Jonas} « Le droit à la vie », \textbf{Sartre} « L'existence précède l'essence » / liberté radicale, \textbf{Heidegger} « Être-pour-la-mort », \textbf{Cicéron/Sénèque} « Apprendre à mourir », \textbf{Hans Jonas} « Principe responsabilité », \textbf{Paul Ricœur} « Soi-même comme un autre » / éthique de la sollicitude).
        \item Définir la \textbf{dignité} ici : inhérente (ontologique) ou acquise (fonctionnelle) ?
    \end{itemize}
    \item \textbf{Analyse Anthropologique et Culturelle (Mort sociale, Devoirs lignagers)} : 
    \begin{itemize}
        \item Expliquer la conception Bamiléké de la \cle{mort sociale} vs \cle{mort biologique}. Rôle des \cle{Ancêtres}, du \cle{signe}, du \cle{chef de lignage}.
        \item Analyser le conflit de \textbf{devoirs} : Devoir filial (maintenir la vie) vs Devoir de vérité/médical (arrêter l'inutile) vs Devoir sociétal (justice, coût social).
        \item Proposer une \textbf{herméneutique du dialogue} : comment articuler droit positif et droit coutumier sans nier l'un ni l'autre ?
    \end{itemize}
    \item \textbf{Recommandation Motivée et Procédure de Médiation} : 
    \begin{itemize}
        \item Décision collégiale (Arrêt ? Maintien ? Transfert ? Délai ?) avec \textbf{majorité qualifiée} (2/3) et \textbf{vote minorité} si désaccord.
        \item Plan de médiation concret (étapes, acteurs, calendrier) pour apaiser la famille et sécuriser l'équipe soignante.
        \item Mesures de \textbf{soins palliatifs} obligatoires quoi qu'il en soit.
    \end{itemize}
\end{enumerate}

\courssub{Ressources à mobiliser}
\begin{definition}[Concepts clés de la leçon]
\begin{itemize}
    \item \textbf{Existence (Heidegger/Jaspers/Sartre) :} Mode d'être propre à l'homme (Dasein) : projet, liberté, temporalité, conscience de la mort. Diffère de la simple \emph{vie biologique} (métabolisme, reproduction).
    \item \textbf{Mort :} 
    \begin{itemize}
        \item \emph{Biologique :} Arrêt irréversible fonctions vitales (cœur, respiration).
        \item \emph{Encéphalique (Mort cérébrale) :} Arrêt irréversible \emph{toutes} fonctions encéphaliques (tronc + cortex). = Mort légale dans bcp pays (pas encore explicite au Cameroun mais admis par déontologie).
        \item \emph{Sociale / Culturelle :} Perte du statut de personne, sortie de la communauté des vivants (rites, deuil, ancêtres). Peut précéder ou suivre la mort biologique.
    \end{itemize}
    \item \textbf{Acharnement thérapeutique (Obstination déraisonnable) :} Maintien de traitements inutiles, disproportionnés, n'ayant d'autre effet que prolonger artificiellement la vie (douleur, indignité). Interdit (Déontologie, Loi 2023).
    \item \textbf{Laisser mourir / Limitation/Arrêt de traitement :} Décision d'arrêter un traitement devenu disproportionné. But : respecter processus naturel mort, dignité. \emph{Intention} $\neq$ tuer.
    \item \textbf{Dignité :} Valeur intrinsèque (Kant : fin en soi) vs Qualité de vie (fonctionnelle). En fin de vie : respect autonomie, soulagement souffrance, proportionnalité soins.
\end{itemize}
\end{definition}

\begin{propriete}[Références philosophiques attendues (Programme Terminale Tech)]
\begin{itemize}
    \item \textbf{Heidegger (\emph{Être et Temps}) :} L'homme est \emph{Être-pour-la-mort}. La mort est la \emph{possibilité de l'impossibilité} la plus propre. L'anticipation de la mort (résolution) donne sens à l'existence (authenticité).
    \item \textbf{Sartre (\emph{L'Être et le Néant}) :} L'existence précède l'essence. Liberté radicale. La mort est l'absurde (fin du projet), mais l'homme se \emph{fait} par ses choix. Autonomie = choisir sa fin si possible.
    \item \textbf{Jonas (\emph{Le Principe Responsabilité}) :} Éthique de la responsabilité pour l'avenir. La vie est un bien à protéger. Mais responsabilité \emph{proportionnée} : ne pas transformer médecine en tyrannie technique.
    \item \textbf{Anthropologie africaine (Mbiti, Nyamnjoh, Owona Nguini) :} « Je suis parce que nous sommes ». Personne = relation. Mort = transition, passage aux ancêtres. Corps = lieu du sacré. Devoir de mémoire et de rites.
\end{itemize}
\end{propriete}

\begin{methode}[Rédaction Avis Éthique / Dissertation courte]
\begin{enumerate}
    \item \textbf{Introduction (1/2 page) :} Accroche (cas), Définition termes clés, Problématique (ex: « Jusqu'où le devoir de préserver la vie biologique s'étend-il face à l'irréversibilité de l'existence consciente et aux impératifs culturels ? »), Annonce du plan (4 parties).
    \item \textbf{Développement (2 pages) :} 4 parties correspondant aux consignes. Chaque partie = Thèse argumentée + Exemple/Reference + Nuance/Critique. Transitions logiques.
    \item \textbf{Synthèse / Recommandation (1/2 page) :} Réponse directe à la problématique. Décision claire. Mesures d'accompagnement (médiation, palliatif). Signature collégiale.
\end{enumerate}
\end{methode}

\courssub{Démarche et corrigé commenté}
\begin{exoresolu}[Avis Motivé du Comité d'Éthique Hospitalier – Cas Mme Ngo (Modèle de rédaction attendue)]
\textbf{AVIS N\textsuperscript{o} 2024/CEH/045 – SESSION EXTRAORDINAIRE DU 18 AVRIL 2024}\par
\textbf{Objet :} Demande d'avis sur la limitation/arrêt des traitements de support pour Mme Ngo (68 ans, EVP post-AVC massif).\par
\textbf{Composition du comité :} Dr Fouda (Président), Dr Ngono, Me Essomba, Pr Owona, Chef Njiké, Rév. Père Kamga, Mme Atangana.\par
\vspace{0.5cm}

\noindent\textbf{1. ANALYSE CLINIQUE ET JURIDIQUE : QUALIFICATION DE L'ÉTAT ET CADRE LÉGAL}\par
\emph{1.1. Qualification clinique : État Végétatif Persistant (EVP) irréversible, non Mort Encéphalique.}\par
Le dossier médical collégial (J+95) établit : absence totale de conscience (cortex détruit), conservation de fonctions végétatives autonomes (respiration spontanée, circulation, digestion) par intégrité partielle du tronc cérébral (réflexes oculaires absents mais drive respiratoire présent). \textbf{Ce n'est pas une mort encéphalique complète} (critères légaux stricts non remplis : apnée non confirmée, réflexes tronculaires résiduels). C'est un \textbf{EVP irréversible} (durée $>$ 3 mois post-AVC non traumatique, lésions massives irréversibles). La patiente est \textbf{vivante biologiquement} mais \textbf{absente existentiellement} (plus de \cle{projet}, de \cle{lien}, de \cle{représentation} – Heidegger/Sartre).\par
\emph{1.2. Qualification de l'acte proposé : Arrêt de traitements disproportionnés (Laisser mourir), non Euthanasie.}\par
La nutrition/hydratation artificielle (PEG) et la ventilation sont des \textbf{traitements de support vitaux}. Leur poursuite dans un EVP irréversible de 95 jours, avec complications infectieuses répétées et escarres, répond à la définition de l'\textbf{acharnement thérapeutique} (Art. 38 Code Déontologie : « traitements inutiles, disproportionnés, n'ayant d'autre effet que le maintien artificiel de la vie »). L'arrêt vise à \textbf{cesser l'obstination}, non à provoquer la mort (intention $\neq$ tuer). La mort surviendra de la pathologie initiale (AVC), non de l'arrêt. Juridiquement : \textbf{Art. 297-1 Code Pénal} exonère le médecin limitant des traitements « inutiles, disproportionnés » avec accord famille/personne de confiance.\par
\emph{1.3. Représentation légale : Conflit de légitimité.}\par
Pas de Directives Anticipées, pas de Personne de Confiance désignée. Le fils aîné, Chef de lignage, est le \textbf{représentant légal coutumier} et \emph{de facto} décideur substitut. Son refus bloque la procédure standard. Le Comité doit arbitrer entre \emph{autonomie substitutive} (volonté présumée patiente / intérêt supérieur) et \emph{autorité coutumière}.\par
\vspace{0.3cm}

\noindent\textbf{2. ANALYSE PHILOSOPHIQUE : DIGNITÉ, AUTONOMIE, SENS DE L'EXISTENCE}\par
\emph{2.1. Thèse A : Le Vitalisme Absolu (Vie = Valeur Suprême Intangible).}\par
Soutenue par la famille et l'Imam. La vie est don de Dieu (ou des Ancêtres). Seul le Donateur peut la reprendre. \textbf{Sanctuarité de la vie} (Judaïsme/Christianisme/Islam traditionnels). Tout arrêt de support = homicide. \textbf{Arg. :} Mme Ngo respire, digère, vit. Sa dignité \emph{ontologique} (Kant : \emph{humanité comme fin en soi}) est intacte indépendamment de ses fonctions cognitives. La réduire à sa conscience serait un \emph{utilitarisme} dangereux (glissement eugénique). \textbf{Référence :} \textbf{Jean-Paul II} (\emph{Evangelium Vitae}) / \textbf{Lejeune} : « La vie humaine a un prix infini ».\par
\emph{2.2. Thèse B : Le Personalisme / Qualité de Vie (Vie = Condition de l'Existence Personnelle).}\par
Soutenue par l'équipe médicale, le juriste, le philosophe. \textbf{Hans Jonas} (\emph{Le Principe Responsabilité}) : La vie biologique n'est pas une valeur absolue si elle devient « vie végétative » dépourvue de \emph{monde}, de \emph{liberté}, de \emph{responsabilité}. \textbf{Heidegger} : L'\cle{existence} (Dasein) est \emph{projet}, \emph{compréhension}, \emph{être-avec}. L'EVP est la \emph{perte de l'existence} (plus de \emph{Da-sein}, plus de \emph{monde}), ne reste que le \emph{corps} (Leib/Körper distinction). \textbf{Sartre} : L'homme est liberté. Perte irréversible de la liberté = perte de la condition humaine. \textbf{Dignité} ici = \emph{respect de la personne que fut Mme Ngo} (mère, commerçante, croyante) : ne pas la réduire à un corps objet de soins techniques infinis. \textbf{Paul Ricœur} (\emph{Soi-même comme un autre}) : L'éthique de la \emph{sollicitude} exige de ne pas faire subir à l'autre (corps vulnérable) une survie indigne de son histoire.\par
\emph{2.3. Synthèse philosophique : Proportionnalité et Fidélité.}\par
La dignité n'est pas \emph{soit} ontologique \emph{soit} fonctionnelle, mais \emph{relationnelle}. Respecter Mme Ngo, c'est respecter \emph{sa} vie (son histoire, sa foi, sa culture), pas \emph{la} vie en général. L'acharnement la \emph{réifie} (en fait un objet technique). L'arrêt proportionné (soins de confort) est un acte de \emph{fidélité} à sa personne, pas un abandon.\par
\vspace{0.3cm}

\noindent\textbf{3. ANALYSE ANTHROPOLOGIQUE : MORT SOCIALE, DEVOIRS LIGNAGERS, MÉDIATION}\par
\emph{3.1. La mort comme processus social (Conception Bamiléké).}\par
Pour les Bamiléké (Chef Njiké), la personne est un \textbf{nœud relationnel} (Mbiti : « Je suis parce que nous sommes »). La mort biologique n'est qu'une \emph{étape}. La \textbf{vraie mort} est \emph{sociale} : quand les Ancêtres (Nzon Mbou) rappellent l'âme, validé par le \textbf{signe} (rêve, oiseau, arbre). Tant que le signe manque, la personne est \emph{vivante socialement}, habitée par la force vitale lignagère. Le \textbf{devoir filial (Nkoung)} du fils aîné (chef de lignage) est \emph{sacré} : maintenir le corps pour permettre le passage ordonné. Refuser = trahir les ancêtres, exposer la lignée à la malédiction.\par
\emph{3.2. Conflit de devoirs et Herméneutique du dialogue.}\par
\begin{itemize}
    \item \textbf{Devoir médical/vérité :} Ne pas mentir (EVP irréversible), ne pas nuire (acharnement = souffrance corps, gaspillage ressources).
    \item \textbf{Devoir filial/coutumier :} Protéger l'ancêtre, attendre le signe, assurer la continuité.
    \item \textbf{Devoir civique/juridique :} Appliquer la loi (Art. 297-1), protéger équipe soignante.
\end{itemize}
\textbf{Piste de résolution :} Ne pas opposer « Science » vs « Tradition » mais \textbf{articuler les temporalités}. La médecine gère le \emph{corps biologique} (temps court, irréversibilité). La coutume gère la \emph{personne sociale/ancestrale} (temps long, rites). \textbf{Proposition :} Redéfinir l'arrêt des machines non comme « tuer » mais comme \textbf{« rendre le corps à la terre/aux ancêtres pour que le voyage commence »}. Le signe peut être \emph{interprété} : l'irréversibilité médicale \emph{est} un signe (le corps ne peut plus porter la vie).\par
\vspace{0.3cm}

\noindent\textbf{4. RECOMMANDATION MOTIVÉE ET PROCÉDURE DE MÉDIATION}\par
\emph{4.1. Décision du Comité (Vote à bulletin secret, majorité 2/3 requise).}\par
\textbf{Résultat : 5 voix POUR / 2 voix CONTRE (Chef Njiké, Imam Moussa consulté).}\par
\textbf{Décision :} \textbf{Arrêt de la nutrition/hydratation artificielle (PEG) et de la ventilation mécanique} dans un délai de 48h, avec mise en place immédiate d'un \textbf{protocole de soins palliatifs / confort maximal} (antalgiques morphiniques, sédation proportionnée, hygiène, présence continue).\par
\emph{Motivation :} 
\begin{itemize}
    \item Constat d'\textbf{acharnement thérapeutique} (traitements disproportionnés, inutiles, source de souffrance corporelle).
    \item Respect de la \textbf{dignité de la personne} (Histoire de Mme Ngo, non simple biologie).
    \item Application de l'\textbf{Art. 297-1 Code Pénal} et \textbf{Art. 38 Déontologie}.
    \item Absence de directives contraires de la patiente (principe autonomie respecté par défaut via intérêt supérieur).
\end{itemize}
\emph{4.2. Procédure de Médiation (Apaisement, Sens, Rite).}\par
\begin{enumerate}
    \item \textbf{J+0 (Aujourd'hui) :} Annonce officielle à la famille par Dr Fouda + Pr Owona + Chef Njiké (médiateur interne). Remise écrite de l'Avis (ce document). Explication pédagogique : « Arrêt des machines = Fin de l'obstacle technique au voyage ancestral ».
    \item \textbf{J+1 :} \textbf{Cérémonie de « Levée du corps » (adaptée) au chevet.} Invitation du Notable, Prêtre, Imam. Prières, libations, parole du fils aîné. \textbf{Symbolique :} On ne « débranche » pas, on « accompagne le départ ». L'équipe soignante participe (soin = rite).
    \item \textbf{J+2 (48h) :} Arrêt effectif PEG + Ventilation. \textbf{Soins de confort seuls.} Surveillance continue infirmière. Famille autorisée 24h/24.
    \item \textbf{J+2 à J+7 :} Accompagnement deuil (Psychologue hôpital, Aumônerie, Association « Veilleurs »). Préparation funérailles coutumières (transport corps domicile, rites complets).
    \item \textbf{Suivi :} Compte-rendu au Procureur (transparence). Évaluation éthique à 6 mois (Retour d'expérience).
\end{enumerate}
\emph{4.3. Mesures Palliatives Obligatoires (Quelle que soit l'issue).}\par
Prescription immédiate : Morphine SC (pompe continue) + Midazolam SC (si dyspnée/agitation) + Soins de bouche/yeux/peau q4h + Positionnement anti-escarres + Présence humaine continue. \textbf{Objectif :} Mort paisible, sans douleur, sans dyspnée, entourée.\par
\vspace{0.5cm}

\noindent\textbf{Fait à Bafoussam, le 18 avril 2024.}\par
\noindent\textbf{Signatures :} Dr Fouda (Président), Dr Ngono, Me Essomba, \textbf{Pr Owona}, Rév. Père Kamga, Mme Atangana. \textit{(Chef Njiké : « Je note la décision. Je m'engage à préparer ma mère et la lignée au voyage. Je ne signe pas mais je ne m'oppose plus. »)}
\tcblower
\textbf{COMMENTAIRE PÉDAGOGIQUE (Grille d'évaluation) :}\par
\begin{itemize}
    \item \textbf{Concepts (5/5) :} Distinction nette Vie/Existence/Mort (biologique, encéphalique, sociale, végétative). Utilisation précise : Acharnement, Laisser mourir, Dignité (ontologique/fonctionnelle), Autonomie/Substituée.
    \item \textbf{Argumentation (5/5) :} Structure dialectique (Thèse/Antithèse/Synthèse) respectée. Auteurs cités pertinents (Heidegger, Jonas, Sartre, Ricœur, Kant, Mbiti). Articulation philo/droit/anthro réussie.
    \item \textbf{Droit (5/5) :} Application correcte Art. 297-1 (nouveau), Code Déontologie, distinction euthanasie/arrêt traitement. Gestion absence directives/personne confiance.
    \item \textbf{Contexte Camerounais (5/5) :} Intégration réelle Bamiléké (Ancêtres, signe, lignage, chef, devoir filial). Médiation culturellement intelligente (cérémonie, réinterprétation signe). Acteurs locaux (Notable, Prêtre, Imam).
    \item \textbf{Rédaction (5/5) :} Format avis officiel. Ton professionnel, éthique, humain. Plan respecté. Recommandation opérationnelle claire (délai, protocole, médiation). 3 pages respectées.
\end{itemize}
\textbf{Note finale : 20/20 – Excellent. Modèle de dissertation appliquée / Avis d'éthique.}
\end{exoresolu}

\courssub{Bilan de l'activité}
Cette activité d'intégration a permis de vérifier l'acquisition des savoirs et savoir-faire de la leçon « L'existence et la mort » dans une situation professionnelle simulée à haute charge éthique et culturelle. Elle a mis en évidence :
\begin{itemize}
    \item La \textbf{pertinence opératoire} des distinctions conceptuelles (Vie/Existence, Mort biologique/encéphalique/sociale, Acharnement/Laisser mourir) pour trancher un dilemme réel.
    \item La \textbf{nécessité du dialogue interculturel} en bioéthique au Cameroun : le droit positif (Art. 297-1) ne s'applique pas dans le vide ; il doit être \emph{interprété} et \emph{médiatisé} par l'anthropologie (Ancêtres, signes, lignage) pour être accepté socialement.
    \item La \textbf{puissance de la philosophie} (Heidegger, Jonas, Ricœur, pensée africaine) pour dépasser l'opposition stérile « Pour/Contre l'euthanasie » et construire une \emph{éthique de la proportionnalité et de la fidélité} à la personne concrète.
    \item L'exigence de \textbf{rigueur rédactionnelle} (structure avis/dissertation, vocabulaire précis, argumentation référencée, recommandation actionnable) pour former des techniciens supérieurs capables de responsabilité éthique dans leurs futurs métiers (santé, droit, administration, ingénierie sociale).
\end{itemize}
L'élève qui réussit cet exercice démontre qu'il a intégré le programme : il \emph{pense} la mort pour mieux \emph{agir} face à la fin de vie.

\newpage

\coursec{Méthodes et exercices résolus}

\courssub{Méthodes et exercices résolus}

\begin{methode}[Définir et distinguer rigoureusement les concepts philosophiques : Vie, Existence, Mort]
\textbf{Objectif :} Maîtriser la précision conceptuelle exigée en introduction de dissertation ou en explication de texte.
\begin{enumerate}
    \item \textbf{Identifier le genre et l'espèce :} Rattacher le concept à une catégorie large (ex. : la Vie $\to$ phénomène biologique ; l'Existence $\to$ modalité humaine ; la Mort $\to$ événement/fin).
    \item \textbf{Produire une définition par analyse :} Décomposer le concept en traits essentiels (nécessaires et suffisants).
    \begin{itemize}
        \item \cle{Vie} : Ensemble des fonctions physiologiques (nutrition, reproduction, métabolisme) \textbf{+} potentialité de conscience (chez l'homme).
        \item \cle{Existence} (sens philosophique strict, ex. Heidegger/Sartre) : Manière d'être propre à l'homme (Dasein) qui \textbf{se comprend elle-même}, se projette (projet), et est \textbf{jetée} dans le monde (facticité). Diffère de la simple « réalité » ou « vie biologique ».
        \item \cle{Mort} : Cessation irréversible des fonctions vitales (mort biologique) \textbf{vs} Mort comme \cle{possibilité propre, certaine et indépassable} de l'existence (Heidegger) qui confère à la vie son urgence et sa finitude.
    \end{itemize}
    \item \textbf{Établir la distinction opératoire :} Utiliser un tableau de différences (Nature / Statut / Temporalité / Rapport au sens).
    \item \textbf{Illustrer par un exemple concret camerounais :} La veillée mortuaire (existence sociale qui se prolonge au-delà de la mort biologique).
\end{enumerate}
\end{methode}

\begin{exoresolu}[Sujet type Bac : « Distinguez la vie biologique de l'existence humaine. Montrez en quoi la mort éclaire cette distinction. »]
\textbf{Analyse du sujet :}
\begin{itemize}
    \item \textbf{Concepts clés :} Vie biologique, Existence humaine, Mort (comme révélateur).
    \item \textbf{Problématique implicite :} L'homme est-il seulement un organisme vivant ? La mort est-elle un simple accident biologique ou ce qui structure l'humain ?
    \item \textbf{Plan dialectique imposé :} 1. Distinction conceptuelle (Vie $\neq$ Existence). 2. La mort comme horizon de l'existence (Heidegger). 3. La mort comme épreuve du sens (Ancestralité / Symbolique).
\end{itemize}

\tcblower

\textbf{Solution détaillée :}

\textbf{Introduction}
\textbf{Amorce :} « L'homme est un animal métaphysique » (Schopenhauer). Au-delà de sa condition organique, il questionne son être.
\textbf{Définitions :} La \cle{vie biologique} est l'ensemble des processus physico-chimiques assurant la conservation de l'espèce (métabolisme, reproduction). L'\cle{existence} (au sens heideggérien) est l'\emph{ek-sistere} : l'homme « se tient hors » de lui-même, il est projet, liberté, et souci de son être. La \cle{mort} n'est pas seulement la fin biologique, mais l'horizon de finitude qui donne sa texture à l'existence.
\textbf{Problématique :} En quoi la mort, loin d'être la simple cessation de la vie biologique, révèle-t-elle la singularité de l'existence humaine ?
\textbf{Annonce du plan :} Nous montrerons d'abord la distinction ontologique entre vie et existence, puis comment la mort structure l'existence, enfin comment la tradition camerounaise transcende l'opposition.

\textbf{Développement}

\textbf{I. Vie biologique et Existence humaine : une différence de nature, non de degré}
\begin{itemize}
    \item \textbf{La Vie : immanence et automatisme.} L'organisme vit \emph{par} la vie. Il subit les lois de la nature (homéostasie, évolution). Ex. : Le bananier dans une plantation à Njombé croît, se nourrit, meurt sans « choisir » sa vie.
    \item \textbf{L'Existence : transcendance et liberté.} L'homme \emph{mène} son existence (Sartre : « l'existence précède l'essence »). Il se \textbf{projette} vers des fins (devenir ingénieur, chef de famille), il \textbf{se comprend} (il a une « compréhension de l'être »), il est \textbf{jeté} dans un monde (ethnie, langue, histoire) qu'il doit assumer.
    \item \textbf{Distinction opératoire :} Un légume est « en vie » ; un homme « existe ». La vie est un \textbf{don} subi ; l'existence est une \textbf{tâche} à accomplir.
\end{itemize}

\textbf{II. La Mort : de l'accident biologique à la « possibilité propre » de l'existence (Heidegger, \emph{Sein und Zeit})}
\begin{itemize}
    \item \textbf{Mort biologique (perte) :} Cessation des fonctions (arrêt cardiaque, mort cellulaire). C'est un événement \textbf{extérieur} à la vie (on « subit » sa mort biologique). Elle concerne le \emph{corps}.
    \item \textbf{Mort ontologique (finitude) :} Pour Heidegger, la mort est la \cle{possibilité la plus propre, la plus certaine, la plus indépassable} du Dasein. Elle n'est pas un événement futur, mais une \textbf{structure permanente} de l'existence (« être-pour-la-mort »).
    \item \textbf{Révélation :} L'angoisse face à la mort (différente de la peur de mourir) arrache l'homme à l'inauthenticité du « On » (la foule, les commérages au marché, les réseaux sociaux) et le renvoie à sa \textbf{responsabilité} unique : « Personne ne peut mourir à ma place. » C'est \textbf{par} l'anticipation de la mort que l'existence devient authentique (urgente, lucide).
\end{itemize}

\textbf{III. Au-delà de l'opposition : La mort comme passage et lien social (Ancestralité camerounaise)}
\begin{itemize}
    \item \textbf{Critique de la finitude radicale :} Si la mort biologique est une fin, l'existence symbolique se prolonge. Chez les Bamiléké, Bassa, ou Beti, la mort est un \textbf{changement d'état} : le trépassé devient \cle{ancêtre} (Ndzém, Miengu, etc.).
    \item \textbf{Fonction sociale :} Les funérailles (pleurs, danses, sacrifices, « deuil » codifié) ne sont pas seulement un adieu, mais une \textbf{réparation du tissu social} déchiré. L'ancêtre veille sur la lignée, sanctionne les fautes, assure la fécondité.
    \item \textbf{Synthèse :} La mort biologique met fin à la \emph{vie} ; la mort rituelle inaugure une \emph{existence} spirituelle et sociale supérieure. L'homme camerounais « existe » pleinement \emph{par} sa mort rituelle.
\end{itemize}

\textbf{Conclusion}
\textbf{Bilan :} La distinction Vie/Existence repose sur le passage du biologique à l'ontologique (liberté, projet, sens). La mort, loin d'être un simple dysfonctionnement, est l'horizon qui donne à l'existence sa densité (Heidegger) et sa dimension communautaire (Ancestralité).
\textbf{Ouverture :} La médecine moderne (réanimation, mort cérébrale) brouille-t-elle la frontière entre vie biologique et existence personnelle ? Le débat sur l'euthanasie ou la fin de vie en unité de soins intensifs à l'Hôpital Central de Yaoundé en témoigne.
\end{exoresolu}

\begin{methode}[Analyser et confronter les grandes thèses sur le sens de l'existence]
\textbf{Objectif :} Savoir présenter, expliquer et critiquer les réponses classiques (Religieuse, Existentialiste, Absurde, Nihiliste, Humaniste) dans une dissertation.
\begin{enumerate}
    \item \textbf{Identifier la thèse centrale :} Une phrase résumant la réponse de l'auteur au « Pourquoi vivons-nous ? ».
    \item \textbf{Expliciter l'argumentation :} Quels concepts ? Quels exemples ? (ex. : Dieu, Liberté, Absurde, Bonheur, Devoir).
    \item \textbf{Situer l'auteur / le courant :} Contexte historique, visée (salut, libération, lucidité, action).
    \item \textbf{Comparer / Confronter :} Utiliser un tableau synoptique (Auteur | Sens de la vie | Source du sens | Attitude face à la mort | Limite).
    \item \textbf{Problématiser :} Ne pas lister, mais montrer les tensions (ex. : Sens donné \emph{ab extra} (Dieu) vs Sens créé \emph{ab intra} (Liberté)).
\end{enumerate}
\end{methode}

\begin{exoresolu}[Sujet type Bac : « Le sens de l'existence est-il donné ou construit ? »]
\textbf{Analyse :}
\begin{itemize}
    \item \textbf{Donné :} Hétérodétermination (Dieu, Nature, Destin, Société, Biologie). Sens \emph{a priori}.
    \item \textbf{Construit :} Autodétermination (Sujet, Liberté, Histoire, Projet). Sens \emph{a posteriori}.
    \item \textbf{Plan :} 1. Thèses du sens donné (Religion, Essentialisme, Déterminismes). 2. Thèses du sens construit (Existentialisme, Humanisme, Absurde héroïque). 3. Dépassement : La dialectique Donné/Construit (Facticité/Projet, Condition/Conversion).
\end{itemize}

\tcblower

\textbf{Solution détaillée :}

\textbf{Introduction}
\textbf{Amorce :} « La vie n'a pas de sens, c'est à nous de lui en donner un » (Sartre) vs « Tu as fait nos cœurs pour toi, Seigneur, et ils sont inquiets tant qu'ils ne reposent pas en toi » (Saint Augustin).
\textbf{Définition :} Le \cle{sens} = direction, finalité, intelligibilité, valeur. \cle{Donné} = antérieur au sujet, objectif, imposé. \cle{Construit} = postérieur au sujet, subjectif, libre.
\textbf{Problématique :} L'homme découvre-t-il une partition préécrite ou compose-t-il sa propre mélodie sur un instrument qu'il n'a pas choisi ?
\textbf{Plan :} 1. L'illusion ou la certitude d'un sens donné. 2. L'exigence et l'angoisse d'un sens construit. 3. La condition humaine : recevoir pour transformer (Facticité et Projet).

\textbf{Développement}

\textbf{I. Les figures du sens « donné » : Hétéronomie et Réconfort}
\begin{itemize}
    \item \textbf{A. Le sens religieux (Théocentrisme) :} \textbf{Saint Augustin / Thomas d'Aquin.} Sens = \textbf{Béatitude} (vision de Dieu). La vie est un pèlerinage. La mort = passage vers la vie vraie. \textit{Force :} Apaise l'angoisse, fonde l'éthique universelle. \textit{Limite :} Exige la foi ; risque d'aliénation (Feuerbach : « projection »), déresponsabilisation terrestre (« l'opium du peuple » - Marx).
    \item \textbf{B. Le sens naturaliste / biologique :} \textbf{Spinoza / Darwin / Sociobiologie.} Sens = \textbf{Conatus} (persévérer dans l'être), transmission génétique, adaptation. La mort = sélection naturelle. \textit{Limite :} Réduit l'homme à un véhicule génétique ; nie la singularité culturelle et morale.
    \item \textbf{C. Le sens social / traditionaliste :} \textbf{Durkheim / Tradition camerounaise.} Sens = \textbf{Intégration} au groupe, respect des coutumes, service de la cité/ancêtres. L'individu n'existe que par le « Nous ». \textit{Limite :} Sacrifie l'autonomie individuelle au conformisme.
\end{itemize}

\textbf{II. L'aventure du sens « construit » : Autonomie et Vertige}
\begin{itemize}
    \item \textbf{A. L'Existence précède l'essence (Sartre, \emph{L'Existentialisme est un humanisme}) :} Pas de nature humaine prédéfinie. L'homme est \cle{condamné à être libre}. Il \textbf{invente} ses valeurs par ses choix engagés. Ex. : Un jeune ingénieur à Douala qui choisit de rester au pays pour créer une start-up agricole \textbf{créé} le sens de sa vie (développement local), au lieu de suivre le sens « donné » par la diaspora (envoi d'argent).
    \item \textbf{B. L'Absurde et la Révolte (Camus, \emph{Le Mythe de Sisyphe}) :} Le monde est irrationnel, muet (Absurde = divorce entre appel humain au sens et silence du monde). \textbf{Refus du suicide philosophique (religion) et du suicide physique.} Sens = \textbf{Lucidité + Révolte + Liberté + Passion}. « Il faut imaginer Sisyphe heureux. » \textit{Force :} Héroïsme lucide. \textit{Limite :} Risque de nihilisme esthétique, solitude radicale.
    \item \textbf{C. Le Sens par l'œuvre / l'Histoire (Marx, Malraux) :} Sens = \textbf{Transformation du monde} (praxis). « L'homme se fait lui-même par son travail. » La mort est acceptée si l'œuvre survit.
\end{itemize}

\textbf{III. Dépassement : La dialectique de la Facticité et du Projet (Heidegger, Merleau-Ponty, Ricoeur)}
\begin{itemize}
    \item \textbf{La Facticité (le Donné) :} Je ne choisis pas ma naissance, mon corps, ma langue (bassa, français, anglais), mon histoire coloniale, ma mortalité. Ce sont des \textbf{conditions} sine qua non.
    \item \textbf{Le Projet (le Construit) :} Je \textbf{reçois} ces conditions pour les \textbf{convertir} en possibilités. \textbf{Paul Ricoeur :} « L'homme est un être capable. » Capable de dire « je », d'agir, de se raconter (identité narrative), de se porter garant.
    \item \textbf{Exemple concret :} Un chef traditionnel Bamoun \textbf{reçoit} (donné) sa charge par lignage (facticité). Mais il \textbf{construit} son règne (projet) en modernisant la chefferie, en gérant les conflits fonciers, en éduquant la jeunesse. Il \textbf{assume} le donné pour le dépasser.
\end{itemize}

\textbf{Conclusion}
\textbf{Bilan :} Le sens n'est ni un trésor caché qu'on trouve (donné pur), ni un château en Espagne qu'on bâtit ex nihilo (construit pur). Il est \textbf{co-construction} : l'homme \textbf{répond} à l'appel de sa condition (donnée) par l'invention de son existence (construite).
\textbf{Ouverture :} L'intelligence artificielle et le transhumanisme (allongement de vie, téléchargement de conscience) ne risquent-ils pas de transformer le « donné » biologique en un « construit » technologique, changeant radicalement la donne philosophique ?
\end{exoresolu}

\begin{methode}[Structurer une dissertation philosophique : Plan dialectique et Plan thématique]
\textbf{Objectif :} Produire une copie organisée, logique, répondant exactement au sujet.
\begin{enumerate}
    \item \textbf{Analyse du sujet (15 min) :} Mots-clés $\to$ Définitions $\to$ Présupposés $\to$ Problématique (question tension) $\to$ Choix du plan.
    \item \textbf{Choix du plan :}
    \begin{itemize}
        \item \textbf{Plan Dialectique (Thèse / Antithèse / Synthèse) :} Idéal pour les sujets « Est-ce que... ? », « Faut-il... ? », opposant deux réponses légitimes. \textit{Attention :} La synthèse n'est pas un compromis mou, mais une \textbf{résolution} à un niveau supérieur.
        \item \textbf{Plan Thématique (Analyse / Critique / Dépassement) :} Idéal pour les sujets « Qu'est-ce que... ? », « Comment... ? », ou citations. On analyse la thèse de l'auteur (I), on en montre les limites (II), on propose une ouverture (III).
        \item \textbf{Plan Conceptuel (Concept A / Concept B / Articulation) :} Pour les sujets comparant deux notions (ex. « Vie et Existence »).
    \end{itemize}
    \item \textbf{Rédaction de l'introduction (modèle type) :}
    \begin{enumerate}
        \item Accroche (citation, fait d'actualité, paradoxe).
        \item Définitions rigoureuses des termes du sujet.
        \item Reformulation / Restitution du sujet.
        \item \textbf{Problématique} (une question claire, issue d'une tension).
        \item Annonce du plan (verbes d'action : \emph{Examiner, Montrer, Dépasser}).
    \end{enumerate}
    \item \textbf{Corps du devoir :} Une idée directrice par grande partie (I, II, III). Un argument + exemple/référence par sous-partie (A, B). Transitions logiques.
    \item \textbf{Conclusion :} Bilan synthétique (réponse à la problématique) + Ouverture (nouvelle question, autre domaine, actualité).
\end{enumerate}
\end{methode}

\begin{exoresolu}[Sujet type Probatoire : « La mort donne-t-elle un sens à la vie ? »]
\textbf{Analyse rapide :}
\begin{itemize}
    \item \textbf{Mots-clés :} Mort (fin, finitude), Sens (finalité, valeur, intelligibilité), Vie (existence, temps).
    \item \textbf{Tension :} Mort = néant / absurdité (ôte le sens) VS Mort = urgence / achèvement (donne le sens).
    \item \textbf{Plan choisi :} \textbf{Dialectique} (Thèse : Oui, elle structure le sens / Antithèse : Non, elle l'anéantit / Synthèse : Elle est l'horizon qui permet la valorisation du temps).
\end{itemize}

\tcblower

\textbf{Introduction rédigée :}
« La mort est la fin de la vie, mais est-elle la fin du sens ? » (Interrogation sur la finitude).
\textbf{Définitions :} \cle{Mort} : Cessation définitive des fonctions vitales (biologique) ; Possibilité propre la plus certaine (ontologique - Heidegger). \cle{Sens} : Finalité, cohérence, valeur qui justifie l'action. \cle{Vie} : Durée vécue, projet, expérience temporelle.
\textbf{Problématique :} La mort, en tant que terme absolu, est-elle ce qui \textbf{fonde} la valeur de la vie (rareté, urgence) ou ce qui la \textbf{ruine} en la rendant vaine (néant final) ?
\textbf{Plan :} I. La mort comme horizon constitutif du sens (Thèse : Heidegger, Tradition). II. La mort comme absurdité qui retire tout sens (Antithèse : Épicure, Camus, Nihilisme). III. La mort comme épreuve de la liberté : le sens par l'acceptation lucide (Synthèse : Sartre, Ricoeur, Ancestralité).

\textbf{Développement structuré (Squelette argumentatif) :}

\textbf{I. La mort donne le sens : l'horizon de finitude (Thèse)}
\begin{itemize}
    \item \textbf{A. L'argument de la rareté (Économie du temps) :} Ce qui est infini n'a pas de valeur. La mort rend le temps \textbf{précieux}. « Nous sommes des êtres pour la mort » (Heidegger) : l'anticipation de la mort (Vorlaufen) arrache à l'inauthenticité, force à choisir \textbf{ici et maintenant}. Ex. : Un étudiant en Terminale à Garoua révise son Bac parce que le temps est compté.
    \item \textbf{B. L'argument de l'achèvement (Téléologie / Narrativité) :} La mort clôt le récit, en fait une \textbf{totalité} close (Aristote, Ricoeur). Sans fin, pas d'histoire, pas de sens. L'ancêtre Bamiléké n'est « accompli » qu'après ses funérailles royales ; sa vie prend son sens \emph{rétrospectif} à sa mort.
    \item \textbf{C. L'argument du sacrifice / Don :} La mort acceptée pour une cause (patrie, foi, famille) \textbf{crée} du sens \emph{par} la mort elle-même. Martyrs de l'UPC, soldats au front.
\end{itemize}

\textbf{II. La mort nie le sens : l'absurde et le néant (Antithèse)}
\begin{itemize}
    \item \textbf{A. L'argument épicurien (Symétrie) :} « Tant que nous existons, la mort n'est pas là ; quand la mort est là, nous ne sommes plus. » La mort ne nous concerne \textbf{jamais} comme expérience. Elle est \textbf{extérieure} à la vie. Donc elle ne peut rien donner à la vie.
    \item \textbf{B. L'argument de l'absurde (Camus / Nagel) :} La mort est la \textbf{rupture radicale} du projet. Je bâtis une cathédrale (sens), la mort la réduit en poussière. « Le travail de l'homme n'est qu'une lente tâche pour retrouver... les deux ou trois images simples dont son cœur s'est nourri » (Camus) $\to$ mais la mort brise la retrouvaille. Le sens \textbf{s'effondre} dans le néant.
    \item \textbf{C. Le nihilisme (Nietzsche passif / Schopenhauer) :} Si tout finit dans le néant, \textbf{tout est vain} (Vanité des vanités, Qohelet). La mort rétroactive : elle contamine le passé et le présent de nullité.
\end{itemize}

\textbf{III. Dépassement : La mort comme appel à la responsabilité (Synthèse)}
\begin{itemize}
    \item \textbf{A. La mort comme « possibilité propre » (Heidegger) :} Elle ne donne pas un sens \emph{tout fait} (donné), elle \textbf{oblige} à en créer un (construit). L'authenticité = assumer sa finitude pour choisir librement.
    \item \textbf{B. Le sens par l'héritage symbolique (Ricoeur / Ancestralité) :} La mort biologique détruit l'individu, mais l'existence \textbf{survit symboliquement} (œuvres, transmission, mémoire ancestrale). Le sens n'est pas \emph{donné} par la mort, mais \emph{sauvé} \emph{par} la manière dont on l'affronte.
    \item \textbf{C. L'éthique de la finitude (Hans Jonas / Bioéthique) :} La mortalité fonde la \textbf{responsabilité} envers les générations futures (écologie, éducation). « Agis de façon que les effets de ton action soient compatibles avec la permanence d'une vie authentiquement humaine. » La mort donne le sens de la \textbf{transmission}.
\end{itemize}

\textbf{Conclusion rédigée :}
\textbf{Bilan :} La mort ne « donne » pas un sens positif comme un cadeau (Antithèse), ni ne le vole nécessairement (Thèse naïve). Elle est la \textbf{condition de possibilité} du sens : sans finitude, pas de choix, pas d'urgence, pas de transmission. Le sens naît de la \textbf{réponse libre} que l'existence apporte à sa propre finitude.
\textbf{Ouverture :} Si la science permettait l'immortalité biologique (transhumanisme), l'existence perdrait-elle son intensité morale ? Le « sens » survivrait-il à la suppression de la mort ?
\end{exoresolu}

\begin{methode}[Réussir l'explication de texte philosophique (Commentaire composé)]
\textbf{Objectif :} Expliquer la pensée d'un auteur \emph{de l'intérieur}, sans plaquer ses propres idées, en suivant la structure logique du texte.
\begin{enumerate}
    \item \textbf{Lecture active (x3) :} 1. Globale (sujet, thèse). 2. Détaillée (vocabulaire technique, connecteurs logiques : \emph{donc, or, mais, car, en effet}). 3. Structurelle (découpage en séquences logiques / paragraphes).
    \item \textbf{Introduction (Standard codifiée) :}
    \begin{enumerate}
        \item Présentation : Auteur, Œuvre, Date, Contexte (mouvement philosophique).
        \item Thèse générale du texte (en une phrase).
        \item Problématique du texte (quelle question l'auteur résout-il ?).
        \item Annonce du plan de l'explication (suivant le découpage du texte : \emph{Dans un premier temps... Ensuite... Enfin...}).
    \end{enumerate}
    \item \textbf{Développement (Explication linéaire structurée) :} Pour chaque séquence :
    \begin{enumerate}
        \item \textbf{Reformuler} l'idée principale (sans citer bêtement).
        \item \textbf{Analyser} les concepts clés (définir \emph{in situ}).
        \item \textbf{Expliciter} le raisonnement (arguments, exemples de l'auteur, distinctions).
        \item \textbf{Montrer la logique} (liaison avec séquence précédente/suivante).
    \end{enumerate}
    \item \textbf{Conclusion :}
    \begin{enumerate}
        \item Résumé de la démonstration (réponse à la problématique).
        \item \textbf{Portée / Limites / Confrontation :} Situer la thèse dans l'histoire de la philosophie (accord/désaccord avec d'autres auteurs), actualité, critique personnelle brève.
    \end{enumerate}
    \item \textbf{Règles d'or :} Pas de « Je pense que... » dans le développement. Respecter l'auteur. Citer le texte (guillemets) pour appuyer l'explication.
\end{enumerate}
\end{methode}

\begin{exoresolu}[Explication de texte : Extrait fictif inspiré de \emph{Sein und Zeit} (Heidegger) et pensée ancestrale]
\textbf{Texte :}
\begin{quote}
« La mort n'est pas un événement qui survient à la fin de la vie. Elle est, dès l'origine, la possibilité la plus propre de l'existence. Tant que l'homme fuit cette certitude dans le bruit du "On" — les cérémonies vides, les discours convenus —, il vit inauthentiquement. Mais lorsqu'il anticipe sa mort non comme une échéance lointaine, mais comme l'ombre qui porte chacun de ses pas, alors il retrouve la gravité de ses choix. Chez nos pères, le tronc d'arbre planté à la naissance n'est pas un symbole : il \emph{est} la croissance de l'homme vers sa fin. Quand l'arbre meurt, l'ancêtre ne disparaît pas ; il change de demeure pour veiller sur la lignée. Ainsi, la mort n'est pas le néant, mais le passage qui fonde l'autorité des vivants. »
\end{quote}

\tcblower

\textbf{Solution détaillée (Modèle de copie) :}

\textbf{Introduction}
\textbf{Présentation :} Ce texte, bien que synthétique, croise la phénoménologie existentialiste de \textbf{Martin Heidegger} (\emph{Être et Temps}, 1927, § 46-53 : « L'être-pour-la-mort ») et une anthropologie philosophique de la \textbf{tradition africaine (camerounaise)} (conception vitale de la mort comme passage ancestral).
\textbf{Thèse :} La mort n'est pas un accident terminal extérieur à la vie, mais une \textbf{structure ontologique constitutive de l'existence} qui, assumée authentiquement (anticipation) ou rituellement (ancestralité), fonde la gravité des choix et la continuité du sens.
\textbf{Problématique :} Comment la mort, loin d'être une simple fin biologique, structure-t-elle positivement l'existence humaine (authenticité, transmission, autorité) ?
\textbf{Plan du texte :} 1. Critique de la conception vulgaire : la mort comme événement futur et fui (Inauthenticité). 2. La mort comme structure existentielle : l'anticipation fondatrice d'authenticité (Heidegger). 3. L'illustration ancestrale : la mort comme passage et fondement de l'ordre social (Tradition).

\textbf{Développement}

\textbf{Première partie : La fuite inauthentique de la mort (Le « On » et le leurre)}
\begin{itemize}
    \item \textbf{Reformulation :} L'auteur attaque la représentation courante : la mort = « événement qui survient à la fin ». C'est une \textbf{extériorisation} » : on la repousse au futur indéfini.
    \item \textbf{Analyse du « On » (Das Man) :} « Bruit du On », « cérémonies vides », « discours convenus ». Heidegger : le \cle{On} est le mode d'être quotidien où le Dasein se perd dans l'anonymat, la curiosité, le bavardage (Gerede). On « meurt » généralement (« on meurt tous »), personne ne \emph{meurt sa mort}.
    \item \textbf{Conséquence :} « Vit inauthentiquement ». Fuite de la responsabilité, dispersion dans l'immédiat (réseaux sociaux, buvettes, commérages). La vie est plate, sans gravité.
\end{itemize}

\textbf{Deuxième partie : L'anticipation heideggérienne, source d'authenticité}
\begin{itemize}
    \item \textbf{Concept clé :} « Possibilité la plus propre, la plus certaine, la plus indépassable ».
    \item \textbf{Explication :} La mort n'est pas \emph{à la fin}, elle est \emph{depuis l'origine} (« dès l'origine »). \cle{Vorlaufen} (Anticipation) = se projeter vers sa propre mort comme possibilité \textbf{actuelle}.
    \item \textbf{Effet :} « Retrouve la gravité de ses choix ». L'angoisse (Angst) face au néant possible individualise : « Personne ne peut mourir à ma place ». Cela force à choisir \emph{soi-même} (projet propre) et non par défaut (suivre la foule). Ex. : Choisir une filière technique par passion et non par pression parentale.
\end{itemize}

\textbf{Troisième partie : L'anthropologie ancestrale : la mort comme passage fondateur}
\begin{itemize}
    \item \textbf{Image centrale :} « Le tronc d'arbre planté à la naissance ». Métaphore organique : vie et mort sont \textbf{continues} (croissance $\to$ mort de l'arbre = mort de l'homme). Pas de coupure brutale.
    \item \textbf{Statut de l'ancêtre :} « Change de demeure pour veiller ». \cle{Ancestralité} : La mort est un \textbf{changement de modalité d'existence} (corporelle $\to$ spirituelle). L'ancêtre (Ndzém, Miengu) garde l'\cle{autorité} (sanction, bénédiction, cohésion).
    \item \textbf{Fonction sociale :} « Fonde l'autorité des vivants ». Les chefs, les aînés tirent leur légitimité de la fidélité aux ancêtres. La mort \textbf{crée} le lien social (solidarité lignagère, respect des interdits).
\end{itemize}

\textbf{Conclusion}
\textbf{Résumé :} Le texte opère un double renversement : 1. Ontologique (Heidegger) : la mort n'est pas la fin de l'existence, mais sa \textbf{condition de possibilité} (finitude = liberté). 2. Culturel (Ancestralité) : la mort n'est pas le néant, mais le \textbf{fondement de l'ordre symbolique} (autorité, transmission).
\textbf{Portée / Confrontation :} Cette convergence est frappante. Heidegger (phénoménologie occidentale) et la sagesse africaine (ontologie vitale) rejoignent l'idée que \textbf{fuir la mort, c'est fuir sa vie}. Différence : Heidegger reste dans la \textbf{solitude} du Dasein (« personne ne peut mourir à ma place »), alors que la tradition africaine inscrit la mort dans la \textbf{communalité} (l'ancêtre veille sur le « Nous »). Aujourd'hui, la « mort médicalisée » à l'hôpital (Yaoundé, Douala) risque de nous faire régresser vers l'inauthenticité (mort cachée, technique, sans rituel), nous privant de cette « gravité des choix » et de cette « autorité des vivants ».
\end{exoresolu}

\begin{methode}[Appliquer les notions à une situation concrète camerounaise (Cas pratique / Étude de cas)]
\textbf{Objectif :} Mobiliser les concepts (Existence, Mort, Ancestralité, Liberté, Aliénation, Sens) pour analyser un fait de société, un fait divers, une pratique culturelle ou une actualité.
\begin{enumerate}
    \item \textbf{Identifier le fait brut :} Décrire objectivement la situation (Qui ? Quoi ? Où ? Quand ? Contexte social).
    \item \textbf{Repérer les concepts opératoires :} Quels concepts du cours éclairent cette situation ? (ex. : Deuil $\to$ Travail du deuil (Freud) / Rite de passage (Van Gennep) / Ancestralité / Inauthenticité).
    \item \textbf{Problématiser la situation :} Quelle tension philosophique révèle-t-elle ? (ex. : Tradition vs Modernité ; Individu vs Communauté ; Corps vs Esprit ; Vérité vs Croyance).
    \item \textbf{Argumenter en 3 temps (Mini-dissertation) :}
    \begin{enumerate}
        \item \textbf{Analyse conceptuelle :} La situation illustre telle thèse (ex. : Les funérailles coûteuses illustrent l'« existence sociale » primant sur l'individu).
        \item \textbf{Mise en tension / Critique :} Limites, contradictions, évolutions (ex. : Endettement des familles, pression sociale, perte de sens du rite).
        \item \textbf{Synthèse / Évaluation éthique :} Quel sens redonner ? Quelle adaptation possible ? (ex. : Réforme des coutumes, « deuil digne » sans ruine).
    \end{enumerate}
    \item \textbf{Conclusion :} Réponse claire à la question posée, ouverture sur l'avenir de la pratique.
\end{enumerate}
\end{methode}

\begin{exoresolu}[Cas pratique : « Enterrement de luxe » et endettement familial à Douala]
\textbf{Énoncé :} « M. K., cadre bancaire, décède. Sa famille organise des funérailles « dignes de son rang » : cercueil importé (3 millions FCFA), repas pour 500 personnes, location de chaises/tentes, boissons, dot posthume. La famille s'endette lourdement (prêts usuriers, vente d'un terrain non bâti). Les héritiers (femme, enfants) se retrouvent précaires. Un cousin philosophe dit : « Il est mort deux fois : biologiquement, et socialement par la ruine de ses proches. » Analysez cette situation à la lumière des notions d'Existence, de Mort et d'Ancestralité. »

\tcblower

\textbf{Solution détaillée :}

\textbf{1. Identification des concepts opératoires}
\begin{itemize}
    \item \cle{Existence sociale} (Heidegger / Mbiti) : « Je suis parce que nous sommes ». L'individu n'existe pleinement que par sa reconnaissance dans le groupe.
    \item \cle{Mort sociale / Deuxième mort} : La mort biologique n'est que la première étape. La « vraie » mort survient quand plus personne ne se souvient ou quand le rite n'a pas lieu.
    \item \cle{Aliénation / Fétichisme du rite} (Marx / Baudrillard) : Le signe (cercueil, festin) remplace le sens (hommage, transmission). La forme détruit le fond.
    \item \cle{Responsabilité / Justice} (Jonas / Rawls) : Devoir envers les vivants (veuve, orphelins) vs Devoir envers le mort (coutume).
\end{itemize}

\textbf{2. Problématisation}
La tradition ancestrale vise à \textbf{assurer la continuité} (l'ancêtre protège la lignée). Ici, le rite \textbf{détruit l'avenir} de la lignée (endettement, vente du patrimoine). \textbf{Paradoxe :} Le rite censé « fonder l'autorité des vivants » (texte Heidegger/Ancêtres) les plonge dans la précarité.

\textbf{3. Analyse argumentée}

\textbf{I. La logique du don et de l'honneur : L'existence par le regard d'autrui (Thèse traditionnaliste)}
\begin{itemize}
    \item Chez les Sawa, Bamiléké, Bassa, « bien enterrer » = \textbf{payer sa dette} envers le lignage. C'est une \textbf{obligation morale} » (Devoir-être).
    \item Le « cercueil importé » et le festin sont des \textbf{signes de statut} : ils disent au monde social : « Cet homme a compté, sa lignée est puissante ».
    \item \textbf{Validité :} Cela maintient la cohésion, honore la mémoire, rassure le mort sur le sort des siens (croyance). L'existence du défunt est \textbf{validée} par la dépense.
\end{itemize}

\textbf{II. La dérive mortifère : Quand le rite tue la vie (Antithèse critique)}
\begin{itemize}
    \item \textbf{Inversion des finalités :} Le rite (moyen d'intégration) devient une \textbf{fin en soi} (concurrence sociale, « enterrement de l'année »).
    \item \textbf{Violation de l'éthique de la responsabilité (Hans Jonas) :} « Agis pour que les effets de ton action soient compatibles avec la permanence de la vie humaine. » Endetter les orphelins \textbf{contredit} la finalité vitale de l'ancestralité (protéger la descendance).
    \item \textbf{Aliénation marchande :} Le sens (amour, respect) est remplacé par la \textbf{valeur d'échange} (prix du cercueil, nombre de bouteilles). C'est la « société du spectacle » (Debord) au village.
    \item \textbf{La « deuxième mort » :} Le cousin a raison. En ruinant les héritiers, on tue l'\textbf{avenir} que l'ancêtre devait protéger. L'ancêtre devient un \textbf{fardeau} (esprit vengeur ?) au lieu d'un protecteur.
\end{itemize}

\textbf{III. Vers une réinvention éthique du rite (Synthèse) : Le « Deuil digne »}
\begin{itemize}
    \item \textbf{Principe :} Séparer \textbf{l'essentiel symbolique} (veillée, prières, libations, parole des aînés, plantation de l'arbre mémorial) du \textbf{superflu marchand} (cercueil de luxe, festin pléthorique, dettes).
    \item \textbf{Exemples concrets au Cameroun :}
    \begin{itemize}
        \item \emph{Églises de Réveil / Catholiques :} « Enterrement simple », collecte pour la veuve/orphelins au lieu du festin.
        \item \emph{Chefferies modernes (ex. Bafoussam, Foumban) :} Édiction de « Règlements funéraires » plafonnant les dépenses, interdisant les cercueils importés, imposant le « deuil productif » (cotisation pour projet générateur de revenus pour la veuve).
        \item \emph{Assurances obsèques / Mutuelles :} Mutualiser le risque financier pour éviter l'endettement individuel.
    \end{itemize}
    \item \textbf{Philosophie :} L'existence de l'ancêtre ne se mesure pas au \textbf{coût} de sa sortie, mais à la \textbf{qualité de la transmission} (éducation des enfants, paix familiale, patrimoine préservé). C'est ça, « veiller sur la lignée ».
\end{itemize}

\textbf{Conclusion}
La situation révèle une \textbf{crise de sens} » : le rite ancestral, conçu pour \textbf{vaincre la mort} par la mémoire et la solidarité, devient sous l'effet de la modernité marchande et de la pression sociale, un facteur de \textbf{mort sociale} pour les vivants. La philosophie (Heidegger, Jonas, Éthique de la responsabilité) invite à \textbf{réformer la coutume} non pas pour la détruire, mais pour \textbf{sauver son esprit} (la vie, la lignée, la dignité) en changeant sa lettre (la dépense ostentatoire). L'ancêtre véritable ne veut pas la ruine de ses enfants.
\end{exoresolu}

\begin{methode}[Rédiger et justifier une démarche argumentative (Savoir-faire transversal : Dissertation / Commentaire / Épreuve orale)]
\textbf{Objectif :} Produire un raisonnement rigoureux, structuré, étayé, sans affirmations gratuites ni subjectivité non fondée.
\begin{enumerate}
    \item \textbf{La structure standard de l'argument (Toulmin simplifié) :}
    \begin{center}
    \textbf{Revendication (Thèse)} $\leftarrow$ \textbf{Fait / Exemple / Citation (Donnée)} + \textbf{Justification (Garant / Concept)} $\rightarrow$ \textbf{Rebut (Condition de validité / Exception)}
    \end{center}
    \item \textbf{Les 3 types de justification attendus :}
    \begin{itemize}
        \item \textbf{Conceptuelle :} Définition précise, distinction (ex. : « Par \emph{existence}, j'entends non la vie biologique mais... »).
        \item \textbf{Autoritaire (Référence) :} Citation exacte ou référence précise d'auteur au programme (Sartre, Heidegger, Camus, Aristote, Épicure, Auteur africain : Mbiti, Kagame, Mudimbe, Oyono-Mbia). \textit{Attention :} L'autorité ne prouve pas, elle \textbf{appuie}.
        \item \textbf{Factuelle / Concrète :} Exemple réel, historique, littéraire, camerounais (ex. : Code de la famille, Coutume Bamiléké, Actualité sanitaire).
    \end{itemize}
    \item \textbf{Les connecteurs logiques obligatoires :}
    \begin{itemize}
        \item \textbf{Enchaînement :} \emph{En premier lieu, De plus, Par ailleurs, En outre.}
        \item \textbf{Opposition / Nuance :} \emph{Toutefois, Cependant, Certes... mais, Si... cependant.}
        \item \textbf{Conséquence :} \emph{Aussi, Par conséquent, C'est pourquoi, Il s'ensuit que.}
        \item \textbf{Explication :} \emph{En effet, Car, Parce que, Puisque.}
    \end{itemize}
    \item \textbf{La phrase de conclusion de paragraphe (Mini-synthèse) :} Doit répondre à la problématique partielle de la partie. Ex. : « Ainsi, la mort, loin d'être une simple fin, structure le temps de l'existence. »
    \item \textbf{Auto-correction (Relecture) :} Ai-je répondu \textbf{exactement} au sujet ? Y a-t-il des « \emph{je pense} » non argumentés ? Les définitions sont-elles claires ? Les auteurs nommés ?
\end{enumerate}
\end{methode}

\begin{exoresolu}[Exercice d'application : « Rédigez un paragraphe argumenté (15-20 lignes) pour défendre ou critiquer cette affirmation : "Les rites funéraires camerounais sont un frein au développement économique." »]
\textbf{Consigne :} Appliquer la méthode ci-dessus. Position : Nuancée (Dialectique). Structure : Thèse partielle (Coût) $\to$ Antithèse (Valeur sociale/humaine) $\to$ Synthèse (Réforme nécessaire).

\tcblower

\textbf{Paragraphe argumenté type (Niveau Bac/Probatoire) :}

\textbf{Revendication (Ouverture) :} Affirmer que les rites funéraires camerounais constituent un frein au développement économique relève d'une lecture \textbf{réductrice et économiciste} qui confond « développement » et « accumulation de capital financier », tout en ignorant la fonction \textbf{structurelle} de ces rites dans la reproduction du lien social.

\textbf{Argument 1 (Donnée factuelle + Concept : Coût d'opportunité) :} \textbf{En effet}, sur le plan strictement comptable, les dépenses funéraires (cercueils, festins, dots posthumes) peuvent représenter plusieurs années de salaire moyen, plongeant des ménages dans l'endettement (vente d'actifs productifs, prêts usuriers). \textbf{Concept :} C'est une \textbf{consommation improductive} (Keynes) ou une « dépense sans retour » (Bataille) qui réduit l'épargne nationale et l'investissement (éducation, santé, PME). \textbf{Exemple :} À Douala, des familles vendent leur unique parcelle pour « bien enterrer », hypothéquant l'avenir des orphelins.

\textbf{Argument 2 (Nuance / Opposition : Valeur d'usage symbolique) :} \textbf{Toutefois}, analyser ces rites sous le seul angle du \textbf{coût monétaire} (valeur d'échange) occulte leur \textbf{valeur d'usage anthropologique}. Suivant \textbf{Marcel Mauss} (\emph{Essai sur le don}) et \textbf{John Mbiti} (\emph{Religions et philosophies africaines}), le rite funéraire n'est pas une « dépense » mais un \textbf{investissement social} : il répare le tissu communautaire déchiré par la mort, réaffirme l'appartenance lignagère, régule les conflits d'héritage et assure la \textbf{transmission intergénérationnelle} des normes (respect, solidarité). \textbf{Concept :} C'est du \textbf{capital social} (Putnam, Bourdieu) indispensable à la confiance nécessaire à toute activité économique.

\textbf{Argument 3 (Synthèse / Dépassement : La distinction forme/fond) :} \textbf{Par conséquent}, le « frein » ne réside pas dans le rite \emph{per se} (le sens), mais dans son \textbf{hypertrophie marchande} (la forme : luxe importé, gigantisme, mimétisme social). Le développement authentique (Sen : « développement comme liberté ») exige de \textbf{découpler} la charge financière de l'obligation symbolique.

\textbf{Mini-conclusion :} Ainsi, loin d'être un frein absolu, le rite funéraire est un \textbf{levier de cohésion} sans lequel l'individu isolé, livré au marché, devient économiquement plus vulnérable. L'enjeu n'est pas l'abolition, mais la \textbf{régulation coutumière} (plafonnement des frais, mutualisation solidaire) pour que l'hommage aux ancêtres ne devienne pas le sacrifice de la descendance.

\textbf{Analyse de la copie (Pourquoi ça marche) :}
\begin{itemize}
    \item \textbf{Concepts définis :} Développement (Sen), Capital social, Valeur d'usage/échange, Don (Mauss).
    \item \textbf{Auteurs cités :} Mauss, Mbiti, Sen, Keynes/Bataille (références culturelles solides).
    \item \textbf{Exemples camerounais :} Douala, parcelle, orphelins, chefferies.
    \item \textbf{Logique :} Certes (Thèse adverse) $\to$ Cependant (Critique conceptuelle) $\to$ Par conséquent (Synthèse opératoire).
    \item \textbf{Ton :} Objectif, analytique, pas de « Je trouve que », pas d'invective.
\end{itemize}
\end{exoresolu}

\begin{attention}[Les 5 erreurs qui coûtent des points au Bac (Probatoire / Baccalauréat Technique)]
\begin{enumerate}
    \item \textbf{Confondre « Vie » et « Existence » (ou « Mort biologique » et « Mort ontologique »).} \textit{Sanction :} Hors-sujet conceptuel (0/4 ou 1/4 en analyse). \textbf{Réflexe :} Définir \emph{avant} d'argumenter. Vie = biologique / Existence = projet, sens, liberté / Mort bio = fin corps / Mort onto = horizon de finitude.
    \item \textbf{Faire un « catalogue d'auteurs » sans problématiser.} \textit{Exemple :} « Platon dit ça, Aristote dit ça, Sartre dit ça. » \textbf{Sanction :} Note de « connaissance » seulement (max 6-8/20), pas de « dissertation ». \textbf{Réflexe :} Les auteurs sont des \textbf{témoins} ou des \textbf{adversaires} dans \emph{votre} démonstration. Les confronter : « \emph{Contre} Épicure, Heidegger montre que... »
    \item \textbf{Négliger l'ancrage camerounais / Exemples concrets.} \textit{Consigne programme :} « Appliquer les notions à des situations concrètes de la vie courante ». \textbf{Sanction :} Pénalité sur la note de « savoir-faire » / « application » (souvent 4 à 6 pts sur 20). \textbf{Réflexe :} Un exemple local par grande partie (Deuil Bamiléké, Hôpital Central, Start-up Jeune, Chefferie, Code de la Famille, Actualité locale).
    \item \textbf{Problématique absente, floue ou hors-sujet.} \textit{Exemple :} Sujet « La mort donne-t-elle un sens à la vie ? » $\to$ Problématique : « Qu'est-ce que la mort ? » \textbf{Sanction :} Plan incohérent, note plafond 8/20. \textbf{Réflexe :} Problématique = \textbf{Question tension} issue des définitions. Elle doit contenir les mots-clés du sujet.
    \item \textbf{Conclusion « ouverture » bâclée ou inexistante / Phrase de conclusion de paragraphe manquante.} \textit{Erreur :} Finir par « En conclusion, la mort est importante. » ou s'arrêter au développement. \textbf{Sanction :} Perte de 2 à 3 points (structure). \textbf{Réflexe :} Conclusion = \textbf{Bilan synthétique} (Réponse directe à la problématique) + \textbf{Ouverture pertinente} (Autre champ : Bioéthique, IA, Écologie, Art, Politique). Chaque grande partie (I, II, III) doit se terminer par une phrase résumant l'apport de la partie à la réponse finale.
\end{enumerate}
\end{attention}

\newpage

\coursec{Exercices}

\courssub{Je vérifie mes connaissances}

\begin{exercice}[QCM : Concepts fondamentaux]
Pour chacune des affirmations suivantes, une seule proposition est exacte. Recopiez le numéro et la lettre correspondante.
\begin{enumerate}
\item La distinction philosophique entre \emph{vie} et \emph{existence} repose principalement sur :
\begin{itemize}
\item[A.] La vie est biologique, l'existence est la vie humaine vécue comme conscience et projet.
\item[B.] La vie concerne les animaux, l'existence concerne seulement les philosophes.
\item[C.] Il n'y a pas de différence, ce sont des synonymes.
\item[D.] L'existence est la durée, la vie est l'intensité.
\end{itemize}
\item Pour Épicure, la mort est :
\begin{itemize}
\item[A.] Le plus grand des maux car elle nous prive des plaisirs futurs.
\item[B.] « Rien pour nous » car « tant que nous existons, la mort n'est pas là, et quand la mort est là, nous ne sommes plus ».
\item[C.] Le passage vers une autre vie meilleure.
\item[D.] Une épreuve initiatoire nécessaire.
\end{itemize}
\item Selon Heidegger (\emph{Sein und Zeit}), la mort est :
\begin{itemize}
\item[A.] Un événement extérieur qui nous arrive comme une maladie.
\item[B.] Une possibilité d'être propre, certaine, indépassable, qui individualise le Dasein dans son être-pour-la-mort.
\item[C.] La fin de la conscience, donc la fin de tout sens.
\item[D.] Un concept purement biologique sans portée ontologique.
\end{itemize}
\item Dans la pensée traditionnelle camerounaise (conception bantoue), la mort est souvent conçue comme :
\begin{itemize}
\item[A.] Une annihilation totale de la personne.
\item[B.] Un passage, un changement d'état vers le monde des ancêtres, où le défunt continue d'agir sur les vivants.
\item[C.] Une injustice réparable par la magie.
\item[D.] La fin définitive des liens sociaux.
\end{itemize}
\end{enumerate}
\end{exercice}

\begin{exercice}[Vrai / Faux justifié]
Indiquez si les propositions sont vraies ou fausses et justifiez brièvement votre réponse en mobilisant une référence philosophique ou culturelle.
\begin{enumerate}
\item « L'existence précède l'essence » (Sartre) signifie que l'homme naît avec une nature prédéfinie qu'il doit réaliser.
\item Pour Jankélévitch, la mort est une « aventure » qui nous concerne au premier chef mais que nous ne pouvons pas faire l'expérience de notre propre mort.
\item Le \emph{Pour-soi} sartrien est une conscience pleine, solide, identique à elle-même, à l'image de l'\emph{En-soi}.
\item Chez les Bamiléké comme chez les Bassa, les funérailles sont un moment purement privé, sans dimension sociale ni politique.
\end{enumerate}
\end{exercice}

\begin{exercice}[Définitions et distinctions]
Définissez précisément les concepts suivants en philosophie de l'existence (2 à 3 lignes chacun) :
\begin{enumerate}
\item L'existence (vs vie biologique).
\item La facticité (Heidegger / Sartre).
\item L'angoisse (vs la peur).
\item L'ancestralité (conception africaine classique).
\end{enumerate}
\end{exercice}

\begin{exercice}[Association Auteurs / Thèses]
Reliez chaque auteur à la thèse qui lui correspond le mieux concernant le sens de l'existence face à la mort.
\begin{center}
\begin{tabularx}{\linewidth}{|l|X|}\hline
\textbf{Auteur} & \textbf{Thèse} \\ \hline
Épicure & 1. La mort donne à la vie son urgence et son authenticité ; assumer sa mortalité fonde la liberté. \\ \hline
Martin Heidegger & 2. L'existence est absurde car l'homme cherche du sens dans un monde silencieux ; il faut imaginer Sisyphe heureux. \\ \hline
Jean-Paul Sartre & 3. La mort ne nous concerne pas ; le sage ne la craint pas car elle est privation de sensation. \\ \hline
Albert Camus & 4. L'homme est condamné à être libre ; la mort est le triomphe du point de vue d'autrui qui me fige en être-pour-autrui. \\ \hline
Vladimir Jankélévitch & 5. La mort est l'irréversible par excellence, l'intouchable qui rend la vie précieuse et chaque instant unique. \\ \hline
\end{tabularx}
\end{center}
\end{exercice}

\courssub{Je m'entraîne}

\begin{exercice}[Analyse de citations]
Pour chacune des citations ci-dessous : identifiez l'auteur, expliquez le sens des concepts soulignés, et précisez la thèse défendue sur le rapport existence/mort.
\begin{enumerate}
\item « L'homme n'est pas seulement un être qui meurt, il est un être \textbf{mortel} ; c'est-à-dire qu'il a la mort pour possibilité propre. » (Heidegger, \emph{Sein und Zeit})
\item « La mort est le \textbf{triomphe du point de vue d'autrui}. » (Sartre, \emph{L'Être et le Néant})
\item « Il n'y a qu'un problème philosophique vraiment sérieux : c'est le \textbf{suicide}. Juger si la vie vaut ou non la peine d'être vée, c'est répondre à la question fondamentale de la philosophie. » (Camus, \emph{Le Mythe de Sisyphe})
\item « La mort est l'\textbf{irréparable}, l'\textbf{irréversible} ; elle est ce qui ne peut pas être rattrapé, ce qui ne peut pas être refait. » (Jankélévitch, \emph{La Mort})
\end{enumerate}
\end{exercice}

\begin{exercice}[Situation concrète : Le deuil au Cameroun]
M. Nguema, chef de famille bassa, vient de perdre son père. La coutume exige une veillée de plusieurs nuits, des libations, des danses de deuil, et un second enterrement (le « deuil officiel ») un an plus tard pour « envoyer » le défunt auprès des ancêtres. M. Nguema, cadre bancaire moderne, trouve ces rites coûteux, longs et incompatibles avec son travail. Il souhaite une cérémonie brève au cimetière.
\begin{enumerate}
\item Identifiez le conflit de valeurs entre la conception \emph{traditionnelle} de la mort (ancestralité, lien social) et la conception \emph{moderne/occidentale} (efficacité, vie privée, finitude biologique).
\item En vous appuyant sur la distinction \emph{Vie / Existence}, expliquez pourquoi les rites funéraires traditionnels visent à assurer la \emph{survie de l'existence} du défunt au-delà de la fin de sa \emph{vie biologique}.
\item Proposez un argumentaire philosophique (3 arguments) pour concilier devoir professionnel et devoir filial/ancestral dans ce cas précis.
\end{enumerate}
\end{exercice}

\begin{exercice}[Comparaison de thèses : Le sens de l'existence face à la mort]
Complétez le tableau comparatif suivant en synthétisant la position de chaque auteur.
\begin{center}
\begin{tabularx}{\linewidth}{|l|X|X|X|}\hline
\textbf{Auteur} & \textbf{Nature de la mort} & \textbf{Attitude face à la mort} & \textbf{Conséquence sur le sens de l'existence} \\ \hline
Épicure & & & \\ \hline
Heidegger & & & \\ \hline
Sartre & & & \\ \hline
Camus & & & \\ \hline
Jankélévitch & & & \\ \hline
\end{tabularx}
\end{center}
\end{exercice}

\begin{exercice}[Argumentation dirigée : La mort donne-t-elle du sens à la vie ?]
Rédigez un paragraphe argumenté (150-200 mots) répondant à la question : « La conscience de la mort est-elle ce qui humanise l'existence ? ».
\vspace{0.5cm}
\noindent \textbf{Contraintes :}
\begin{itemize}
\item Intégrez la définition de l'existence comme projet/liberté.
\item Opposez la thèse de Heidegger (mort comme possibilité propre qui individualise) à celle d'Épicure (mort comme néant indifférent).
\item Concluez par une nuance personnelle (ex: Jankélévitch, Camus, ou perspective africaine).
\end{itemize}
\end{exercice}

\courssub{Je me prépare au Bac}

\begin{exercice}[Sujet de Dissertation — Type Bac Tech]
\textbf{Sujet :} \emph{« La conscience de la mort est-elle ce qui humanise l'existence ? »}

\vspace{0.3cm}
\noindent \textbf{Consignes :}
\begin{enumerate}
\item Analysez les termes du sujet : définissez \emph{conscience}, \emph{mort}, \emph{humaniser}, \emph{existence}. Distinguez existence (humaine) et vie (biologique).
\item Problématisez : en quoi la mort (finitude) menace-t-elle le sens de l'existence, et paradoxalement, en quoi sa conscience pourrait-elle la fonder ?
\item Construisez un plan dialectique en trois parties (Thèse / Antithèse / Synthèse) en intégrant obligatoirement :
\begin{itemize}
\item \textbf{Thèse :} La conscience de la mort comme condition de l'authenticité et de la liberté (Heidegger : être-pour-la-mort ; Jankélévitch : prix de la vie).
\item \textbf{Antithèse :} La conscience de la mort comme source d'angoisse, d'absurde ou d'aliénation (Épicure : ne pas s'en soucier ; Sartre : mort comme triomphe d'autrui, contingence absurde ; Camus : l'absurde).
\item \textbf{Synthèse :} Dépasser l'opposition : la mort comme horizon qui oblige à l'invention de sens (projet, transmission, ancestralité). Exemple camerounais : la mort comme passage et responsabilité intergénérationnelle.
\end{itemize}
\item Rédigez l'introduction complète (accroche, définitions, problématique, annonce du plan) et la conclusion (réponse synthétique, ouverture).
\end{enumerate}
\end{exercice}

\begin{exercice}[Explication de Texte — Type Bac Tech]
\textbf{Texte :} Extrait de \emph{L'Être et le Néant}, Jean-Paul Sartre (1943), partie IV, chapitre 1, « La Mort ».

\vspace{0.2cm}
\noindent « La mort n'est pas une menace ; elle est la fin de tous les possibles. [...] La mort est le triomphe du point de vue d'autrui. C'est par la mort que l'autre me dispose. [...] Je suis mort pour autrui. Mais je ne puis pas assumer ma mort ; elle m'échappe. Elle est cette existence qui m'échappe et qui m'appartient. [...] L'être-pour-la-mort n'est pas un être-vers-la-fin ; c'est un être-vers-le-commencement de ma fin. »

\vspace{0.3cm}
\noindent \textbf{Questions :}
\begin{enumerate}
\item \textbf{Thèse :} Dégagez la thèse principale du texte en une phrase.
\item \textbf{Concepts :} Expliquez les concepts suivants \emph{dans le texte} : « triomphe du point de vue d'autrui », « l'autre me dispose », « être-pour-la-mort » (différence avec Heidegger).
\item \textbf{Argumentation :} Relevez et reformulez les deux arguments principaux qui soutiennent la thèse.
\item \textbf{Discussion critique :} Ce texte affirme que je ne peux pas assumer ma mort. Discutez cette affirmation en la confrontant à la position de Heidegger (l'anticipation résolue) et à la sagesse épicurienne. La mort peut-elle être « mon » acte ?
\end{enumerate}
\end{exercice}

\begin{exercice}[Cas Pratique / Invention — Type Bac Tech]
\textbf{Situation :}
Vous êtes le Dr Fouda, médecin chef du service de soins palliatifs de l'Hôpital Central de Yaoundé. M. K., 68 ans, atteint d'un cancer du pancréas au stade terminal, souffre de douleurs réfractaires aux antalgiques de palier 3. Il est lucide, compétent, et vous demande formellement, par écrit, une « aide active à mourir » (injection létale) pour abréger ses souffrances. Sa femme et ses enfants adultes s'y opposent farouchement : « Chez nous (peuple Bamiléké), la vie appartient à la lignée et aux ancêtres ; seul Dieu/les ancêtres décident du rappel. Demander la mort est une offense aux ancêtres et un suicide, c'est tabou. Nous voulons qu'il vive ses derniers jours parmi nous, même dans la douleur, c'est son devoir. »

\vspace{0.3cm}
\noindent \textbf{Travail demandé :}
Rédigez le \textbf{raisonnement éthique structuré} (argumentaire) que vous présentez au Comité d'Éthique Hospitalier pour justifier votre décision finale (accorder ou refuser la demande). Votre texte doit comporter :
\begin{enumerate}
\item L'identification claire du \textbf{conflit de valeurs} (Autonomie du patient / Bienfaisance / Non-malfaisance vs Tradition communautaire / Sainteté de la vie / Volonté des ancêtres).
\item L'analyse du \textbf{cadre légal camerounais} (Code de déontologie médicale, Code pénal — euthanasie active interdite, sédation profonde et continue autorisée ?).
\item L'argumentation philosophique mobilisant :
\begin{itemize}
\item La notion d'\emph{existence} comme projet et autonomie (Sartre, Kant : dignité, fin en soi).
\item La conception \emph{communautaire/ancestrale} de la personne (Mbiti, Wiredu : « Je suis parce que nous sommes », corps social).
\item La distinction \emph{laisser mourir / faire mourir} et la doctrine de l'effet double (sédation palliative).
\end{itemize}
\item Une \textbf{décision motivée} (proposition concrète : ex. refus de l'euthanasie active, mise en place sédation palliative profonde, accompagnement psychologique/spirituel intégrant la famille, médiation culturelle).
\item Une \textbf{ouverture} sur la nécessaire évolution du dialogue entre bioéthique universelle et éthiques africaines.
\end{enumerate}
\end{exercice}

\newpage

\coursec{Corrigés des exercices}

\begin{corrige}[Exercice 1 — QCM : Concepts fondamentaux]
\begin{enumerate}
\item \textbf{Réponse A.} La \cle{vie} désigne le phénomène biologique commun à tous les vivants (métabolisme, reproduction, mort cellulaire) ; l'\cle{existence} désigne le mode d'être spécifiquement humain : une vie vécue en conscience, capable de se rapporter à elle-même, de se projeter et de se choisir (Heidegger : le \emph{Dasein} est cet être qui a à être son être). B est faux (l'existence concerne tout homme), C est faux (la distinction est fondatrice de la philosophie de l'existence), D est une confusion sans fondement philosophique.
\item \textbf{Réponse B.} Pour Épicure (\emph{Lettre à Ménécée}), la mort est privation de sensation ; or le mal suppose un sujet qui éprouve. Donc « la mort n'est rien pour nous » : vivant, je ne la rencontre pas ; mort, je ne suis plus là pour la subir. A est la thèse inverse, C relève du platonisme ou des religions, D de l'orphisme.
\item \textbf{Réponse B.} Dans \emph{Sein und Zeit} (1927), la mort est la possibilité « la plus propre, inconditionnée, indépassable » du \emph{Dasein}. L'anticipation de cette possibilité (l'\emph{être-pour-la-mort}) arrache l'homme à l'anonymat du « On » et l'individualise. A, C et D réduisent la mort à un événement biologique extérieur, ce que Heidegger refuse précisément.
\item \textbf{Réponse B.} Dans les conceptions bantoues (Mbiti : « Je suis parce que nous sommes »), la mort est un \emph{passage} : le défunt rejoint les ancêtres, demeure membre de la communauté et continue d'agir sur les vivants (bénédictions, protections, sanctions). A et D contredisent la continuité du lien social ; C est une caricature.
\end{enumerate}
\textbf{Points de vigilance :} exiger la justification de chaque choix par une référence précise (œuvre, concept) ; ne pas accepter la seule lettre.
\end{corrige}

\begin{corrige}[Exercice 2 — Vrai / Faux justifié]
\begin{enumerate}
\item \textbf{Faux.} « L'existence précède l'essence » (Sartre, \emph{L'Existentialisme est un humanisme}, 1946) signifie exactement l'inverse : l'homme existe d'abord, se jette dans le monde, et ne se définit qu'ensuite par ses actes. Il n'y a pas de nature humaine prédéfinie (pas de « fabricant » divin, pas d'essence à réaliser) : « l'homme n'est rien d'autre que ce qu'il se fait ».
\item \textbf{Vrai.} Jankélévitch (\emph{La Mort}, 1966) distingue la mort des autres (la « deuxième mort », dont je peux faire l'expérience du deuil) et ma propre mort (la « troisième mort »), que je ne peux jamais vivre puisque je cesse d'être au moment où elle advient. Ma mort m'est à la fois la chose la plus mienne et la plus insaisissable : une « aventure » sans récit possible.
\item \textbf{Faux.} Chez Sartre (\emph{L'Être et le Néant}), le \cle{Pour-soi} (la conscience) est néantisation, manque, perpétuel écart avec soi : il est « ce qu'il n'est pas et n'est pas ce qu'il est ». C'est l'\emph{En-soi} (la chose, l'être plein) qui est massif, coïncidant avec soi. La proposition inverse les deux concepts.
\item \textbf{Faux.} Les funérailles bamiléké (comme bassa) sont un événement social et politique majeur : elles rassemblent la lignée, règlent la succession, réaffirment les alliances et la hiérarchie, et assurent le passage du défunt au statut d'ancêtre. Le deuil est une affaire de communauté, non de sphère privée.
\end{enumerate}
\textbf{Points de vigilance :} la justification doit citer l'auteur et l'œuvre ; une réponse « vrai/faux » sans justification ne vaut aucun point.
\end{corrige}

\begin{corrige}[Exercice 3 — Définitions et distinctions]
\begin{enumerate}
\item \textbf{L'existence (vs vie biologique).} La vie est le processus organique (naissance, croissance, mort) partagé par tous les vivants. L'existence est le mode d'être proprement humain : une vie qui se sait vivante, se rapporte à elle-même, se projette vers des possibles et se définit par ses choix. Exister, c'est « avoir à être » son être (Heidegger), c'est-à-dire être en question pour soi-même.
\item \textbf{La facticité.} Ensemble des données que l'existence n'a pas choisies et qu'elle doit assumer : mon corps, ma naissance, mon époque, ma langue, ma famille, ma mort. Heidegger parle de « déréliction » (\emph{Geworfenheit}, être-jeté) ; Sartre montre que la liberté n'existe qu'en situation : je ne choisis pas ma facticité, mais je choisis le sens que je lui donne.
\item \textbf{L'angoisse (vs la peur).} La peur a un objet déterminé et extérieur (un lion, un examen). L'angoisse n'a pas d'objet : elle est le sentiment de l'être lui-même — de ma liberté vertigineuse (Sartre : « l'angoisse, c'est la conscience de ma liberté ») ou de ma finitude face au néant (Heidegger : l'angoisse individualise et révèle l'être-pour-la-mort).
\item \textbf{L'ancestralité.} Dans les traditions africaines (Mbiti, conception bantoue), rapport continu entre vivants et morts : les ancêtres, morts « vivants », demeurent membres de la communauté, protègent et sanctionnent les vivants, qui leur doivent rites et mémoire. La mort est passage, non annihilation ; la personne se définit par son insertion dans la chaîne des générations.
\end{enumerate}
\textbf{Points de vigilance :} chaque définition doit comporter la distinction demandée (existence/vie, angoisse/peur) et une référence ; éviter les définitions par simple synonyme.
\end{corrige}

\begin{corrige}[Exercice 4 — Association Auteurs / Thèses]
\begin{center}
\begin{tabularx}{\linewidth}{|l|X|}\hline
\textbf{Auteur} & \textbf{Thèse correspondante} \\ \hline
Épicure & \textbf{3.} La mort ne nous concerne pas ; le sage ne la craint pas car elle est privation de sensation (\emph{Lettre à Ménécée}). \\ \hline
Martin Heidegger & \textbf{1.} La mort donne à la vie son urgence et son authenticité ; l'être-pour-la-mort individualise et fonde la liberté (\emph{Sein und Zeit}). \\ \hline
Jean-Paul Sartre & \textbf{4.} L'homme est condamné à être libre ; la mort me fige en être-pour-autrui, elle est « le triomphe du point de vue d'autrui » (\emph{L'Être et le Néant}). \\ \hline
Albert Camus & \textbf{2.} L'existence est absurde (confrontation de l'appel humain et du silence du monde) ; « il faut imaginer Sisyphe heureux » (\emph{Le Mythe de Sisyphe}). \\ \hline
Vladimir Jankélévitch & \textbf{5.} La mort est l'irréversible par excellence ; c'est elle qui rend chaque instant unique et la vie précieuse (\emph{La Mort}). \\ \hline
\end{tabularx}
\end{center}
\textbf{Points de vigilance :} attention au piège classique Heidegger/Sartre : chez Heidegger la mort est ma possibilité la plus propre (je peux l'anticiper) ; chez Sartre elle m'échappe et appartient à autrui. Confondre les deux est l'erreur la plus fréquente.
\end{corrige}

\begin{corrige}[Exercice 5 — Analyse de citations]
\begin{enumerate}
\item \textbf{Heidegger, \emph{Sein und Zeit}.} \emph{Mourir} est le sort commun des vivants (événement biologique) ; être \emph{mortel} est une structure ontologique : l'homme seul se sait mortel et rapporte sa mort à soi comme \emph{possibilité propre}. Thèse : la mort n'est pas un accident qui survient, mais une possibilité constitutive de l'existence ; l'anticiper (être-pour-la-mort) arrache le \emph{Dasein} au « On » anonyme et le rend authentique et libre.
\item \textbf{Sartre, \emph{L'Être et le Néant}.} Le \emph{triomphe du point de vue d'autrui} signifie qu'à ma mort, mon existence devient un objet entre les mains des autres : ils fixent le sens de ma vie (souvenirs, jugements, biographies) sans que je puisse le modifier. Thèse : la mort n'est pas ma possibilité (contre Heidegger) mais une contingence absurde qui me dépossède de mon sens ; elle ne donne pas de sens à la vie, elle le retire à mon contrôle.
\item \textbf{Camus, \emph{Le Mythe de Sisyphe}.} Le \emph{suicide} est la question fondamentale car il pose le problème de la valeur de la vie face à l'absurde : le divorce entre mon désir de sens et le silence du monde. Thèse : il ne faut ni suicide (fuite) ni espoir religieux (suicide philosophique), mais la révolte lucide — vivre sans appel, en gardant l'absurde éveillé : « il faut imaginer Sisyphe heureux ».
\item \textbf{Jankélévitch, \emph{La Mort}.} L'\emph{irréparable} (ce qui ne peut être rattrapé) et l'\emph{irréversible} (ce qui ne peut être refait) définissent la mort comme l'intouchable absolu. Thèse : c'est précisément cette irréversibilité qui donne son prix à chaque instant : parce que tout est unique et non recommençable, la vie est précieuse et engageante. La mort est le « trésor » qui fonde la valeur de l'existence.
\end{enumerate}
\textbf{Points de vigilance :} trois étapes exigées par citation : auteur + œuvre, explication des concepts soulignés, thèse sur le rapport existence/mort. Ne pas paraphraser la citation : il faut dégager la thèse.
\end{corrige}

\begin{corrige}[Exercice 6 — Situation concrète : Le deuil au Cameroun]
\begin{enumerate}
\item \textbf{Conflit de valeurs.} D'un côté, la conception traditionnelle : la mort est un passage vers l'ancestralité ; le défunt reste membre actif de la communauté ; les rites (veillées, libations, second deuil) sont des devoirs qui maintiennent le lien entre vivants et ancêtres et la cohésion du groupe. De l'autre, la conception moderne : la mort est un événement biologique privé ; prime l'efficacité, le temps de travail, l'autonomie individuelle ; les rites longs et coûteux apparaissent comme des contraintes irrationnelles. Le conflit oppose donc \emph{devoir envers la communauté et les ancêtres} à \emph{devoir professionnel et rationalité économique}.
\item \textbf{Vie / Existence.} La vie biologique du père de M. Nguema est finie : l'organisme a cessé de fonctionner. Mais son \emph{existence} — ce qu'il a été, son nom, sa parole, sa place dans la lignée — peut survivre dans la mémoire et l'action de la communauté. Les rites funéraires sont précisément la technique culturelle de cette survie existentielle : le second deuil « envoie » le défunt chez les ancêtres, c'est-à-dire transforme la fin biologique en nouveau statut social et spirituel. Supprimer les rites, c'est réduire le défunt à sa seule vie biologique et consommer sa « seconde mort » (l'oubli).
\item \textbf{Argumentaire de conciliation.}
\begin{itemize}
\item \emph{Argument 1 (hiérarchie des devoirs)} : le devoir filial et ancestral est constitutif de l'identité de M. Nguema (« je suis parce que nous sommes ») ; y renoncer totalement, c'est se couper de ce qui le fonde. Le devoir professionnel, lui, peut être aménagé (congés, délégation) sans être détruit.
\item \emph{Argument 2 (Sartre : liberté en situation)} : la facticité (sa culture, son travail) n'abolit pas sa liberté ; M. Nguema peut choisir librement une forme réduite mais authentique des rites (présence aux moments essentiels, participation financière raisonnable), assumant ainsi sa situation sans mauvaise foi.
\item \emph{Argument 3 (sens de l'existence)} : accomplir les rites, c'est aussi préparer son propre statut futur d'ancêtre et transmettre à ses enfants le sens de la continuité ; l'investissement temporel n'est pas une perte mais une inscription dans la durée intergénérationnelle.
\end{itemize}
\end{enumerate}
\textbf{Points de vigilance :} appliquer explicitement les notions du cours (vie/existence, facticité, ancestralité) à la situation ; éviter le jugement ethnocentrique (« les traditions sont dépassées ») comme le traditionalisme fermé.
\end{corrige}

\begin{corrige}[Exercice 7 — Comparaison de thèses]
\begin{center}
\begin{tabularx}{\linewidth}{|l|X|X|X|}\hline
\textbf{Auteur} & \textbf{Nature de la mort} & \textbf{Attitude face à la mort} & \textbf{Conséquence sur le sens de l'existence} \\ \hline
Épicure & Privation de sensation ; « rien pour nous » & Ne pas la craindre ; sérénité (ataraxie) du sage & Vivre heureux ici et maintenant, sans crainte ni espoir d'au-delà \\ \hline
Heidegger & Possibilité la plus propre, certaine, indépassable & Anticipation résolue (être-pour-la-mort) & La mort individualise, arrache au « On », fonde l'authenticité et la liberté \\ \hline
Sartre & Contingence absurde ; triomphe du point de vue d'autrui & On ne peut l'assumer ; elle m'échappe & La mort ne donne pas de sens ; c'est ma liberté seule qui fait sens de mon vivant \\ \hline
Camus & Révélatrice de l'absurde (monde sans réponse) & Révolte lucide : ni suicide ni espoir & Vivre sans appel, multiplier les expériences ; « Sisyphe heureux » \\ \hline
Jankélévitch & L'irréversible, l'irréparable, l'intouchable & Lucidité grave ; ne pas la banaliser & La mort rend chaque instant unique : elle est le prix de la vie \\ \hline
\end{tabularx}
\end{center}
\textbf{Points de vigilance :} chaque case doit contenir une formule synthétique et exacte ; vérifier la cohérence interne de chaque ligne (la nature de la mort doit commander l'attitude et la conséquence). Barème indicatif : 1 point par ligne complète et cohérente.
\end{corrige}

\begin{corrige}[Exercice 8 — Argumentation dirigée : La mort donne-t-elle du sens à la vie ?]
\textbf{Paragraphe corrigé (modèle).}

L'existence humaine se distingue de la vie biologique en ce qu'elle est conscience de soi, liberté et projet : l'homme n'est pas, il a à être, il se fait en se choisissant. Or c'est précisément la conscience de la mort qui donne à ce projet son urgence. Pour Heidegger, la mort est ma possibilité la plus propre, certaine et indépassable : anticiper ma finitude (l'être-pour-la-mort) m'arrache à l'anonymat du « On », m'individualise et me rend à ma liberté authentique — je cesse de vivre « comme tout le monde » pour assumer une vie qui est mienne et non recommençable. Épicure objecterait que la mort n'est « rien pour nous » : tant que je suis, elle n'est pas ; quand elle est, je ne suis plus. S'en soucier serait donc une crainte vaine, et la sérénité du sage consisterait à l'ignorer. Mais cette indifférence, si elle délivre de la peur, risque aussi de dissoudre l'urgence de vivre : si la mort ne me concerne jamais, rien ne m'oblige à hiérarchiser mes possibles. Avec Jankélévitch, on dira plutôt que c'est l'irréversibilité de la mort qui fait le prix de chaque instant ; et les traditions camerounaises confirment que penser sa mort, c'est penser sa place dans la chaîne des générations. La conscience de la mort humanise donc l'existence non en la terrifiant, mais en la rendant unique, limitée et responsable.

\textbf{Barème indicatif (sur 5) :} définition existence/projet (1 pt) ; thèse de Heidegger exacte (1 pt) ; objection épicurienne exacte (1 pt) ; nuance personnelle argumentée (1 pt) ; cohérence, langue, respect du volume (1 pt).

\textbf{Points de vigilance :} respecter les 150-200 mots ; mobiliser les deux auteurs imposés ; la nuance finale ne doit pas contredire mais dépasser l'opposition ; éviter la paraphrase des cours sans articulation argumentative.
\end{corrige}

\begin{corrige}[Exercice 9 — Sujet de Dissertation : « La conscience de la mort est-elle ce qui humanise l'existence ? »]
\textbf{1. Analyse des termes.}
\begin{itemize}
\item \emph{Conscience} : rapport de l'esprit à soi et au monde ; ici, savoir lucide de sa finitude (non simple information, mais prise en charge existentielle).
\item \emph{Mort} : fin de la vie biologique ; mais aussi possibilité propre (Heidegger), néant indifférent (Épicure), passage vers l'ancestralité (traditions africaines).
\item \emph{Humaniser} : rendre proprement humain, distinguer l'homme de l'animal ; ce qui fait de la vie une existence.
\item \emph{Existence} : mode d'être humain comme conscience, liberté, projet — à distinguer de la vie biologique.
\end{itemize}
\textbf{2. Problématique.} La mort semble d'abord nier tout sens : si tout finit, à quoi bon ? (absurde, angoisse). Mais paradoxalement, un vivant immortel n'aurait ni urgence, ni choix décisifs, ni valeur : tout serait indifféremment recommençable. La conscience de la mort ne serait-elle donc pas la condition même du sérieux, de la liberté et de la valeur de l'existence ? Faut-il, pour être pleinement homme, penser sa mort — ou au contraire l'écarter pour vivre ?
\textbf{3. Plan dialectique.}
\begin{itemize}
\item \textbf{I. Thèse} : la conscience de la mort fonde l'authenticité. Heidegger : l'être-pour-la-mort individualise et libère du « On » ; Jankélévitch : l'irréversible donne son prix à chaque instant ; sans finitude, pas de choix véritables.
\item \textbf{II. Antithèse} : la conscience de la mort aliène ou angoisse. Épicure : elle est une crainte vaine à éliminer ; Sartre : la mort est contingence absurde, triomphe d'autrui, elle ne se laisse pas assumer ; Camus : elle révèle l'absurde qui menace tout sens.
\item \textbf{III. Synthèse} : la mort n'est ni à fuir ni à subir : elle est l'horizon qui oblige à inventer le sens — projet personnel (Sartre : la liberté fait sens de son vivant), révolte lucide (Camus), transmission et ancestralité (la mort comme passage et responsabilité intergénérationnelle dans les cultures camerounaises : rites, mémoire, devoir envers les descendants).
\end{itemize}
\textbf{4. Introduction (modèle).} Tout homme sait qu'il mourra ; aucun animal ne le sait. Ce savoir, qui distingue l'existence humaine de la simple vie biologique, est aussi un scandale : comment vivre et agir sous la menace d'une fin qui anéantira tout ? La \emph{mort} est la cessation de la vie organique ; la \emph{conscience} de la mort est le rapport lucide de l'homme à sa finitude ; \emph{humaniser} signifie constituer l'existence comme existence proprement humaine, c'est-à-dire comme liberté et projet. Dès lors, la question s'impose : la conscience de la mort est-elle ce qui humanise l'existence ? En quoi la finitude, qui semble détruire tout sens, pourrait-elle paradoxalement le fonder ? Nous montrerons d'abord que la conscience de la mort conditionne l'authenticité et la liberté ; nous verrons ensuite qu'elle peut aussi être source d'angoisse et d'absurde ; nous dépasserons enfin cette opposition en montrant que la mort est l'horizon qui oblige l'homme à inventer et transmettre le sens.
\textbf{Conclusion (modèle).} La conscience de la mort humanise l'existence, non parce qu'elle la terrifie, mais parce qu'elle la limite : c'est la finitude qui rend chaque choix décisif, chaque instant unique, chaque vie irremplaçable. Encore faut-il transformer ce savoir en tâche : vivre authentiquement, créer du sens, transmettre. Reste ouverte une question que nos sociétés camerounaises posent avec acuité : le dialogue entre la finitude individuelle et la survie dans la mémoire des ancêtres ne dessine-t-il pas une autre manière, proprement africaine, de vaincre la mort ?

\textbf{Barème indicatif (sur 20) :} introduction (accroche, définitions, problématique, plan) : 4 pts ; thèse argumentée avec références : 4 pts ; antithèse argumentée avec références : 4 pts ; synthèse personnelle et dépassement : 4 pts ; conclusion avec réponse et ouverture : 2 pts ; langue, organisation, exemples camerounais : 2 pts.

\textbf{Points de vigilance :} ne pas transformer la dissertation en catalogue d'auteurs ; chaque référence doit servir une étape de l'argumentation ; la synthèse doit être un dépassement (pas un compromis mou) ; l'exemple camerounais (ancestralité) est exigé par le sujet de la leçon.
\end{corrige}

\begin{corrige}[Exercice 10 — Explication de Texte : Sartre, \emph{L'Être et le Néant}]
\textbf{1. Thèse.} La mort n'est pas une possibilité que je puisse assumer comme mienne (contre Heidegger) : elle est une contingence absurde qui m'échappe et livre le sens de ma vie au regard d'autrui.
\textbf{2. Concepts.}
\begin{itemize}
\item \emph{Triomphe du point de vue d'autrui} : une fois mort, je deviens un objet pour les autres ; ce sont eux qui fixent désormais le sens de mon existence (mémoire, jugements, récits) sans que je puisse rien y changer. Ma vie entière devient un « être-pour-autrui » figé.
\item \emph{L'autre me dispose} : les survivants disposent de moi comme d'une chose — de mon corps, de mes biens, de ma réputation, du sens de mes actes. Je perds toute maîtrise sur ce que je suis.
\item \emph{Être-pour-la-mort} : chez Sartre, ce n'est pas (comme chez Heidegger) une anticipation qui me rend authentique, mais un « être-vers-le-commencement de ma fin » : je suis à chaque instant exposé à une mort toujours possible, imprévisible, qui peut interrompre mes projets à tout moment — sans que cette attente me donne prise sur elle.
\end{itemize}
\textbf{3. Arguments.}
\begin{itemize}
\item \emph{Argument 1 (imprévisibilité)} : la mort peut survenir à tout instant, avant l'accomplissement de mes projets ; elle est donc toujours « trop tôt » ou absurde, jamais mon acte : je ne peux pas l'assumer car elle m'échappe par définition (mort, je ne suis plus là pour l'assumer).
\item \emph{Argument 2 (dépossession du sens)} : en mourant, je cesse d'être un pour-soi qui se définit ; le sens de ma vie passe entre les mains d'autrui. La mort est donc la fin de tous mes possibles et le triomphe du point de vue d'autrui.
\end{itemize}
\textbf{4. Discussion critique.} Sartre a raison contre Heidegger sur un point : je ne peux pas faire l'expérience de ma mort ni la choisir comme un projet parmi d'autres ; l'« anticipation résolue » heideggérienne court le risque de transformer en possibilité intérieure ce qui est d'abord un événement qui m'anéantit. Cependant, Heidegger répondrait que l'être-pour-la-mort ne prétend pas maîtriser la mort, mais seulement se l'approprier comme possibilité, ce qui suffit à individualiser et à libérer. De son côté, Épicure dissout le problème : si la mort n'est « rien pour nous », la question de l'assumer ne se pose même pas ; mais cette sérénité se paie d'une indifférence qui peut vider l'existence de son urgence. La mort ne peut être « mon » acte au sens fort ; elle peut en revanche devenir mon horizon : je ne la maîtrise pas, mais je décide de vivre en sachant qu'elle vient — ce que confirme la sagesse africaine, pour qui bien mourir suppose d'avoir transmis.

\textbf{Barème indicatif (sur 20) :} thèse dégagée (3 pts) ; concepts expliqués dans le texte (6 pts) ; arguments relevés et reformulés (5 pts) ; discussion critique confrontant Heidegger et Épicure (4 pts) ; langue et méthode (2 pts).

\textbf{Points de vigilance :} expliquer les concepts \emph{dans le texte} (ne pas réciter le cours général sur Sartre) ; la discussion doit confronter explicitement les trois positions demandées ; distinguer soigneusement « assumer » (Sartre : impossible) et « anticiper » (Heidegger : possible).
\end{corrige}

\begin{corrige}[Exercice 11 — Cas Pratique : L'aide à mourir à l'Hôpital Central de Yaoundé]
\textbf{Raisonnement éthique structuré (modèle de copie).}

\textbf{1. Conflit de valeurs.} Le cas oppose deux univers éthiques. D'un côté, les principes de la bioéthique moderne : l'\emph{autonomie} du patient lucide (M. K. demande par écrit de mourir), la \emph{bienfaisance} (soulager une souffrance réfractaire) et la \emph{non-malfaisance} (ne pas prolonger une agonie douloureuse). De l'autre, l'éthique communautaire bamiléké : la vie n'appartient pas à l'individu mais à la lignée et aux ancêtres ; demander la mort est un suicide, un tabou, une offense aux ancêtres ; la famille se constitue gardienne de cette valeur. Le conflit est donc celui de la \emph{personne-individu autonome} contre la \emph{personne-insérée-dans-la-communauté}.

\textbf{2. Cadre légal camerounais.} Le Code pénal camerounais réprime l'homicide, y compris l'homicide sur demande de la victime : l'euthanasie active (injection létale) est un acte illégal, quel que soit le consentement du patient. Le Code de déontologie médicale interdit au médecin de provoquer délibérément la mort, mais l'oblige à soulager les souffrances et à refuser l'acharnement thérapeutique. La sédation profonde et continue en phase terminale, lorsqu'elle vise la souffrance et non la mort, relève des soins palliatifs licites.

\textbf{3. Argumentation philosophique.}
\begin{itemize}
\item \emph{Autonomie et dignité} : pour Kant, la personne est une fin en soi, jamais un simple moyen ; pour Sartre, l'existence est projet et liberté : nier la volonté lucide de M. K., c'est le traiter comme une chose (en-soi) et nier ce qui fait son humanité. Son consentement éclairé pèse donc d'un poids moral majeur.
\item \emph{Conception communautaire} : mais pour Mbiti (« je suis parce que nous sommes ») et Wiredu, la personne africaine se constitue dans le corps social et la chaîne des générations ; la mort de M. K. n'est pas un événement privé : elle engage la lignée, les ancêtres, les rites. Ignorer la famille serait une violence culturelle et une négation de son identité même.
\item \emph{Laisser mourir / faire mourir} : la doctrine de l'effet double distingue l'intention et l'effet prévisible : administrer une sédation profonde \emph{pour soulager} la douleur, même si elle risque d'abréger la vie, est moralement et légalement distinct de l'injection \emph{destinée à tuer}. Cette distinction permet de respecter à la fois la non-malfaisance, la loi et l'interdit du suicide.
\end{itemize}
\textbf{4. Décision motivée.} Je propose au Comité : (a) le \emph{refus de l'euthanasie active}, illégale et incompatible avec l'interdit communautaire ; (b) la mise en place d'une \emph{sédation palliative profonde et continue} visant la souffrance, conforme à la loi et à la déontologie ; (c) l'arrêt de tout acharnement thérapeutique ; (d) un accompagnement psychologique et spirituel du patient, avec médiation culturelle : implication de la famille, des notables et, si le patient le souhaite, des autorités traditionnelles et religieuses, afin que ses derniers jours demeurent insérés dans le lien lignager ; (e) la consignation écrite de la volonté du patient et du consentement de la famille au protocole palliatif. Cette solution honore l'autonomie de M. K. (sa souffrance est prise au sérieux, sa parole entendue) sans transformer le médecin en donneur de mort, et sans rompre le pacte entre les vivants et les ancêtres.
\textbf{5. Ouverture.} Ce cas montre que la bioéthique universelle des principes ne peut s'appliquer mécaniquement : elle doit dialoguer avec les éthiques africaines de la personne et de la mort. Une loi camerounaise sur la fin de vie, élaborée avec les traditions, reste à penser.

\textbf{Barème indicatif (sur 20) :} conflit de valeurs identifié (3 pts) ; cadre légal exact (3 pts) ; arguments philosophiques (Kant/Sartre, Mbiti/Wiredu, effet double) (6 pts) ; décision concrète et cohérente (4 pts) ; ouverture (2 pts) ; clarté et structure (2 pts).

\textbf{Points de vigilance :} ne jamais proposer l'euthanasie active comme solution finale (illégale au Cameroun) ; la sédation palliative doit être justifiée par la doctrine de l'effet double ; traiter la position familiale avec respect (pas de mépris pour la tradition) ; la décision doit être concrète, argumentée et réaliste.
\end{corrige}

\newpage

\coursec{Fiche bilan}

\begin{popfigure}
\begin{center}
\begin{tikzpicture}[mindmap, grow cyclic, every node/.style={concept, font=\small},
  root concept/.append style={concept color=popink, font=\bfseries},
  level 1/.append style={level distance=3.4cm, sibling angle=72},
  level 1 concept/.append style={font=\footnotesize}]
\node[root concept]{Existence et mort}
  child[concept color=popblue]{ node {D\'efinitions} 
    child[concept color=popblueL]{ node[concept, font=\scriptsize] {Vie : fait biologique} }
    child[concept color=popblueL]{ node[concept, font=\scriptsize] {Existence : vie consciente et libre} }
  }
  child[concept color=poporange]{ node {Sens de l'existence}
    child[concept color=poporangeL]{ node[concept, font=\scriptsize] {Absurde : Camus} }
    child[concept color=poporangeL]{ node[concept, font=\scriptsize] {Projet : Sartre} }
  }
  child[concept color=popgreen]{ node {La mort}
    child[concept color=popgreenL]{ node[concept, font=\scriptsize] {Fin radicale : \'Epicure} }
    child[concept color=popgreenL]{ node[concept, font=\scriptsize] {Horizon de sens : Heidegger} }
  }
  child[concept color=popgold]{ node {Ancestralit\'e au Cameroun}
    child[concept color=popgoldL]{ node[concept, font=\scriptsize] {Culte des anc\^etres, fun\'erailles} }
  }
  child[concept color=poppurple]{ node {M\'ethodes}
    child[concept color=poppurpleL]{ node[concept, font=\scriptsize] {Dissertation : probl\'ematique, plan} }
    child[concept color=poppurpleL]{ node[concept, font=\scriptsize] {Explication de texte} }
  };
\end{tikzpicture}
\end{center}
\legende{Carte mentale de la le\c{c}on : les cinq branches essentielles \`a ma\^itriser pour le Probatoire et le Baccalaur\'eat technique.}
\end{popfigure}

\begin{bilan}
\textbf{D\'efinitions cl\'es}
\begin{itemize}
\item \cle{Vie} : ensemble des ph\'enom\`enes biologiques (naissance, croissance, reproduction, mort) communs \`a tous les \^etres vivants.
\item \cle{Existence} : mani\`ere proprement humaine d'\^etre au monde ; la vie v\'ecue par un \^etre conscient, libre et capable de se questionner sur le sens de sa vie.
\item \cle{Mort} : cessation d\'efinitive des fonctions vitales ; pour l'homme, elle est aussi un \emph{\'ev\'enement pens\'e} qui donne son prix au temps.
\item \cle{Absurde} : conflit entre le d\'esir humain de sens et le silence du monde (Camus).
\item \cle{Ancestralit\'e} : survie symbolique du d\'efunt dans la m\'emoire et les rites de la communaut\'e.
\end{itemize}

\textbf{Th\`eses d'auteurs \`a conna\^itre par c\oe ur}
\begin{center}
\begin{tabularx}{\linewidth}{|l|X|}
\hline
\rowcolor{popblueL}\textbf{Auteur} & \textbf{Th\`ese} \\
\hline
\'Epicure & << La mort n'est rien pour nous >> : tant que nous sommes, la mort n'est pas ; quand elle est l\`a, nous ne sommes plus. \\
\hline
S\'en\`eque & Apprendre \`a mourir, c'est apprendre \`a vivre libre : la m\'editation de la mort d\'elivre de la peur. \\
\hline
Heidegger & L'homme est un << \^etre-pour-la-mort >> : la conscience de la finitude rend l'existence authentique. \\
\hline
Camus & L'absurde na\^it du divorce entre l'homme et le monde ; il faut se r\'evolter et vivre sans espoir illusoire. \\
\hline
Sartre & L'existence pr\'ec\`ede l'essence : l'homme se fait lui-m\^eme par ses choix et ses projets. \\
\hline
\end{tabularx}
\end{center}

\textbf{M\'ethodes express}
\begin{itemize}
\item \emph{Dissertation} : analyser le sujet $\rightarrow$ d\'egager la probl\'ematique $\rightarrow$ plan dialectique (th\`ese / antith\`ese / synth\`ese) $\rightarrow$ introduction (accroche, d\'efinitions, probl\'ematique, annonce du plan) $\rightarrow$ d\'eveloppement argument\'e avec exemples $\rightarrow$ conclusion.
\item \emph{Explication de texte} : situer l'auteur et l'\oe uvre $\rightarrow$ d\'egager la th\`ese $\rightarrow$ suivre la structure du texte $\rightarrow$ expliquer les arguments $\rightarrow$ discuter bri\`evement.
\item \emph{Toujours} : d\'efinir les termes du sujet d\`es l'introduction et illustrer par des exemples concrets (fun\'erailles, deuil, rites d'ancestralit\'e au Cameroun).
\end{itemize}
\end{bilan}

\begin{aretenir}[Checklist avant le Bac]
\begin{itemize}
\item[$\square$] Je sais d\'efinir \cle{vie}, \cle{existence} et \cle{mort} et les distinguer clairement.
\item[$\square$] Je sais expliquer pourquoi l'existence humaine est une vie \emph{consciente et libre}.
\item[$\square$] Je connais la th\`ese d'\'Epicure sur la mort et je sais la citer.
\item[$\square$] Je connais la th\`ese de Heidegger : l'\^etre-pour-la-mort.
\item[$\square$] Je sais expliquer la notion d'\cle{absurde} chez Camus.
\item[$\square$] Je sais expliquer << l'existence pr\'ec\`ede l'essence >> chez Sartre.
\item[$\square$] Je peux mobiliser un exemple camerounais (fun\'erailles, culte des anc\^etres) pour illustrer la survie symbolique.
\item[$\square$] Je sais construire une probl\'ematique \`a partir d'un sujet sur l'existence ou la mort.
\item[$\square$] Je ma\^itriser les \'etapes de l'introduction de dissertation.
\item[$\square$] Je sais r\'ediger un plan dialectique en trois parties.
\item[$\square$] Je sais confronter deux th\`eses oppos\'ees (fin radicale / horizon de sens) et justifier ma conclusion.
\item[$\square$] Je relis ma copie : orthographe, articulations logiques, exemples pr\'ecis.
\end{itemize}
\end{aretenir}

\findecours

\end{document}
